% Shared definitions for the badness brand artifacts.
%
% \input by the per-artifact .tex files (logo.tex, wordmark.tex, og.tex).
% Anything \providecommand'd here is overridable by the caller — define
% your override BEFORE % Shared definitions for the badness brand artifacts.
%
% \input by the per-artifact .tex files (logo.tex, wordmark.tex, og.tex).
% Anything \providecommand'd here is overridable by the caller — define
% your override BEFORE % Shared definitions for the badness brand artifacts.
%
% \input by the per-artifact .tex files (logo.tex, wordmark.tex, og.tex).
% Anything \providecommand'd here is overridable by the caller — define
% your override BEFORE % Shared definitions for the badness brand artifacts.
%
% \input by the per-artifact .tex files (logo.tex, wordmark.tex, og.tex).
% Anything \providecommand'd here is overridable by the caller — define
% your override BEFORE \input{_common.tex}.

\usepackage{tikz}
\usepackage{xcolor}

% --- Font dispatch ----------------------------------------------------------
%
% \bnFont selects the sans-serif family. Set it in your artifact file
% before \input{_common.tex}. Available values:
%
%   plex       IBM Plex Sans         — modern technical (default)
%   tgheros    TeX Gyre Heros        — Helvetica-clone, neutral
%   source     Source Sans Pro       — humanist, friendly
%   roboto     Roboto                — Google flat-ui geometric
%   fira       Fira Sans             — Mozilla, slightly quirky
%   inter      Inter                 — UI workhorse, very tight
%   montserrat Montserrat            — wide geometric, posterish
%   lato       Lato                  — warm humanist
%   cmbright   Computer Modern Bright — light, scholarly
%   lm         Latin Modern Sans     — LaTeX default, neutral
%
\providecommand{\bnFont}{lm}

\makeatletter
\def\bn@font@plex      {\RequirePackage{plex-sans}}
\def\bn@font@tgheros   {\RequirePackage{tgheros}}
\def\bn@font@source    {\RequirePackage[default]{sourcesanspro}}
\def\bn@font@roboto    {\RequirePackage[sfdefault]{roboto}}
\def\bn@font@fira      {\RequirePackage[sfdefault]{FiraSans}}
\def\bn@font@inter     {\RequirePackage{inter}}
\def\bn@font@montserrat{\RequirePackage[defaultfam,tabular,lining]{montserrat}}
\def\bn@font@lato      {\RequirePackage[default]{lato}}
\def\bn@font@cmbright  {\RequirePackage{cmbright}}
\def\bn@font@lm        {\RequirePackage{lmodern}}
\@nameuse{bn@font@\bnFont}
\makeatother
\renewcommand{\familydefault}{\sfdefault}

% --- Palette ----------------------------------------------------------------
%
% Five role-named slots. \bnDarkMode=1 swaps to a dark variant.
% Override any individual color by \definecolor'ing it AFTER \input{_common.tex}.

\providecommand{\bnDarkMode}{0}

\ifnum\bnDarkMode=1
  % Dark variant
  \definecolor{bnInk}    {HTML}{F2EFE9} % "ink" reads light on dark
  \definecolor{bnBad}    {HTML}{F08A57} % brighter orange so it pops on dark
  \definecolor{bnPaper}  {HTML}{1B1B1B} % near-black canvas
  \definecolor{bnMargin} {HTML}{3D3833} % subtle margin guide
  \definecolor{bnMuted}  {HTML}{A39A8C} % secondary text
\else
  % Light variant (default)
  \definecolor{bnInk}    {HTML}{1B1B1B}
  \definecolor{bnBad}    {HTML}{D9633F}
  \definecolor{bnPaper}  {HTML}{FAFAF7}
  \definecolor{bnMargin} {HTML}{B5B0A5}
  \definecolor{bnMuted}  {HTML}{6F665A}
\fi

% --- Logo geometry (all in cm; tweak freely) --------------------------------

\providecommand{\bnBarCount}{4}     % number of horizontal bars
\providecommand{\bnBarWidth}{4.0}   % length of an in-margin bar
\providecommand{\bnBarHeight}{0.55} % bar thickness
\providecommand{\bnBarGap}{0.32}    % vertical gap between bars
\providecommand{\bnJutWhich}{3}     % 1..bnBarCount — which bar overflows
\providecommand{\bnJutExtra}{1.15}  % how far it sticks out
\providecommand{\bnRound}{0.07}     % corner rounding
\providecommand{\bnShowMargin}{0}   % 1 = draw the dashed margin guide, 0 = hide

% --- The mark itself --------------------------------------------------------
%
% Draws into the current tikzpicture, anchored so the bottom-left of the
% non-jutting bars is at (0,0). The jutting bar extends to the right.

\newcommand{\bnDrawMark}{%
  \pgfmathtruncatemacro{\bnLast}{\bnBarCount-1}%
  \pgfmathsetmacro{\bnTotalH}{\bnBarCount*\bnBarHeight + \bnLast*\bnBarGap}%
  \ifnum\bnShowMargin=1
    % The right-margin guide: a dashed vertical line at the in-margin width,
    % running the full height of the bar stack with a touch of overshoot.
    \draw[bnMargin, dashed, line width=0.4pt, dash pattern=on 2pt off 2pt]
    (\bnBarWidth, -0.12) -- (\bnBarWidth, \bnTotalH+0.12);
  \fi
  \foreach \i in {0,...,\bnLast}{%
      \pgfmathsetmacro{\bny}{(\bnLast-\i)*(\bnBarHeight+\bnBarGap)}%
      \pgfmathtruncatemacro{\bnNum}{\i+1}%
      \ifnum\bnNum=\bnJutWhich
        \fill[bnBad, rounded corners=\bnRound cm]
        (0,\bny) rectangle (\bnBarWidth+\bnJutExtra, \bny+\bnBarHeight);
      \else
        \fill[bnInk, rounded corners=\bnRound cm]
        (0,\bny) rectangle (\bnBarWidth, \bny+\bnBarHeight);
      \fi
    }%
}
.

\usepackage{tikz}
\usepackage{xcolor}

% --- Font dispatch ----------------------------------------------------------
%
% \bnFont selects the sans-serif family. Set it in your artifact file
% before % Shared definitions for the badness brand artifacts.
%
% \input by the per-artifact .tex files (logo.tex, wordmark.tex, og.tex).
% Anything \providecommand'd here is overridable by the caller — define
% your override BEFORE \input{_common.tex}.

\usepackage{tikz}
\usepackage{xcolor}

% --- Font dispatch ----------------------------------------------------------
%
% \bnFont selects the sans-serif family. Set it in your artifact file
% before \input{_common.tex}. Available values:
%
%   plex       IBM Plex Sans         — modern technical (default)
%   tgheros    TeX Gyre Heros        — Helvetica-clone, neutral
%   source     Source Sans Pro       — humanist, friendly
%   roboto     Roboto                — Google flat-ui geometric
%   fira       Fira Sans             — Mozilla, slightly quirky
%   inter      Inter                 — UI workhorse, very tight
%   montserrat Montserrat            — wide geometric, posterish
%   lato       Lato                  — warm humanist
%   cmbright   Computer Modern Bright — light, scholarly
%   lm         Latin Modern Sans     — LaTeX default, neutral
%
\providecommand{\bnFont}{lm}

\makeatletter
\def\bn@font@plex      {\RequirePackage{plex-sans}}
\def\bn@font@tgheros   {\RequirePackage{tgheros}}
\def\bn@font@source    {\RequirePackage[default]{sourcesanspro}}
\def\bn@font@roboto    {\RequirePackage[sfdefault]{roboto}}
\def\bn@font@fira      {\RequirePackage[sfdefault]{FiraSans}}
\def\bn@font@inter     {\RequirePackage{inter}}
\def\bn@font@montserrat{\RequirePackage[defaultfam,tabular,lining]{montserrat}}
\def\bn@font@lato      {\RequirePackage[default]{lato}}
\def\bn@font@cmbright  {\RequirePackage{cmbright}}
\def\bn@font@lm        {\RequirePackage{lmodern}}
\@nameuse{bn@font@\bnFont}
\makeatother
\renewcommand{\familydefault}{\sfdefault}

% --- Palette ----------------------------------------------------------------
%
% Five role-named slots. \bnDarkMode=1 swaps to a dark variant.
% Override any individual color by \definecolor'ing it AFTER \input{_common.tex}.

\providecommand{\bnDarkMode}{0}

\ifnum\bnDarkMode=1
  % Dark variant
  \definecolor{bnInk}    {HTML}{F2EFE9} % "ink" reads light on dark
  \definecolor{bnBad}    {HTML}{F08A57} % brighter orange so it pops on dark
  \definecolor{bnPaper}  {HTML}{1B1B1B} % near-black canvas
  \definecolor{bnMargin} {HTML}{3D3833} % subtle margin guide
  \definecolor{bnMuted}  {HTML}{A39A8C} % secondary text
\else
  % Light variant (default)
  \definecolor{bnInk}    {HTML}{1B1B1B}
  \definecolor{bnBad}    {HTML}{D9633F}
  \definecolor{bnPaper}  {HTML}{FAFAF7}
  \definecolor{bnMargin} {HTML}{B5B0A5}
  \definecolor{bnMuted}  {HTML}{6F665A}
\fi

% --- Logo geometry (all in cm; tweak freely) --------------------------------

\providecommand{\bnBarCount}{4}     % number of horizontal bars
\providecommand{\bnBarWidth}{4.0}   % length of an in-margin bar
\providecommand{\bnBarHeight}{0.55} % bar thickness
\providecommand{\bnBarGap}{0.32}    % vertical gap between bars
\providecommand{\bnJutWhich}{3}     % 1..bnBarCount — which bar overflows
\providecommand{\bnJutExtra}{1.15}  % how far it sticks out
\providecommand{\bnRound}{0.07}     % corner rounding
\providecommand{\bnShowMargin}{0}   % 1 = draw the dashed margin guide, 0 = hide

% --- The mark itself --------------------------------------------------------
%
% Draws into the current tikzpicture, anchored so the bottom-left of the
% non-jutting bars is at (0,0). The jutting bar extends to the right.

\newcommand{\bnDrawMark}{%
  \pgfmathtruncatemacro{\bnLast}{\bnBarCount-1}%
  \pgfmathsetmacro{\bnTotalH}{\bnBarCount*\bnBarHeight + \bnLast*\bnBarGap}%
  \ifnum\bnShowMargin=1
    % The right-margin guide: a dashed vertical line at the in-margin width,
    % running the full height of the bar stack with a touch of overshoot.
    \draw[bnMargin, dashed, line width=0.4pt, dash pattern=on 2pt off 2pt]
    (\bnBarWidth, -0.12) -- (\bnBarWidth, \bnTotalH+0.12);
  \fi
  \foreach \i in {0,...,\bnLast}{%
      \pgfmathsetmacro{\bny}{(\bnLast-\i)*(\bnBarHeight+\bnBarGap)}%
      \pgfmathtruncatemacro{\bnNum}{\i+1}%
      \ifnum\bnNum=\bnJutWhich
        \fill[bnBad, rounded corners=\bnRound cm]
        (0,\bny) rectangle (\bnBarWidth+\bnJutExtra, \bny+\bnBarHeight);
      \else
        \fill[bnInk, rounded corners=\bnRound cm]
        (0,\bny) rectangle (\bnBarWidth, \bny+\bnBarHeight);
      \fi
    }%
}
. Available values:
%
%   plex       IBM Plex Sans         — modern technical (default)
%   tgheros    TeX Gyre Heros        — Helvetica-clone, neutral
%   source     Source Sans Pro       — humanist, friendly
%   roboto     Roboto                — Google flat-ui geometric
%   fira       Fira Sans             — Mozilla, slightly quirky
%   inter      Inter                 — UI workhorse, very tight
%   montserrat Montserrat            — wide geometric, posterish
%   lato       Lato                  — warm humanist
%   cmbright   Computer Modern Bright — light, scholarly
%   lm         Latin Modern Sans     — LaTeX default, neutral
%
\providecommand{\bnFont}{lm}

\makeatletter
\def\bn@font@plex      {\RequirePackage{plex-sans}}
\def\bn@font@tgheros   {\RequirePackage{tgheros}}
\def\bn@font@source    {\RequirePackage[default]{sourcesanspro}}
\def\bn@font@roboto    {\RequirePackage[sfdefault]{roboto}}
\def\bn@font@fira      {\RequirePackage[sfdefault]{FiraSans}}
\def\bn@font@inter     {\RequirePackage{inter}}
\def\bn@font@montserrat{\RequirePackage[defaultfam,tabular,lining]{montserrat}}
\def\bn@font@lato      {\RequirePackage[default]{lato}}
\def\bn@font@cmbright  {\RequirePackage{cmbright}}
\def\bn@font@lm        {\RequirePackage{lmodern}}
\@nameuse{bn@font@\bnFont}
\makeatother
\renewcommand{\familydefault}{\sfdefault}

% --- Palette ----------------------------------------------------------------
%
% Five role-named slots. \bnDarkMode=1 swaps to a dark variant.
% Override any individual color by \definecolor'ing it AFTER % Shared definitions for the badness brand artifacts.
%
% \input by the per-artifact .tex files (logo.tex, wordmark.tex, og.tex).
% Anything \providecommand'd here is overridable by the caller — define
% your override BEFORE \input{_common.tex}.

\usepackage{tikz}
\usepackage{xcolor}

% --- Font dispatch ----------------------------------------------------------
%
% \bnFont selects the sans-serif family. Set it in your artifact file
% before \input{_common.tex}. Available values:
%
%   plex       IBM Plex Sans         — modern technical (default)
%   tgheros    TeX Gyre Heros        — Helvetica-clone, neutral
%   source     Source Sans Pro       — humanist, friendly
%   roboto     Roboto                — Google flat-ui geometric
%   fira       Fira Sans             — Mozilla, slightly quirky
%   inter      Inter                 — UI workhorse, very tight
%   montserrat Montserrat            — wide geometric, posterish
%   lato       Lato                  — warm humanist
%   cmbright   Computer Modern Bright — light, scholarly
%   lm         Latin Modern Sans     — LaTeX default, neutral
%
\providecommand{\bnFont}{lm}

\makeatletter
\def\bn@font@plex      {\RequirePackage{plex-sans}}
\def\bn@font@tgheros   {\RequirePackage{tgheros}}
\def\bn@font@source    {\RequirePackage[default]{sourcesanspro}}
\def\bn@font@roboto    {\RequirePackage[sfdefault]{roboto}}
\def\bn@font@fira      {\RequirePackage[sfdefault]{FiraSans}}
\def\bn@font@inter     {\RequirePackage{inter}}
\def\bn@font@montserrat{\RequirePackage[defaultfam,tabular,lining]{montserrat}}
\def\bn@font@lato      {\RequirePackage[default]{lato}}
\def\bn@font@cmbright  {\RequirePackage{cmbright}}
\def\bn@font@lm        {\RequirePackage{lmodern}}
\@nameuse{bn@font@\bnFont}
\makeatother
\renewcommand{\familydefault}{\sfdefault}

% --- Palette ----------------------------------------------------------------
%
% Five role-named slots. \bnDarkMode=1 swaps to a dark variant.
% Override any individual color by \definecolor'ing it AFTER \input{_common.tex}.

\providecommand{\bnDarkMode}{0}

\ifnum\bnDarkMode=1
  % Dark variant
  \definecolor{bnInk}    {HTML}{F2EFE9} % "ink" reads light on dark
  \definecolor{bnBad}    {HTML}{F08A57} % brighter orange so it pops on dark
  \definecolor{bnPaper}  {HTML}{1B1B1B} % near-black canvas
  \definecolor{bnMargin} {HTML}{3D3833} % subtle margin guide
  \definecolor{bnMuted}  {HTML}{A39A8C} % secondary text
\else
  % Light variant (default)
  \definecolor{bnInk}    {HTML}{1B1B1B}
  \definecolor{bnBad}    {HTML}{D9633F}
  \definecolor{bnPaper}  {HTML}{FAFAF7}
  \definecolor{bnMargin} {HTML}{B5B0A5}
  \definecolor{bnMuted}  {HTML}{6F665A}
\fi

% --- Logo geometry (all in cm; tweak freely) --------------------------------

\providecommand{\bnBarCount}{4}     % number of horizontal bars
\providecommand{\bnBarWidth}{4.0}   % length of an in-margin bar
\providecommand{\bnBarHeight}{0.55} % bar thickness
\providecommand{\bnBarGap}{0.32}    % vertical gap between bars
\providecommand{\bnJutWhich}{3}     % 1..bnBarCount — which bar overflows
\providecommand{\bnJutExtra}{1.15}  % how far it sticks out
\providecommand{\bnRound}{0.07}     % corner rounding
\providecommand{\bnShowMargin}{0}   % 1 = draw the dashed margin guide, 0 = hide

% --- The mark itself --------------------------------------------------------
%
% Draws into the current tikzpicture, anchored so the bottom-left of the
% non-jutting bars is at (0,0). The jutting bar extends to the right.

\newcommand{\bnDrawMark}{%
  \pgfmathtruncatemacro{\bnLast}{\bnBarCount-1}%
  \pgfmathsetmacro{\bnTotalH}{\bnBarCount*\bnBarHeight + \bnLast*\bnBarGap}%
  \ifnum\bnShowMargin=1
    % The right-margin guide: a dashed vertical line at the in-margin width,
    % running the full height of the bar stack with a touch of overshoot.
    \draw[bnMargin, dashed, line width=0.4pt, dash pattern=on 2pt off 2pt]
    (\bnBarWidth, -0.12) -- (\bnBarWidth, \bnTotalH+0.12);
  \fi
  \foreach \i in {0,...,\bnLast}{%
      \pgfmathsetmacro{\bny}{(\bnLast-\i)*(\bnBarHeight+\bnBarGap)}%
      \pgfmathtruncatemacro{\bnNum}{\i+1}%
      \ifnum\bnNum=\bnJutWhich
        \fill[bnBad, rounded corners=\bnRound cm]
        (0,\bny) rectangle (\bnBarWidth+\bnJutExtra, \bny+\bnBarHeight);
      \else
        \fill[bnInk, rounded corners=\bnRound cm]
        (0,\bny) rectangle (\bnBarWidth, \bny+\bnBarHeight);
      \fi
    }%
}
.

\providecommand{\bnDarkMode}{0}

\ifnum\bnDarkMode=1
  % Dark variant
  \definecolor{bnInk}    {HTML}{F2EFE9} % "ink" reads light on dark
  \definecolor{bnBad}    {HTML}{F08A57} % brighter orange so it pops on dark
  \definecolor{bnPaper}  {HTML}{1B1B1B} % near-black canvas
  \definecolor{bnMargin} {HTML}{3D3833} % subtle margin guide
  \definecolor{bnMuted}  {HTML}{A39A8C} % secondary text
\else
  % Light variant (default)
  \definecolor{bnInk}    {HTML}{1B1B1B}
  \definecolor{bnBad}    {HTML}{D9633F}
  \definecolor{bnPaper}  {HTML}{FAFAF7}
  \definecolor{bnMargin} {HTML}{B5B0A5}
  \definecolor{bnMuted}  {HTML}{6F665A}
\fi

% --- Logo geometry (all in cm; tweak freely) --------------------------------

\providecommand{\bnBarCount}{4}     % number of horizontal bars
\providecommand{\bnBarWidth}{4.0}   % length of an in-margin bar
\providecommand{\bnBarHeight}{0.55} % bar thickness
\providecommand{\bnBarGap}{0.32}    % vertical gap between bars
\providecommand{\bnJutWhich}{3}     % 1..bnBarCount — which bar overflows
\providecommand{\bnJutExtra}{1.15}  % how far it sticks out
\providecommand{\bnRound}{0.07}     % corner rounding
\providecommand{\bnShowMargin}{0}   % 1 = draw the dashed margin guide, 0 = hide

% --- The mark itself --------------------------------------------------------
%
% Draws into the current tikzpicture, anchored so the bottom-left of the
% non-jutting bars is at (0,0). The jutting bar extends to the right.

\newcommand{\bnDrawMark}{%
  \pgfmathtruncatemacro{\bnLast}{\bnBarCount-1}%
  \pgfmathsetmacro{\bnTotalH}{\bnBarCount*\bnBarHeight + \bnLast*\bnBarGap}%
  \ifnum\bnShowMargin=1
    % The right-margin guide: a dashed vertical line at the in-margin width,
    % running the full height of the bar stack with a touch of overshoot.
    \draw[bnMargin, dashed, line width=0.4pt, dash pattern=on 2pt off 2pt]
    (\bnBarWidth, -0.12) -- (\bnBarWidth, \bnTotalH+0.12);
  \fi
  \foreach \i in {0,...,\bnLast}{%
      \pgfmathsetmacro{\bny}{(\bnLast-\i)*(\bnBarHeight+\bnBarGap)}%
      \pgfmathtruncatemacro{\bnNum}{\i+1}%
      \ifnum\bnNum=\bnJutWhich
        \fill[bnBad, rounded corners=\bnRound cm]
        (0,\bny) rectangle (\bnBarWidth+\bnJutExtra, \bny+\bnBarHeight);
      \else
        \fill[bnInk, rounded corners=\bnRound cm]
        (0,\bny) rectangle (\bnBarWidth, \bny+\bnBarHeight);
      \fi
    }%
}
.

\usepackage{tikz}
\usepackage{xcolor}

% --- Font dispatch ----------------------------------------------------------
%
% \bnFont selects the sans-serif family. Set it in your artifact file
% before % Shared definitions for the badness brand artifacts.
%
% \input by the per-artifact .tex files (logo.tex, wordmark.tex, og.tex).
% Anything \providecommand'd here is overridable by the caller — define
% your override BEFORE % Shared definitions for the badness brand artifacts.
%
% \input by the per-artifact .tex files (logo.tex, wordmark.tex, og.tex).
% Anything \providecommand'd here is overridable by the caller — define
% your override BEFORE \input{_common.tex}.

\usepackage{tikz}
\usepackage{xcolor}

% --- Font dispatch ----------------------------------------------------------
%
% \bnFont selects the sans-serif family. Set it in your artifact file
% before \input{_common.tex}. Available values:
%
%   plex       IBM Plex Sans         — modern technical (default)
%   tgheros    TeX Gyre Heros        — Helvetica-clone, neutral
%   source     Source Sans Pro       — humanist, friendly
%   roboto     Roboto                — Google flat-ui geometric
%   fira       Fira Sans             — Mozilla, slightly quirky
%   inter      Inter                 — UI workhorse, very tight
%   montserrat Montserrat            — wide geometric, posterish
%   lato       Lato                  — warm humanist
%   cmbright   Computer Modern Bright — light, scholarly
%   lm         Latin Modern Sans     — LaTeX default, neutral
%
\providecommand{\bnFont}{lm}

\makeatletter
\def\bn@font@plex      {\RequirePackage{plex-sans}}
\def\bn@font@tgheros   {\RequirePackage{tgheros}}
\def\bn@font@source    {\RequirePackage[default]{sourcesanspro}}
\def\bn@font@roboto    {\RequirePackage[sfdefault]{roboto}}
\def\bn@font@fira      {\RequirePackage[sfdefault]{FiraSans}}
\def\bn@font@inter     {\RequirePackage{inter}}
\def\bn@font@montserrat{\RequirePackage[defaultfam,tabular,lining]{montserrat}}
\def\bn@font@lato      {\RequirePackage[default]{lato}}
\def\bn@font@cmbright  {\RequirePackage{cmbright}}
\def\bn@font@lm        {\RequirePackage{lmodern}}
\@nameuse{bn@font@\bnFont}
\makeatother
\renewcommand{\familydefault}{\sfdefault}

% --- Palette ----------------------------------------------------------------
%
% Five role-named slots. \bnDarkMode=1 swaps to a dark variant.
% Override any individual color by \definecolor'ing it AFTER \input{_common.tex}.

\providecommand{\bnDarkMode}{0}

\ifnum\bnDarkMode=1
  % Dark variant
  \definecolor{bnInk}    {HTML}{F2EFE9} % "ink" reads light on dark
  \definecolor{bnBad}    {HTML}{F08A57} % brighter orange so it pops on dark
  \definecolor{bnPaper}  {HTML}{1B1B1B} % near-black canvas
  \definecolor{bnMargin} {HTML}{3D3833} % subtle margin guide
  \definecolor{bnMuted}  {HTML}{A39A8C} % secondary text
\else
  % Light variant (default)
  \definecolor{bnInk}    {HTML}{1B1B1B}
  \definecolor{bnBad}    {HTML}{D9633F}
  \definecolor{bnPaper}  {HTML}{FAFAF7}
  \definecolor{bnMargin} {HTML}{B5B0A5}
  \definecolor{bnMuted}  {HTML}{6F665A}
\fi

% --- Logo geometry (all in cm; tweak freely) --------------------------------

\providecommand{\bnBarCount}{4}     % number of horizontal bars
\providecommand{\bnBarWidth}{4.0}   % length of an in-margin bar
\providecommand{\bnBarHeight}{0.55} % bar thickness
\providecommand{\bnBarGap}{0.32}    % vertical gap between bars
\providecommand{\bnJutWhich}{3}     % 1..bnBarCount — which bar overflows
\providecommand{\bnJutExtra}{1.15}  % how far it sticks out
\providecommand{\bnRound}{0.07}     % corner rounding
\providecommand{\bnShowMargin}{0}   % 1 = draw the dashed margin guide, 0 = hide

% --- The mark itself --------------------------------------------------------
%
% Draws into the current tikzpicture, anchored so the bottom-left of the
% non-jutting bars is at (0,0). The jutting bar extends to the right.

\newcommand{\bnDrawMark}{%
  \pgfmathtruncatemacro{\bnLast}{\bnBarCount-1}%
  \pgfmathsetmacro{\bnTotalH}{\bnBarCount*\bnBarHeight + \bnLast*\bnBarGap}%
  \ifnum\bnShowMargin=1
    % The right-margin guide: a dashed vertical line at the in-margin width,
    % running the full height of the bar stack with a touch of overshoot.
    \draw[bnMargin, dashed, line width=0.4pt, dash pattern=on 2pt off 2pt]
    (\bnBarWidth, -0.12) -- (\bnBarWidth, \bnTotalH+0.12);
  \fi
  \foreach \i in {0,...,\bnLast}{%
      \pgfmathsetmacro{\bny}{(\bnLast-\i)*(\bnBarHeight+\bnBarGap)}%
      \pgfmathtruncatemacro{\bnNum}{\i+1}%
      \ifnum\bnNum=\bnJutWhich
        \fill[bnBad, rounded corners=\bnRound cm]
        (0,\bny) rectangle (\bnBarWidth+\bnJutExtra, \bny+\bnBarHeight);
      \else
        \fill[bnInk, rounded corners=\bnRound cm]
        (0,\bny) rectangle (\bnBarWidth, \bny+\bnBarHeight);
      \fi
    }%
}
.

\usepackage{tikz}
\usepackage{xcolor}

% --- Font dispatch ----------------------------------------------------------
%
% \bnFont selects the sans-serif family. Set it in your artifact file
% before % Shared definitions for the badness brand artifacts.
%
% \input by the per-artifact .tex files (logo.tex, wordmark.tex, og.tex).
% Anything \providecommand'd here is overridable by the caller — define
% your override BEFORE \input{_common.tex}.

\usepackage{tikz}
\usepackage{xcolor}

% --- Font dispatch ----------------------------------------------------------
%
% \bnFont selects the sans-serif family. Set it in your artifact file
% before \input{_common.tex}. Available values:
%
%   plex       IBM Plex Sans         — modern technical (default)
%   tgheros    TeX Gyre Heros        — Helvetica-clone, neutral
%   source     Source Sans Pro       — humanist, friendly
%   roboto     Roboto                — Google flat-ui geometric
%   fira       Fira Sans             — Mozilla, slightly quirky
%   inter      Inter                 — UI workhorse, very tight
%   montserrat Montserrat            — wide geometric, posterish
%   lato       Lato                  — warm humanist
%   cmbright   Computer Modern Bright — light, scholarly
%   lm         Latin Modern Sans     — LaTeX default, neutral
%
\providecommand{\bnFont}{lm}

\makeatletter
\def\bn@font@plex      {\RequirePackage{plex-sans}}
\def\bn@font@tgheros   {\RequirePackage{tgheros}}
\def\bn@font@source    {\RequirePackage[default]{sourcesanspro}}
\def\bn@font@roboto    {\RequirePackage[sfdefault]{roboto}}
\def\bn@font@fira      {\RequirePackage[sfdefault]{FiraSans}}
\def\bn@font@inter     {\RequirePackage{inter}}
\def\bn@font@montserrat{\RequirePackage[defaultfam,tabular,lining]{montserrat}}
\def\bn@font@lato      {\RequirePackage[default]{lato}}
\def\bn@font@cmbright  {\RequirePackage{cmbright}}
\def\bn@font@lm        {\RequirePackage{lmodern}}
\@nameuse{bn@font@\bnFont}
\makeatother
\renewcommand{\familydefault}{\sfdefault}

% --- Palette ----------------------------------------------------------------
%
% Five role-named slots. \bnDarkMode=1 swaps to a dark variant.
% Override any individual color by \definecolor'ing it AFTER \input{_common.tex}.

\providecommand{\bnDarkMode}{0}

\ifnum\bnDarkMode=1
  % Dark variant
  \definecolor{bnInk}    {HTML}{F2EFE9} % "ink" reads light on dark
  \definecolor{bnBad}    {HTML}{F08A57} % brighter orange so it pops on dark
  \definecolor{bnPaper}  {HTML}{1B1B1B} % near-black canvas
  \definecolor{bnMargin} {HTML}{3D3833} % subtle margin guide
  \definecolor{bnMuted}  {HTML}{A39A8C} % secondary text
\else
  % Light variant (default)
  \definecolor{bnInk}    {HTML}{1B1B1B}
  \definecolor{bnBad}    {HTML}{D9633F}
  \definecolor{bnPaper}  {HTML}{FAFAF7}
  \definecolor{bnMargin} {HTML}{B5B0A5}
  \definecolor{bnMuted}  {HTML}{6F665A}
\fi

% --- Logo geometry (all in cm; tweak freely) --------------------------------

\providecommand{\bnBarCount}{4}     % number of horizontal bars
\providecommand{\bnBarWidth}{4.0}   % length of an in-margin bar
\providecommand{\bnBarHeight}{0.55} % bar thickness
\providecommand{\bnBarGap}{0.32}    % vertical gap between bars
\providecommand{\bnJutWhich}{3}     % 1..bnBarCount — which bar overflows
\providecommand{\bnJutExtra}{1.15}  % how far it sticks out
\providecommand{\bnRound}{0.07}     % corner rounding
\providecommand{\bnShowMargin}{0}   % 1 = draw the dashed margin guide, 0 = hide

% --- The mark itself --------------------------------------------------------
%
% Draws into the current tikzpicture, anchored so the bottom-left of the
% non-jutting bars is at (0,0). The jutting bar extends to the right.

\newcommand{\bnDrawMark}{%
  \pgfmathtruncatemacro{\bnLast}{\bnBarCount-1}%
  \pgfmathsetmacro{\bnTotalH}{\bnBarCount*\bnBarHeight + \bnLast*\bnBarGap}%
  \ifnum\bnShowMargin=1
    % The right-margin guide: a dashed vertical line at the in-margin width,
    % running the full height of the bar stack with a touch of overshoot.
    \draw[bnMargin, dashed, line width=0.4pt, dash pattern=on 2pt off 2pt]
    (\bnBarWidth, -0.12) -- (\bnBarWidth, \bnTotalH+0.12);
  \fi
  \foreach \i in {0,...,\bnLast}{%
      \pgfmathsetmacro{\bny}{(\bnLast-\i)*(\bnBarHeight+\bnBarGap)}%
      \pgfmathtruncatemacro{\bnNum}{\i+1}%
      \ifnum\bnNum=\bnJutWhich
        \fill[bnBad, rounded corners=\bnRound cm]
        (0,\bny) rectangle (\bnBarWidth+\bnJutExtra, \bny+\bnBarHeight);
      \else
        \fill[bnInk, rounded corners=\bnRound cm]
        (0,\bny) rectangle (\bnBarWidth, \bny+\bnBarHeight);
      \fi
    }%
}
. Available values:
%
%   plex       IBM Plex Sans         — modern technical (default)
%   tgheros    TeX Gyre Heros        — Helvetica-clone, neutral
%   source     Source Sans Pro       — humanist, friendly
%   roboto     Roboto                — Google flat-ui geometric
%   fira       Fira Sans             — Mozilla, slightly quirky
%   inter      Inter                 — UI workhorse, very tight
%   montserrat Montserrat            — wide geometric, posterish
%   lato       Lato                  — warm humanist
%   cmbright   Computer Modern Bright — light, scholarly
%   lm         Latin Modern Sans     — LaTeX default, neutral
%
\providecommand{\bnFont}{lm}

\makeatletter
\def\bn@font@plex      {\RequirePackage{plex-sans}}
\def\bn@font@tgheros   {\RequirePackage{tgheros}}
\def\bn@font@source    {\RequirePackage[default]{sourcesanspro}}
\def\bn@font@roboto    {\RequirePackage[sfdefault]{roboto}}
\def\bn@font@fira      {\RequirePackage[sfdefault]{FiraSans}}
\def\bn@font@inter     {\RequirePackage{inter}}
\def\bn@font@montserrat{\RequirePackage[defaultfam,tabular,lining]{montserrat}}
\def\bn@font@lato      {\RequirePackage[default]{lato}}
\def\bn@font@cmbright  {\RequirePackage{cmbright}}
\def\bn@font@lm        {\RequirePackage{lmodern}}
\@nameuse{bn@font@\bnFont}
\makeatother
\renewcommand{\familydefault}{\sfdefault}

% --- Palette ----------------------------------------------------------------
%
% Five role-named slots. \bnDarkMode=1 swaps to a dark variant.
% Override any individual color by \definecolor'ing it AFTER % Shared definitions for the badness brand artifacts.
%
% \input by the per-artifact .tex files (logo.tex, wordmark.tex, og.tex).
% Anything \providecommand'd here is overridable by the caller — define
% your override BEFORE \input{_common.tex}.

\usepackage{tikz}
\usepackage{xcolor}

% --- Font dispatch ----------------------------------------------------------
%
% \bnFont selects the sans-serif family. Set it in your artifact file
% before \input{_common.tex}. Available values:
%
%   plex       IBM Plex Sans         — modern technical (default)
%   tgheros    TeX Gyre Heros        — Helvetica-clone, neutral
%   source     Source Sans Pro       — humanist, friendly
%   roboto     Roboto                — Google flat-ui geometric
%   fira       Fira Sans             — Mozilla, slightly quirky
%   inter      Inter                 — UI workhorse, very tight
%   montserrat Montserrat            — wide geometric, posterish
%   lato       Lato                  — warm humanist
%   cmbright   Computer Modern Bright — light, scholarly
%   lm         Latin Modern Sans     — LaTeX default, neutral
%
\providecommand{\bnFont}{lm}

\makeatletter
\def\bn@font@plex      {\RequirePackage{plex-sans}}
\def\bn@font@tgheros   {\RequirePackage{tgheros}}
\def\bn@font@source    {\RequirePackage[default]{sourcesanspro}}
\def\bn@font@roboto    {\RequirePackage[sfdefault]{roboto}}
\def\bn@font@fira      {\RequirePackage[sfdefault]{FiraSans}}
\def\bn@font@inter     {\RequirePackage{inter}}
\def\bn@font@montserrat{\RequirePackage[defaultfam,tabular,lining]{montserrat}}
\def\bn@font@lato      {\RequirePackage[default]{lato}}
\def\bn@font@cmbright  {\RequirePackage{cmbright}}
\def\bn@font@lm        {\RequirePackage{lmodern}}
\@nameuse{bn@font@\bnFont}
\makeatother
\renewcommand{\familydefault}{\sfdefault}

% --- Palette ----------------------------------------------------------------
%
% Five role-named slots. \bnDarkMode=1 swaps to a dark variant.
% Override any individual color by \definecolor'ing it AFTER \input{_common.tex}.

\providecommand{\bnDarkMode}{0}

\ifnum\bnDarkMode=1
  % Dark variant
  \definecolor{bnInk}    {HTML}{F2EFE9} % "ink" reads light on dark
  \definecolor{bnBad}    {HTML}{F08A57} % brighter orange so it pops on dark
  \definecolor{bnPaper}  {HTML}{1B1B1B} % near-black canvas
  \definecolor{bnMargin} {HTML}{3D3833} % subtle margin guide
  \definecolor{bnMuted}  {HTML}{A39A8C} % secondary text
\else
  % Light variant (default)
  \definecolor{bnInk}    {HTML}{1B1B1B}
  \definecolor{bnBad}    {HTML}{D9633F}
  \definecolor{bnPaper}  {HTML}{FAFAF7}
  \definecolor{bnMargin} {HTML}{B5B0A5}
  \definecolor{bnMuted}  {HTML}{6F665A}
\fi

% --- Logo geometry (all in cm; tweak freely) --------------------------------

\providecommand{\bnBarCount}{4}     % number of horizontal bars
\providecommand{\bnBarWidth}{4.0}   % length of an in-margin bar
\providecommand{\bnBarHeight}{0.55} % bar thickness
\providecommand{\bnBarGap}{0.32}    % vertical gap between bars
\providecommand{\bnJutWhich}{3}     % 1..bnBarCount — which bar overflows
\providecommand{\bnJutExtra}{1.15}  % how far it sticks out
\providecommand{\bnRound}{0.07}     % corner rounding
\providecommand{\bnShowMargin}{0}   % 1 = draw the dashed margin guide, 0 = hide

% --- The mark itself --------------------------------------------------------
%
% Draws into the current tikzpicture, anchored so the bottom-left of the
% non-jutting bars is at (0,0). The jutting bar extends to the right.

\newcommand{\bnDrawMark}{%
  \pgfmathtruncatemacro{\bnLast}{\bnBarCount-1}%
  \pgfmathsetmacro{\bnTotalH}{\bnBarCount*\bnBarHeight + \bnLast*\bnBarGap}%
  \ifnum\bnShowMargin=1
    % The right-margin guide: a dashed vertical line at the in-margin width,
    % running the full height of the bar stack with a touch of overshoot.
    \draw[bnMargin, dashed, line width=0.4pt, dash pattern=on 2pt off 2pt]
    (\bnBarWidth, -0.12) -- (\bnBarWidth, \bnTotalH+0.12);
  \fi
  \foreach \i in {0,...,\bnLast}{%
      \pgfmathsetmacro{\bny}{(\bnLast-\i)*(\bnBarHeight+\bnBarGap)}%
      \pgfmathtruncatemacro{\bnNum}{\i+1}%
      \ifnum\bnNum=\bnJutWhich
        \fill[bnBad, rounded corners=\bnRound cm]
        (0,\bny) rectangle (\bnBarWidth+\bnJutExtra, \bny+\bnBarHeight);
      \else
        \fill[bnInk, rounded corners=\bnRound cm]
        (0,\bny) rectangle (\bnBarWidth, \bny+\bnBarHeight);
      \fi
    }%
}
.

\providecommand{\bnDarkMode}{0}

\ifnum\bnDarkMode=1
  % Dark variant
  \definecolor{bnInk}    {HTML}{F2EFE9} % "ink" reads light on dark
  \definecolor{bnBad}    {HTML}{F08A57} % brighter orange so it pops on dark
  \definecolor{bnPaper}  {HTML}{1B1B1B} % near-black canvas
  \definecolor{bnMargin} {HTML}{3D3833} % subtle margin guide
  \definecolor{bnMuted}  {HTML}{A39A8C} % secondary text
\else
  % Light variant (default)
  \definecolor{bnInk}    {HTML}{1B1B1B}
  \definecolor{bnBad}    {HTML}{D9633F}
  \definecolor{bnPaper}  {HTML}{FAFAF7}
  \definecolor{bnMargin} {HTML}{B5B0A5}
  \definecolor{bnMuted}  {HTML}{6F665A}
\fi

% --- Logo geometry (all in cm; tweak freely) --------------------------------

\providecommand{\bnBarCount}{4}     % number of horizontal bars
\providecommand{\bnBarWidth}{4.0}   % length of an in-margin bar
\providecommand{\bnBarHeight}{0.55} % bar thickness
\providecommand{\bnBarGap}{0.32}    % vertical gap between bars
\providecommand{\bnJutWhich}{3}     % 1..bnBarCount — which bar overflows
\providecommand{\bnJutExtra}{1.15}  % how far it sticks out
\providecommand{\bnRound}{0.07}     % corner rounding
\providecommand{\bnShowMargin}{0}   % 1 = draw the dashed margin guide, 0 = hide

% --- The mark itself --------------------------------------------------------
%
% Draws into the current tikzpicture, anchored so the bottom-left of the
% non-jutting bars is at (0,0). The jutting bar extends to the right.

\newcommand{\bnDrawMark}{%
  \pgfmathtruncatemacro{\bnLast}{\bnBarCount-1}%
  \pgfmathsetmacro{\bnTotalH}{\bnBarCount*\bnBarHeight + \bnLast*\bnBarGap}%
  \ifnum\bnShowMargin=1
    % The right-margin guide: a dashed vertical line at the in-margin width,
    % running the full height of the bar stack with a touch of overshoot.
    \draw[bnMargin, dashed, line width=0.4pt, dash pattern=on 2pt off 2pt]
    (\bnBarWidth, -0.12) -- (\bnBarWidth, \bnTotalH+0.12);
  \fi
  \foreach \i in {0,...,\bnLast}{%
      \pgfmathsetmacro{\bny}{(\bnLast-\i)*(\bnBarHeight+\bnBarGap)}%
      \pgfmathtruncatemacro{\bnNum}{\i+1}%
      \ifnum\bnNum=\bnJutWhich
        \fill[bnBad, rounded corners=\bnRound cm]
        (0,\bny) rectangle (\bnBarWidth+\bnJutExtra, \bny+\bnBarHeight);
      \else
        \fill[bnInk, rounded corners=\bnRound cm]
        (0,\bny) rectangle (\bnBarWidth, \bny+\bnBarHeight);
      \fi
    }%
}
. Available values:
%
%   plex       IBM Plex Sans         — modern technical (default)
%   tgheros    TeX Gyre Heros        — Helvetica-clone, neutral
%   source     Source Sans Pro       — humanist, friendly
%   roboto     Roboto                — Google flat-ui geometric
%   fira       Fira Sans             — Mozilla, slightly quirky
%   inter      Inter                 — UI workhorse, very tight
%   montserrat Montserrat            — wide geometric, posterish
%   lato       Lato                  — warm humanist
%   cmbright   Computer Modern Bright — light, scholarly
%   lm         Latin Modern Sans     — LaTeX default, neutral
%
\providecommand{\bnFont}{lm}

\makeatletter
\def\bn@font@plex      {\RequirePackage{plex-sans}}
\def\bn@font@tgheros   {\RequirePackage{tgheros}}
\def\bn@font@source    {\RequirePackage[default]{sourcesanspro}}
\def\bn@font@roboto    {\RequirePackage[sfdefault]{roboto}}
\def\bn@font@fira      {\RequirePackage[sfdefault]{FiraSans}}
\def\bn@font@inter     {\RequirePackage{inter}}
\def\bn@font@montserrat{\RequirePackage[defaultfam,tabular,lining]{montserrat}}
\def\bn@font@lato      {\RequirePackage[default]{lato}}
\def\bn@font@cmbright  {\RequirePackage{cmbright}}
\def\bn@font@lm        {\RequirePackage{lmodern}}
\@nameuse{bn@font@\bnFont}
\makeatother
\renewcommand{\familydefault}{\sfdefault}

% --- Palette ----------------------------------------------------------------
%
% Five role-named slots. \bnDarkMode=1 swaps to a dark variant.
% Override any individual color by \definecolor'ing it AFTER % Shared definitions for the badness brand artifacts.
%
% \input by the per-artifact .tex files (logo.tex, wordmark.tex, og.tex).
% Anything \providecommand'd here is overridable by the caller — define
% your override BEFORE % Shared definitions for the badness brand artifacts.
%
% \input by the per-artifact .tex files (logo.tex, wordmark.tex, og.tex).
% Anything \providecommand'd here is overridable by the caller — define
% your override BEFORE \input{_common.tex}.

\usepackage{tikz}
\usepackage{xcolor}

% --- Font dispatch ----------------------------------------------------------
%
% \bnFont selects the sans-serif family. Set it in your artifact file
% before \input{_common.tex}. Available values:
%
%   plex       IBM Plex Sans         — modern technical (default)
%   tgheros    TeX Gyre Heros        — Helvetica-clone, neutral
%   source     Source Sans Pro       — humanist, friendly
%   roboto     Roboto                — Google flat-ui geometric
%   fira       Fira Sans             — Mozilla, slightly quirky
%   inter      Inter                 — UI workhorse, very tight
%   montserrat Montserrat            — wide geometric, posterish
%   lato       Lato                  — warm humanist
%   cmbright   Computer Modern Bright — light, scholarly
%   lm         Latin Modern Sans     — LaTeX default, neutral
%
\providecommand{\bnFont}{lm}

\makeatletter
\def\bn@font@plex      {\RequirePackage{plex-sans}}
\def\bn@font@tgheros   {\RequirePackage{tgheros}}
\def\bn@font@source    {\RequirePackage[default]{sourcesanspro}}
\def\bn@font@roboto    {\RequirePackage[sfdefault]{roboto}}
\def\bn@font@fira      {\RequirePackage[sfdefault]{FiraSans}}
\def\bn@font@inter     {\RequirePackage{inter}}
\def\bn@font@montserrat{\RequirePackage[defaultfam,tabular,lining]{montserrat}}
\def\bn@font@lato      {\RequirePackage[default]{lato}}
\def\bn@font@cmbright  {\RequirePackage{cmbright}}
\def\bn@font@lm        {\RequirePackage{lmodern}}
\@nameuse{bn@font@\bnFont}
\makeatother
\renewcommand{\familydefault}{\sfdefault}

% --- Palette ----------------------------------------------------------------
%
% Five role-named slots. \bnDarkMode=1 swaps to a dark variant.
% Override any individual color by \definecolor'ing it AFTER \input{_common.tex}.

\providecommand{\bnDarkMode}{0}

\ifnum\bnDarkMode=1
  % Dark variant
  \definecolor{bnInk}    {HTML}{F2EFE9} % "ink" reads light on dark
  \definecolor{bnBad}    {HTML}{F08A57} % brighter orange so it pops on dark
  \definecolor{bnPaper}  {HTML}{1B1B1B} % near-black canvas
  \definecolor{bnMargin} {HTML}{3D3833} % subtle margin guide
  \definecolor{bnMuted}  {HTML}{A39A8C} % secondary text
\else
  % Light variant (default)
  \definecolor{bnInk}    {HTML}{1B1B1B}
  \definecolor{bnBad}    {HTML}{D9633F}
  \definecolor{bnPaper}  {HTML}{FAFAF7}
  \definecolor{bnMargin} {HTML}{B5B0A5}
  \definecolor{bnMuted}  {HTML}{6F665A}
\fi

% --- Logo geometry (all in cm; tweak freely) --------------------------------

\providecommand{\bnBarCount}{4}     % number of horizontal bars
\providecommand{\bnBarWidth}{4.0}   % length of an in-margin bar
\providecommand{\bnBarHeight}{0.55} % bar thickness
\providecommand{\bnBarGap}{0.32}    % vertical gap between bars
\providecommand{\bnJutWhich}{3}     % 1..bnBarCount — which bar overflows
\providecommand{\bnJutExtra}{1.15}  % how far it sticks out
\providecommand{\bnRound}{0.07}     % corner rounding
\providecommand{\bnShowMargin}{0}   % 1 = draw the dashed margin guide, 0 = hide

% --- The mark itself --------------------------------------------------------
%
% Draws into the current tikzpicture, anchored so the bottom-left of the
% non-jutting bars is at (0,0). The jutting bar extends to the right.

\newcommand{\bnDrawMark}{%
  \pgfmathtruncatemacro{\bnLast}{\bnBarCount-1}%
  \pgfmathsetmacro{\bnTotalH}{\bnBarCount*\bnBarHeight + \bnLast*\bnBarGap}%
  \ifnum\bnShowMargin=1
    % The right-margin guide: a dashed vertical line at the in-margin width,
    % running the full height of the bar stack with a touch of overshoot.
    \draw[bnMargin, dashed, line width=0.4pt, dash pattern=on 2pt off 2pt]
    (\bnBarWidth, -0.12) -- (\bnBarWidth, \bnTotalH+0.12);
  \fi
  \foreach \i in {0,...,\bnLast}{%
      \pgfmathsetmacro{\bny}{(\bnLast-\i)*(\bnBarHeight+\bnBarGap)}%
      \pgfmathtruncatemacro{\bnNum}{\i+1}%
      \ifnum\bnNum=\bnJutWhich
        \fill[bnBad, rounded corners=\bnRound cm]
        (0,\bny) rectangle (\bnBarWidth+\bnJutExtra, \bny+\bnBarHeight);
      \else
        \fill[bnInk, rounded corners=\bnRound cm]
        (0,\bny) rectangle (\bnBarWidth, \bny+\bnBarHeight);
      \fi
    }%
}
.

\usepackage{tikz}
\usepackage{xcolor}

% --- Font dispatch ----------------------------------------------------------
%
% \bnFont selects the sans-serif family. Set it in your artifact file
% before % Shared definitions for the badness brand artifacts.
%
% \input by the per-artifact .tex files (logo.tex, wordmark.tex, og.tex).
% Anything \providecommand'd here is overridable by the caller — define
% your override BEFORE \input{_common.tex}.

\usepackage{tikz}
\usepackage{xcolor}

% --- Font dispatch ----------------------------------------------------------
%
% \bnFont selects the sans-serif family. Set it in your artifact file
% before \input{_common.tex}. Available values:
%
%   plex       IBM Plex Sans         — modern technical (default)
%   tgheros    TeX Gyre Heros        — Helvetica-clone, neutral
%   source     Source Sans Pro       — humanist, friendly
%   roboto     Roboto                — Google flat-ui geometric
%   fira       Fira Sans             — Mozilla, slightly quirky
%   inter      Inter                 — UI workhorse, very tight
%   montserrat Montserrat            — wide geometric, posterish
%   lato       Lato                  — warm humanist
%   cmbright   Computer Modern Bright — light, scholarly
%   lm         Latin Modern Sans     — LaTeX default, neutral
%
\providecommand{\bnFont}{lm}

\makeatletter
\def\bn@font@plex      {\RequirePackage{plex-sans}}
\def\bn@font@tgheros   {\RequirePackage{tgheros}}
\def\bn@font@source    {\RequirePackage[default]{sourcesanspro}}
\def\bn@font@roboto    {\RequirePackage[sfdefault]{roboto}}
\def\bn@font@fira      {\RequirePackage[sfdefault]{FiraSans}}
\def\bn@font@inter     {\RequirePackage{inter}}
\def\bn@font@montserrat{\RequirePackage[defaultfam,tabular,lining]{montserrat}}
\def\bn@font@lato      {\RequirePackage[default]{lato}}
\def\bn@font@cmbright  {\RequirePackage{cmbright}}
\def\bn@font@lm        {\RequirePackage{lmodern}}
\@nameuse{bn@font@\bnFont}
\makeatother
\renewcommand{\familydefault}{\sfdefault}

% --- Palette ----------------------------------------------------------------
%
% Five role-named slots. \bnDarkMode=1 swaps to a dark variant.
% Override any individual color by \definecolor'ing it AFTER \input{_common.tex}.

\providecommand{\bnDarkMode}{0}

\ifnum\bnDarkMode=1
  % Dark variant
  \definecolor{bnInk}    {HTML}{F2EFE9} % "ink" reads light on dark
  \definecolor{bnBad}    {HTML}{F08A57} % brighter orange so it pops on dark
  \definecolor{bnPaper}  {HTML}{1B1B1B} % near-black canvas
  \definecolor{bnMargin} {HTML}{3D3833} % subtle margin guide
  \definecolor{bnMuted}  {HTML}{A39A8C} % secondary text
\else
  % Light variant (default)
  \definecolor{bnInk}    {HTML}{1B1B1B}
  \definecolor{bnBad}    {HTML}{D9633F}
  \definecolor{bnPaper}  {HTML}{FAFAF7}
  \definecolor{bnMargin} {HTML}{B5B0A5}
  \definecolor{bnMuted}  {HTML}{6F665A}
\fi

% --- Logo geometry (all in cm; tweak freely) --------------------------------

\providecommand{\bnBarCount}{4}     % number of horizontal bars
\providecommand{\bnBarWidth}{4.0}   % length of an in-margin bar
\providecommand{\bnBarHeight}{0.55} % bar thickness
\providecommand{\bnBarGap}{0.32}    % vertical gap between bars
\providecommand{\bnJutWhich}{3}     % 1..bnBarCount — which bar overflows
\providecommand{\bnJutExtra}{1.15}  % how far it sticks out
\providecommand{\bnRound}{0.07}     % corner rounding
\providecommand{\bnShowMargin}{0}   % 1 = draw the dashed margin guide, 0 = hide

% --- The mark itself --------------------------------------------------------
%
% Draws into the current tikzpicture, anchored so the bottom-left of the
% non-jutting bars is at (0,0). The jutting bar extends to the right.

\newcommand{\bnDrawMark}{%
  \pgfmathtruncatemacro{\bnLast}{\bnBarCount-1}%
  \pgfmathsetmacro{\bnTotalH}{\bnBarCount*\bnBarHeight + \bnLast*\bnBarGap}%
  \ifnum\bnShowMargin=1
    % The right-margin guide: a dashed vertical line at the in-margin width,
    % running the full height of the bar stack with a touch of overshoot.
    \draw[bnMargin, dashed, line width=0.4pt, dash pattern=on 2pt off 2pt]
    (\bnBarWidth, -0.12) -- (\bnBarWidth, \bnTotalH+0.12);
  \fi
  \foreach \i in {0,...,\bnLast}{%
      \pgfmathsetmacro{\bny}{(\bnLast-\i)*(\bnBarHeight+\bnBarGap)}%
      \pgfmathtruncatemacro{\bnNum}{\i+1}%
      \ifnum\bnNum=\bnJutWhich
        \fill[bnBad, rounded corners=\bnRound cm]
        (0,\bny) rectangle (\bnBarWidth+\bnJutExtra, \bny+\bnBarHeight);
      \else
        \fill[bnInk, rounded corners=\bnRound cm]
        (0,\bny) rectangle (\bnBarWidth, \bny+\bnBarHeight);
      \fi
    }%
}
. Available values:
%
%   plex       IBM Plex Sans         — modern technical (default)
%   tgheros    TeX Gyre Heros        — Helvetica-clone, neutral
%   source     Source Sans Pro       — humanist, friendly
%   roboto     Roboto                — Google flat-ui geometric
%   fira       Fira Sans             — Mozilla, slightly quirky
%   inter      Inter                 — UI workhorse, very tight
%   montserrat Montserrat            — wide geometric, posterish
%   lato       Lato                  — warm humanist
%   cmbright   Computer Modern Bright — light, scholarly
%   lm         Latin Modern Sans     — LaTeX default, neutral
%
\providecommand{\bnFont}{lm}

\makeatletter
\def\bn@font@plex      {\RequirePackage{plex-sans}}
\def\bn@font@tgheros   {\RequirePackage{tgheros}}
\def\bn@font@source    {\RequirePackage[default]{sourcesanspro}}
\def\bn@font@roboto    {\RequirePackage[sfdefault]{roboto}}
\def\bn@font@fira      {\RequirePackage[sfdefault]{FiraSans}}
\def\bn@font@inter     {\RequirePackage{inter}}
\def\bn@font@montserrat{\RequirePackage[defaultfam,tabular,lining]{montserrat}}
\def\bn@font@lato      {\RequirePackage[default]{lato}}
\def\bn@font@cmbright  {\RequirePackage{cmbright}}
\def\bn@font@lm        {\RequirePackage{lmodern}}
\@nameuse{bn@font@\bnFont}
\makeatother
\renewcommand{\familydefault}{\sfdefault}

% --- Palette ----------------------------------------------------------------
%
% Five role-named slots. \bnDarkMode=1 swaps to a dark variant.
% Override any individual color by \definecolor'ing it AFTER % Shared definitions for the badness brand artifacts.
%
% \input by the per-artifact .tex files (logo.tex, wordmark.tex, og.tex).
% Anything \providecommand'd here is overridable by the caller — define
% your override BEFORE \input{_common.tex}.

\usepackage{tikz}
\usepackage{xcolor}

% --- Font dispatch ----------------------------------------------------------
%
% \bnFont selects the sans-serif family. Set it in your artifact file
% before \input{_common.tex}. Available values:
%
%   plex       IBM Plex Sans         — modern technical (default)
%   tgheros    TeX Gyre Heros        — Helvetica-clone, neutral
%   source     Source Sans Pro       — humanist, friendly
%   roboto     Roboto                — Google flat-ui geometric
%   fira       Fira Sans             — Mozilla, slightly quirky
%   inter      Inter                 — UI workhorse, very tight
%   montserrat Montserrat            — wide geometric, posterish
%   lato       Lato                  — warm humanist
%   cmbright   Computer Modern Bright — light, scholarly
%   lm         Latin Modern Sans     — LaTeX default, neutral
%
\providecommand{\bnFont}{lm}

\makeatletter
\def\bn@font@plex      {\RequirePackage{plex-sans}}
\def\bn@font@tgheros   {\RequirePackage{tgheros}}
\def\bn@font@source    {\RequirePackage[default]{sourcesanspro}}
\def\bn@font@roboto    {\RequirePackage[sfdefault]{roboto}}
\def\bn@font@fira      {\RequirePackage[sfdefault]{FiraSans}}
\def\bn@font@inter     {\RequirePackage{inter}}
\def\bn@font@montserrat{\RequirePackage[defaultfam,tabular,lining]{montserrat}}
\def\bn@font@lato      {\RequirePackage[default]{lato}}
\def\bn@font@cmbright  {\RequirePackage{cmbright}}
\def\bn@font@lm        {\RequirePackage{lmodern}}
\@nameuse{bn@font@\bnFont}
\makeatother
\renewcommand{\familydefault}{\sfdefault}

% --- Palette ----------------------------------------------------------------
%
% Five role-named slots. \bnDarkMode=1 swaps to a dark variant.
% Override any individual color by \definecolor'ing it AFTER \input{_common.tex}.

\providecommand{\bnDarkMode}{0}

\ifnum\bnDarkMode=1
  % Dark variant
  \definecolor{bnInk}    {HTML}{F2EFE9} % "ink" reads light on dark
  \definecolor{bnBad}    {HTML}{F08A57} % brighter orange so it pops on dark
  \definecolor{bnPaper}  {HTML}{1B1B1B} % near-black canvas
  \definecolor{bnMargin} {HTML}{3D3833} % subtle margin guide
  \definecolor{bnMuted}  {HTML}{A39A8C} % secondary text
\else
  % Light variant (default)
  \definecolor{bnInk}    {HTML}{1B1B1B}
  \definecolor{bnBad}    {HTML}{D9633F}
  \definecolor{bnPaper}  {HTML}{FAFAF7}
  \definecolor{bnMargin} {HTML}{B5B0A5}
  \definecolor{bnMuted}  {HTML}{6F665A}
\fi

% --- Logo geometry (all in cm; tweak freely) --------------------------------

\providecommand{\bnBarCount}{4}     % number of horizontal bars
\providecommand{\bnBarWidth}{4.0}   % length of an in-margin bar
\providecommand{\bnBarHeight}{0.55} % bar thickness
\providecommand{\bnBarGap}{0.32}    % vertical gap between bars
\providecommand{\bnJutWhich}{3}     % 1..bnBarCount — which bar overflows
\providecommand{\bnJutExtra}{1.15}  % how far it sticks out
\providecommand{\bnRound}{0.07}     % corner rounding
\providecommand{\bnShowMargin}{0}   % 1 = draw the dashed margin guide, 0 = hide

% --- The mark itself --------------------------------------------------------
%
% Draws into the current tikzpicture, anchored so the bottom-left of the
% non-jutting bars is at (0,0). The jutting bar extends to the right.

\newcommand{\bnDrawMark}{%
  \pgfmathtruncatemacro{\bnLast}{\bnBarCount-1}%
  \pgfmathsetmacro{\bnTotalH}{\bnBarCount*\bnBarHeight + \bnLast*\bnBarGap}%
  \ifnum\bnShowMargin=1
    % The right-margin guide: a dashed vertical line at the in-margin width,
    % running the full height of the bar stack with a touch of overshoot.
    \draw[bnMargin, dashed, line width=0.4pt, dash pattern=on 2pt off 2pt]
    (\bnBarWidth, -0.12) -- (\bnBarWidth, \bnTotalH+0.12);
  \fi
  \foreach \i in {0,...,\bnLast}{%
      \pgfmathsetmacro{\bny}{(\bnLast-\i)*(\bnBarHeight+\bnBarGap)}%
      \pgfmathtruncatemacro{\bnNum}{\i+1}%
      \ifnum\bnNum=\bnJutWhich
        \fill[bnBad, rounded corners=\bnRound cm]
        (0,\bny) rectangle (\bnBarWidth+\bnJutExtra, \bny+\bnBarHeight);
      \else
        \fill[bnInk, rounded corners=\bnRound cm]
        (0,\bny) rectangle (\bnBarWidth, \bny+\bnBarHeight);
      \fi
    }%
}
.

\providecommand{\bnDarkMode}{0}

\ifnum\bnDarkMode=1
  % Dark variant
  \definecolor{bnInk}    {HTML}{F2EFE9} % "ink" reads light on dark
  \definecolor{bnBad}    {HTML}{F08A57} % brighter orange so it pops on dark
  \definecolor{bnPaper}  {HTML}{1B1B1B} % near-black canvas
  \definecolor{bnMargin} {HTML}{3D3833} % subtle margin guide
  \definecolor{bnMuted}  {HTML}{A39A8C} % secondary text
\else
  % Light variant (default)
  \definecolor{bnInk}    {HTML}{1B1B1B}
  \definecolor{bnBad}    {HTML}{D9633F}
  \definecolor{bnPaper}  {HTML}{FAFAF7}
  \definecolor{bnMargin} {HTML}{B5B0A5}
  \definecolor{bnMuted}  {HTML}{6F665A}
\fi

% --- Logo geometry (all in cm; tweak freely) --------------------------------

\providecommand{\bnBarCount}{4}     % number of horizontal bars
\providecommand{\bnBarWidth}{4.0}   % length of an in-margin bar
\providecommand{\bnBarHeight}{0.55} % bar thickness
\providecommand{\bnBarGap}{0.32}    % vertical gap between bars
\providecommand{\bnJutWhich}{3}     % 1..bnBarCount — which bar overflows
\providecommand{\bnJutExtra}{1.15}  % how far it sticks out
\providecommand{\bnRound}{0.07}     % corner rounding
\providecommand{\bnShowMargin}{0}   % 1 = draw the dashed margin guide, 0 = hide

% --- The mark itself --------------------------------------------------------
%
% Draws into the current tikzpicture, anchored so the bottom-left of the
% non-jutting bars is at (0,0). The jutting bar extends to the right.

\newcommand{\bnDrawMark}{%
  \pgfmathtruncatemacro{\bnLast}{\bnBarCount-1}%
  \pgfmathsetmacro{\bnTotalH}{\bnBarCount*\bnBarHeight + \bnLast*\bnBarGap}%
  \ifnum\bnShowMargin=1
    % The right-margin guide: a dashed vertical line at the in-margin width,
    % running the full height of the bar stack with a touch of overshoot.
    \draw[bnMargin, dashed, line width=0.4pt, dash pattern=on 2pt off 2pt]
    (\bnBarWidth, -0.12) -- (\bnBarWidth, \bnTotalH+0.12);
  \fi
  \foreach \i in {0,...,\bnLast}{%
      \pgfmathsetmacro{\bny}{(\bnLast-\i)*(\bnBarHeight+\bnBarGap)}%
      \pgfmathtruncatemacro{\bnNum}{\i+1}%
      \ifnum\bnNum=\bnJutWhich
        \fill[bnBad, rounded corners=\bnRound cm]
        (0,\bny) rectangle (\bnBarWidth+\bnJutExtra, \bny+\bnBarHeight);
      \else
        \fill[bnInk, rounded corners=\bnRound cm]
        (0,\bny) rectangle (\bnBarWidth, \bny+\bnBarHeight);
      \fi
    }%
}
.

\providecommand{\bnDarkMode}{0}

\ifnum\bnDarkMode=1
  % Dark variant
  \definecolor{bnInk}    {HTML}{F2EFE9} % "ink" reads light on dark
  \definecolor{bnBad}    {HTML}{F08A57} % brighter orange so it pops on dark
  \definecolor{bnPaper}  {HTML}{1B1B1B} % near-black canvas
  \definecolor{bnMargin} {HTML}{3D3833} % subtle margin guide
  \definecolor{bnMuted}  {HTML}{A39A8C} % secondary text
\else
  % Light variant (default)
  \definecolor{bnInk}    {HTML}{1B1B1B}
  \definecolor{bnBad}    {HTML}{D9633F}
  \definecolor{bnPaper}  {HTML}{FAFAF7}
  \definecolor{bnMargin} {HTML}{B5B0A5}
  \definecolor{bnMuted}  {HTML}{6F665A}
\fi

% --- Logo geometry (all in cm; tweak freely) --------------------------------

\providecommand{\bnBarCount}{4}     % number of horizontal bars
\providecommand{\bnBarWidth}{4.0}   % length of an in-margin bar
\providecommand{\bnBarHeight}{0.55} % bar thickness
\providecommand{\bnBarGap}{0.32}    % vertical gap between bars
\providecommand{\bnJutWhich}{3}     % 1..bnBarCount — which bar overflows
\providecommand{\bnJutExtra}{1.15}  % how far it sticks out
\providecommand{\bnRound}{0.07}     % corner rounding
\providecommand{\bnShowMargin}{0}   % 1 = draw the dashed margin guide, 0 = hide

% --- The mark itself --------------------------------------------------------
%
% Draws into the current tikzpicture, anchored so the bottom-left of the
% non-jutting bars is at (0,0). The jutting bar extends to the right.

\newcommand{\bnDrawMark}{%
  \pgfmathtruncatemacro{\bnLast}{\bnBarCount-1}%
  \pgfmathsetmacro{\bnTotalH}{\bnBarCount*\bnBarHeight + \bnLast*\bnBarGap}%
  \ifnum\bnShowMargin=1
    % The right-margin guide: a dashed vertical line at the in-margin width,
    % running the full height of the bar stack with a touch of overshoot.
    \draw[bnMargin, dashed, line width=0.4pt, dash pattern=on 2pt off 2pt]
    (\bnBarWidth, -0.12) -- (\bnBarWidth, \bnTotalH+0.12);
  \fi
  \foreach \i in {0,...,\bnLast}{%
      \pgfmathsetmacro{\bny}{(\bnLast-\i)*(\bnBarHeight+\bnBarGap)}%
      \pgfmathtruncatemacro{\bnNum}{\i+1}%
      \ifnum\bnNum=\bnJutWhich
        \fill[bnBad, rounded corners=\bnRound cm]
        (0,\bny) rectangle (\bnBarWidth+\bnJutExtra, \bny+\bnBarHeight);
      \else
        \fill[bnInk, rounded corners=\bnRound cm]
        (0,\bny) rectangle (\bnBarWidth, \bny+\bnBarHeight);
      \fi
    }%
}
.

\usepackage{tikz}
\usepackage{xcolor}

% --- Font dispatch ----------------------------------------------------------
%
% \bnFont selects the sans-serif family. Set it in your artifact file
% before % Shared definitions for the badness brand artifacts.
%
% \input by the per-artifact .tex files (logo.tex, wordmark.tex, og.tex).
% Anything \providecommand'd here is overridable by the caller — define
% your override BEFORE % Shared definitions for the badness brand artifacts.
%
% \input by the per-artifact .tex files (logo.tex, wordmark.tex, og.tex).
% Anything \providecommand'd here is overridable by the caller — define
% your override BEFORE % Shared definitions for the badness brand artifacts.
%
% \input by the per-artifact .tex files (logo.tex, wordmark.tex, og.tex).
% Anything \providecommand'd here is overridable by the caller — define
% your override BEFORE \input{_common.tex}.

\usepackage{tikz}
\usepackage{xcolor}

% --- Font dispatch ----------------------------------------------------------
%
% \bnFont selects the sans-serif family. Set it in your artifact file
% before \input{_common.tex}. Available values:
%
%   plex       IBM Plex Sans         — modern technical (default)
%   tgheros    TeX Gyre Heros        — Helvetica-clone, neutral
%   source     Source Sans Pro       — humanist, friendly
%   roboto     Roboto                — Google flat-ui geometric
%   fira       Fira Sans             — Mozilla, slightly quirky
%   inter      Inter                 — UI workhorse, very tight
%   montserrat Montserrat            — wide geometric, posterish
%   lato       Lato                  — warm humanist
%   cmbright   Computer Modern Bright — light, scholarly
%   lm         Latin Modern Sans     — LaTeX default, neutral
%
\providecommand{\bnFont}{lm}

\makeatletter
\def\bn@font@plex      {\RequirePackage{plex-sans}}
\def\bn@font@tgheros   {\RequirePackage{tgheros}}
\def\bn@font@source    {\RequirePackage[default]{sourcesanspro}}
\def\bn@font@roboto    {\RequirePackage[sfdefault]{roboto}}
\def\bn@font@fira      {\RequirePackage[sfdefault]{FiraSans}}
\def\bn@font@inter     {\RequirePackage{inter}}
\def\bn@font@montserrat{\RequirePackage[defaultfam,tabular,lining]{montserrat}}
\def\bn@font@lato      {\RequirePackage[default]{lato}}
\def\bn@font@cmbright  {\RequirePackage{cmbright}}
\def\bn@font@lm        {\RequirePackage{lmodern}}
\@nameuse{bn@font@\bnFont}
\makeatother
\renewcommand{\familydefault}{\sfdefault}

% --- Palette ----------------------------------------------------------------
%
% Five role-named slots. \bnDarkMode=1 swaps to a dark variant.
% Override any individual color by \definecolor'ing it AFTER \input{_common.tex}.

\providecommand{\bnDarkMode}{0}

\ifnum\bnDarkMode=1
  % Dark variant
  \definecolor{bnInk}    {HTML}{F2EFE9} % "ink" reads light on dark
  \definecolor{bnBad}    {HTML}{F08A57} % brighter orange so it pops on dark
  \definecolor{bnPaper}  {HTML}{1B1B1B} % near-black canvas
  \definecolor{bnMargin} {HTML}{3D3833} % subtle margin guide
  \definecolor{bnMuted}  {HTML}{A39A8C} % secondary text
\else
  % Light variant (default)
  \definecolor{bnInk}    {HTML}{1B1B1B}
  \definecolor{bnBad}    {HTML}{D9633F}
  \definecolor{bnPaper}  {HTML}{FAFAF7}
  \definecolor{bnMargin} {HTML}{B5B0A5}
  \definecolor{bnMuted}  {HTML}{6F665A}
\fi

% --- Logo geometry (all in cm; tweak freely) --------------------------------

\providecommand{\bnBarCount}{4}     % number of horizontal bars
\providecommand{\bnBarWidth}{4.0}   % length of an in-margin bar
\providecommand{\bnBarHeight}{0.55} % bar thickness
\providecommand{\bnBarGap}{0.32}    % vertical gap between bars
\providecommand{\bnJutWhich}{3}     % 1..bnBarCount — which bar overflows
\providecommand{\bnJutExtra}{1.15}  % how far it sticks out
\providecommand{\bnRound}{0.07}     % corner rounding
\providecommand{\bnShowMargin}{0}   % 1 = draw the dashed margin guide, 0 = hide

% --- The mark itself --------------------------------------------------------
%
% Draws into the current tikzpicture, anchored so the bottom-left of the
% non-jutting bars is at (0,0). The jutting bar extends to the right.

\newcommand{\bnDrawMark}{%
  \pgfmathtruncatemacro{\bnLast}{\bnBarCount-1}%
  \pgfmathsetmacro{\bnTotalH}{\bnBarCount*\bnBarHeight + \bnLast*\bnBarGap}%
  \ifnum\bnShowMargin=1
    % The right-margin guide: a dashed vertical line at the in-margin width,
    % running the full height of the bar stack with a touch of overshoot.
    \draw[bnMargin, dashed, line width=0.4pt, dash pattern=on 2pt off 2pt]
    (\bnBarWidth, -0.12) -- (\bnBarWidth, \bnTotalH+0.12);
  \fi
  \foreach \i in {0,...,\bnLast}{%
      \pgfmathsetmacro{\bny}{(\bnLast-\i)*(\bnBarHeight+\bnBarGap)}%
      \pgfmathtruncatemacro{\bnNum}{\i+1}%
      \ifnum\bnNum=\bnJutWhich
        \fill[bnBad, rounded corners=\bnRound cm]
        (0,\bny) rectangle (\bnBarWidth+\bnJutExtra, \bny+\bnBarHeight);
      \else
        \fill[bnInk, rounded corners=\bnRound cm]
        (0,\bny) rectangle (\bnBarWidth, \bny+\bnBarHeight);
      \fi
    }%
}
.

\usepackage{tikz}
\usepackage{xcolor}

% --- Font dispatch ----------------------------------------------------------
%
% \bnFont selects the sans-serif family. Set it in your artifact file
% before % Shared definitions for the badness brand artifacts.
%
% \input by the per-artifact .tex files (logo.tex, wordmark.tex, og.tex).
% Anything \providecommand'd here is overridable by the caller — define
% your override BEFORE \input{_common.tex}.

\usepackage{tikz}
\usepackage{xcolor}

% --- Font dispatch ----------------------------------------------------------
%
% \bnFont selects the sans-serif family. Set it in your artifact file
% before \input{_common.tex}. Available values:
%
%   plex       IBM Plex Sans         — modern technical (default)
%   tgheros    TeX Gyre Heros        — Helvetica-clone, neutral
%   source     Source Sans Pro       — humanist, friendly
%   roboto     Roboto                — Google flat-ui geometric
%   fira       Fira Sans             — Mozilla, slightly quirky
%   inter      Inter                 — UI workhorse, very tight
%   montserrat Montserrat            — wide geometric, posterish
%   lato       Lato                  — warm humanist
%   cmbright   Computer Modern Bright — light, scholarly
%   lm         Latin Modern Sans     — LaTeX default, neutral
%
\providecommand{\bnFont}{lm}

\makeatletter
\def\bn@font@plex      {\RequirePackage{plex-sans}}
\def\bn@font@tgheros   {\RequirePackage{tgheros}}
\def\bn@font@source    {\RequirePackage[default]{sourcesanspro}}
\def\bn@font@roboto    {\RequirePackage[sfdefault]{roboto}}
\def\bn@font@fira      {\RequirePackage[sfdefault]{FiraSans}}
\def\bn@font@inter     {\RequirePackage{inter}}
\def\bn@font@montserrat{\RequirePackage[defaultfam,tabular,lining]{montserrat}}
\def\bn@font@lato      {\RequirePackage[default]{lato}}
\def\bn@font@cmbright  {\RequirePackage{cmbright}}
\def\bn@font@lm        {\RequirePackage{lmodern}}
\@nameuse{bn@font@\bnFont}
\makeatother
\renewcommand{\familydefault}{\sfdefault}

% --- Palette ----------------------------------------------------------------
%
% Five role-named slots. \bnDarkMode=1 swaps to a dark variant.
% Override any individual color by \definecolor'ing it AFTER \input{_common.tex}.

\providecommand{\bnDarkMode}{0}

\ifnum\bnDarkMode=1
  % Dark variant
  \definecolor{bnInk}    {HTML}{F2EFE9} % "ink" reads light on dark
  \definecolor{bnBad}    {HTML}{F08A57} % brighter orange so it pops on dark
  \definecolor{bnPaper}  {HTML}{1B1B1B} % near-black canvas
  \definecolor{bnMargin} {HTML}{3D3833} % subtle margin guide
  \definecolor{bnMuted}  {HTML}{A39A8C} % secondary text
\else
  % Light variant (default)
  \definecolor{bnInk}    {HTML}{1B1B1B}
  \definecolor{bnBad}    {HTML}{D9633F}
  \definecolor{bnPaper}  {HTML}{FAFAF7}
  \definecolor{bnMargin} {HTML}{B5B0A5}
  \definecolor{bnMuted}  {HTML}{6F665A}
\fi

% --- Logo geometry (all in cm; tweak freely) --------------------------------

\providecommand{\bnBarCount}{4}     % number of horizontal bars
\providecommand{\bnBarWidth}{4.0}   % length of an in-margin bar
\providecommand{\bnBarHeight}{0.55} % bar thickness
\providecommand{\bnBarGap}{0.32}    % vertical gap between bars
\providecommand{\bnJutWhich}{3}     % 1..bnBarCount — which bar overflows
\providecommand{\bnJutExtra}{1.15}  % how far it sticks out
\providecommand{\bnRound}{0.07}     % corner rounding
\providecommand{\bnShowMargin}{0}   % 1 = draw the dashed margin guide, 0 = hide

% --- The mark itself --------------------------------------------------------
%
% Draws into the current tikzpicture, anchored so the bottom-left of the
% non-jutting bars is at (0,0). The jutting bar extends to the right.

\newcommand{\bnDrawMark}{%
  \pgfmathtruncatemacro{\bnLast}{\bnBarCount-1}%
  \pgfmathsetmacro{\bnTotalH}{\bnBarCount*\bnBarHeight + \bnLast*\bnBarGap}%
  \ifnum\bnShowMargin=1
    % The right-margin guide: a dashed vertical line at the in-margin width,
    % running the full height of the bar stack with a touch of overshoot.
    \draw[bnMargin, dashed, line width=0.4pt, dash pattern=on 2pt off 2pt]
    (\bnBarWidth, -0.12) -- (\bnBarWidth, \bnTotalH+0.12);
  \fi
  \foreach \i in {0,...,\bnLast}{%
      \pgfmathsetmacro{\bny}{(\bnLast-\i)*(\bnBarHeight+\bnBarGap)}%
      \pgfmathtruncatemacro{\bnNum}{\i+1}%
      \ifnum\bnNum=\bnJutWhich
        \fill[bnBad, rounded corners=\bnRound cm]
        (0,\bny) rectangle (\bnBarWidth+\bnJutExtra, \bny+\bnBarHeight);
      \else
        \fill[bnInk, rounded corners=\bnRound cm]
        (0,\bny) rectangle (\bnBarWidth, \bny+\bnBarHeight);
      \fi
    }%
}
. Available values:
%
%   plex       IBM Plex Sans         — modern technical (default)
%   tgheros    TeX Gyre Heros        — Helvetica-clone, neutral
%   source     Source Sans Pro       — humanist, friendly
%   roboto     Roboto                — Google flat-ui geometric
%   fira       Fira Sans             — Mozilla, slightly quirky
%   inter      Inter                 — UI workhorse, very tight
%   montserrat Montserrat            — wide geometric, posterish
%   lato       Lato                  — warm humanist
%   cmbright   Computer Modern Bright — light, scholarly
%   lm         Latin Modern Sans     — LaTeX default, neutral
%
\providecommand{\bnFont}{lm}

\makeatletter
\def\bn@font@plex      {\RequirePackage{plex-sans}}
\def\bn@font@tgheros   {\RequirePackage{tgheros}}
\def\bn@font@source    {\RequirePackage[default]{sourcesanspro}}
\def\bn@font@roboto    {\RequirePackage[sfdefault]{roboto}}
\def\bn@font@fira      {\RequirePackage[sfdefault]{FiraSans}}
\def\bn@font@inter     {\RequirePackage{inter}}
\def\bn@font@montserrat{\RequirePackage[defaultfam,tabular,lining]{montserrat}}
\def\bn@font@lato      {\RequirePackage[default]{lato}}
\def\bn@font@cmbright  {\RequirePackage{cmbright}}
\def\bn@font@lm        {\RequirePackage{lmodern}}
\@nameuse{bn@font@\bnFont}
\makeatother
\renewcommand{\familydefault}{\sfdefault}

% --- Palette ----------------------------------------------------------------
%
% Five role-named slots. \bnDarkMode=1 swaps to a dark variant.
% Override any individual color by \definecolor'ing it AFTER % Shared definitions for the badness brand artifacts.
%
% \input by the per-artifact .tex files (logo.tex, wordmark.tex, og.tex).
% Anything \providecommand'd here is overridable by the caller — define
% your override BEFORE \input{_common.tex}.

\usepackage{tikz}
\usepackage{xcolor}

% --- Font dispatch ----------------------------------------------------------
%
% \bnFont selects the sans-serif family. Set it in your artifact file
% before \input{_common.tex}. Available values:
%
%   plex       IBM Plex Sans         — modern technical (default)
%   tgheros    TeX Gyre Heros        — Helvetica-clone, neutral
%   source     Source Sans Pro       — humanist, friendly
%   roboto     Roboto                — Google flat-ui geometric
%   fira       Fira Sans             — Mozilla, slightly quirky
%   inter      Inter                 — UI workhorse, very tight
%   montserrat Montserrat            — wide geometric, posterish
%   lato       Lato                  — warm humanist
%   cmbright   Computer Modern Bright — light, scholarly
%   lm         Latin Modern Sans     — LaTeX default, neutral
%
\providecommand{\bnFont}{lm}

\makeatletter
\def\bn@font@plex      {\RequirePackage{plex-sans}}
\def\bn@font@tgheros   {\RequirePackage{tgheros}}
\def\bn@font@source    {\RequirePackage[default]{sourcesanspro}}
\def\bn@font@roboto    {\RequirePackage[sfdefault]{roboto}}
\def\bn@font@fira      {\RequirePackage[sfdefault]{FiraSans}}
\def\bn@font@inter     {\RequirePackage{inter}}
\def\bn@font@montserrat{\RequirePackage[defaultfam,tabular,lining]{montserrat}}
\def\bn@font@lato      {\RequirePackage[default]{lato}}
\def\bn@font@cmbright  {\RequirePackage{cmbright}}
\def\bn@font@lm        {\RequirePackage{lmodern}}
\@nameuse{bn@font@\bnFont}
\makeatother
\renewcommand{\familydefault}{\sfdefault}

% --- Palette ----------------------------------------------------------------
%
% Five role-named slots. \bnDarkMode=1 swaps to a dark variant.
% Override any individual color by \definecolor'ing it AFTER \input{_common.tex}.

\providecommand{\bnDarkMode}{0}

\ifnum\bnDarkMode=1
  % Dark variant
  \definecolor{bnInk}    {HTML}{F2EFE9} % "ink" reads light on dark
  \definecolor{bnBad}    {HTML}{F08A57} % brighter orange so it pops on dark
  \definecolor{bnPaper}  {HTML}{1B1B1B} % near-black canvas
  \definecolor{bnMargin} {HTML}{3D3833} % subtle margin guide
  \definecolor{bnMuted}  {HTML}{A39A8C} % secondary text
\else
  % Light variant (default)
  \definecolor{bnInk}    {HTML}{1B1B1B}
  \definecolor{bnBad}    {HTML}{D9633F}
  \definecolor{bnPaper}  {HTML}{FAFAF7}
  \definecolor{bnMargin} {HTML}{B5B0A5}
  \definecolor{bnMuted}  {HTML}{6F665A}
\fi

% --- Logo geometry (all in cm; tweak freely) --------------------------------

\providecommand{\bnBarCount}{4}     % number of horizontal bars
\providecommand{\bnBarWidth}{4.0}   % length of an in-margin bar
\providecommand{\bnBarHeight}{0.55} % bar thickness
\providecommand{\bnBarGap}{0.32}    % vertical gap between bars
\providecommand{\bnJutWhich}{3}     % 1..bnBarCount — which bar overflows
\providecommand{\bnJutExtra}{1.15}  % how far it sticks out
\providecommand{\bnRound}{0.07}     % corner rounding
\providecommand{\bnShowMargin}{0}   % 1 = draw the dashed margin guide, 0 = hide

% --- The mark itself --------------------------------------------------------
%
% Draws into the current tikzpicture, anchored so the bottom-left of the
% non-jutting bars is at (0,0). The jutting bar extends to the right.

\newcommand{\bnDrawMark}{%
  \pgfmathtruncatemacro{\bnLast}{\bnBarCount-1}%
  \pgfmathsetmacro{\bnTotalH}{\bnBarCount*\bnBarHeight + \bnLast*\bnBarGap}%
  \ifnum\bnShowMargin=1
    % The right-margin guide: a dashed vertical line at the in-margin width,
    % running the full height of the bar stack with a touch of overshoot.
    \draw[bnMargin, dashed, line width=0.4pt, dash pattern=on 2pt off 2pt]
    (\bnBarWidth, -0.12) -- (\bnBarWidth, \bnTotalH+0.12);
  \fi
  \foreach \i in {0,...,\bnLast}{%
      \pgfmathsetmacro{\bny}{(\bnLast-\i)*(\bnBarHeight+\bnBarGap)}%
      \pgfmathtruncatemacro{\bnNum}{\i+1}%
      \ifnum\bnNum=\bnJutWhich
        \fill[bnBad, rounded corners=\bnRound cm]
        (0,\bny) rectangle (\bnBarWidth+\bnJutExtra, \bny+\bnBarHeight);
      \else
        \fill[bnInk, rounded corners=\bnRound cm]
        (0,\bny) rectangle (\bnBarWidth, \bny+\bnBarHeight);
      \fi
    }%
}
.

\providecommand{\bnDarkMode}{0}

\ifnum\bnDarkMode=1
  % Dark variant
  \definecolor{bnInk}    {HTML}{F2EFE9} % "ink" reads light on dark
  \definecolor{bnBad}    {HTML}{F08A57} % brighter orange so it pops on dark
  \definecolor{bnPaper}  {HTML}{1B1B1B} % near-black canvas
  \definecolor{bnMargin} {HTML}{3D3833} % subtle margin guide
  \definecolor{bnMuted}  {HTML}{A39A8C} % secondary text
\else
  % Light variant (default)
  \definecolor{bnInk}    {HTML}{1B1B1B}
  \definecolor{bnBad}    {HTML}{D9633F}
  \definecolor{bnPaper}  {HTML}{FAFAF7}
  \definecolor{bnMargin} {HTML}{B5B0A5}
  \definecolor{bnMuted}  {HTML}{6F665A}
\fi

% --- Logo geometry (all in cm; tweak freely) --------------------------------

\providecommand{\bnBarCount}{4}     % number of horizontal bars
\providecommand{\bnBarWidth}{4.0}   % length of an in-margin bar
\providecommand{\bnBarHeight}{0.55} % bar thickness
\providecommand{\bnBarGap}{0.32}    % vertical gap between bars
\providecommand{\bnJutWhich}{3}     % 1..bnBarCount — which bar overflows
\providecommand{\bnJutExtra}{1.15}  % how far it sticks out
\providecommand{\bnRound}{0.07}     % corner rounding
\providecommand{\bnShowMargin}{0}   % 1 = draw the dashed margin guide, 0 = hide

% --- The mark itself --------------------------------------------------------
%
% Draws into the current tikzpicture, anchored so the bottom-left of the
% non-jutting bars is at (0,0). The jutting bar extends to the right.

\newcommand{\bnDrawMark}{%
  \pgfmathtruncatemacro{\bnLast}{\bnBarCount-1}%
  \pgfmathsetmacro{\bnTotalH}{\bnBarCount*\bnBarHeight + \bnLast*\bnBarGap}%
  \ifnum\bnShowMargin=1
    % The right-margin guide: a dashed vertical line at the in-margin width,
    % running the full height of the bar stack with a touch of overshoot.
    \draw[bnMargin, dashed, line width=0.4pt, dash pattern=on 2pt off 2pt]
    (\bnBarWidth, -0.12) -- (\bnBarWidth, \bnTotalH+0.12);
  \fi
  \foreach \i in {0,...,\bnLast}{%
      \pgfmathsetmacro{\bny}{(\bnLast-\i)*(\bnBarHeight+\bnBarGap)}%
      \pgfmathtruncatemacro{\bnNum}{\i+1}%
      \ifnum\bnNum=\bnJutWhich
        \fill[bnBad, rounded corners=\bnRound cm]
        (0,\bny) rectangle (\bnBarWidth+\bnJutExtra, \bny+\bnBarHeight);
      \else
        \fill[bnInk, rounded corners=\bnRound cm]
        (0,\bny) rectangle (\bnBarWidth, \bny+\bnBarHeight);
      \fi
    }%
}
.

\usepackage{tikz}
\usepackage{xcolor}

% --- Font dispatch ----------------------------------------------------------
%
% \bnFont selects the sans-serif family. Set it in your artifact file
% before % Shared definitions for the badness brand artifacts.
%
% \input by the per-artifact .tex files (logo.tex, wordmark.tex, og.tex).
% Anything \providecommand'd here is overridable by the caller — define
% your override BEFORE % Shared definitions for the badness brand artifacts.
%
% \input by the per-artifact .tex files (logo.tex, wordmark.tex, og.tex).
% Anything \providecommand'd here is overridable by the caller — define
% your override BEFORE \input{_common.tex}.

\usepackage{tikz}
\usepackage{xcolor}

% --- Font dispatch ----------------------------------------------------------
%
% \bnFont selects the sans-serif family. Set it in your artifact file
% before \input{_common.tex}. Available values:
%
%   plex       IBM Plex Sans         — modern technical (default)
%   tgheros    TeX Gyre Heros        — Helvetica-clone, neutral
%   source     Source Sans Pro       — humanist, friendly
%   roboto     Roboto                — Google flat-ui geometric
%   fira       Fira Sans             — Mozilla, slightly quirky
%   inter      Inter                 — UI workhorse, very tight
%   montserrat Montserrat            — wide geometric, posterish
%   lato       Lato                  — warm humanist
%   cmbright   Computer Modern Bright — light, scholarly
%   lm         Latin Modern Sans     — LaTeX default, neutral
%
\providecommand{\bnFont}{lm}

\makeatletter
\def\bn@font@plex      {\RequirePackage{plex-sans}}
\def\bn@font@tgheros   {\RequirePackage{tgheros}}
\def\bn@font@source    {\RequirePackage[default]{sourcesanspro}}
\def\bn@font@roboto    {\RequirePackage[sfdefault]{roboto}}
\def\bn@font@fira      {\RequirePackage[sfdefault]{FiraSans}}
\def\bn@font@inter     {\RequirePackage{inter}}
\def\bn@font@montserrat{\RequirePackage[defaultfam,tabular,lining]{montserrat}}
\def\bn@font@lato      {\RequirePackage[default]{lato}}
\def\bn@font@cmbright  {\RequirePackage{cmbright}}
\def\bn@font@lm        {\RequirePackage{lmodern}}
\@nameuse{bn@font@\bnFont}
\makeatother
\renewcommand{\familydefault}{\sfdefault}

% --- Palette ----------------------------------------------------------------
%
% Five role-named slots. \bnDarkMode=1 swaps to a dark variant.
% Override any individual color by \definecolor'ing it AFTER \input{_common.tex}.

\providecommand{\bnDarkMode}{0}

\ifnum\bnDarkMode=1
  % Dark variant
  \definecolor{bnInk}    {HTML}{F2EFE9} % "ink" reads light on dark
  \definecolor{bnBad}    {HTML}{F08A57} % brighter orange so it pops on dark
  \definecolor{bnPaper}  {HTML}{1B1B1B} % near-black canvas
  \definecolor{bnMargin} {HTML}{3D3833} % subtle margin guide
  \definecolor{bnMuted}  {HTML}{A39A8C} % secondary text
\else
  % Light variant (default)
  \definecolor{bnInk}    {HTML}{1B1B1B}
  \definecolor{bnBad}    {HTML}{D9633F}
  \definecolor{bnPaper}  {HTML}{FAFAF7}
  \definecolor{bnMargin} {HTML}{B5B0A5}
  \definecolor{bnMuted}  {HTML}{6F665A}
\fi

% --- Logo geometry (all in cm; tweak freely) --------------------------------

\providecommand{\bnBarCount}{4}     % number of horizontal bars
\providecommand{\bnBarWidth}{4.0}   % length of an in-margin bar
\providecommand{\bnBarHeight}{0.55} % bar thickness
\providecommand{\bnBarGap}{0.32}    % vertical gap between bars
\providecommand{\bnJutWhich}{3}     % 1..bnBarCount — which bar overflows
\providecommand{\bnJutExtra}{1.15}  % how far it sticks out
\providecommand{\bnRound}{0.07}     % corner rounding
\providecommand{\bnShowMargin}{0}   % 1 = draw the dashed margin guide, 0 = hide

% --- The mark itself --------------------------------------------------------
%
% Draws into the current tikzpicture, anchored so the bottom-left of the
% non-jutting bars is at (0,0). The jutting bar extends to the right.

\newcommand{\bnDrawMark}{%
  \pgfmathtruncatemacro{\bnLast}{\bnBarCount-1}%
  \pgfmathsetmacro{\bnTotalH}{\bnBarCount*\bnBarHeight + \bnLast*\bnBarGap}%
  \ifnum\bnShowMargin=1
    % The right-margin guide: a dashed vertical line at the in-margin width,
    % running the full height of the bar stack with a touch of overshoot.
    \draw[bnMargin, dashed, line width=0.4pt, dash pattern=on 2pt off 2pt]
    (\bnBarWidth, -0.12) -- (\bnBarWidth, \bnTotalH+0.12);
  \fi
  \foreach \i in {0,...,\bnLast}{%
      \pgfmathsetmacro{\bny}{(\bnLast-\i)*(\bnBarHeight+\bnBarGap)}%
      \pgfmathtruncatemacro{\bnNum}{\i+1}%
      \ifnum\bnNum=\bnJutWhich
        \fill[bnBad, rounded corners=\bnRound cm]
        (0,\bny) rectangle (\bnBarWidth+\bnJutExtra, \bny+\bnBarHeight);
      \else
        \fill[bnInk, rounded corners=\bnRound cm]
        (0,\bny) rectangle (\bnBarWidth, \bny+\bnBarHeight);
      \fi
    }%
}
.

\usepackage{tikz}
\usepackage{xcolor}

% --- Font dispatch ----------------------------------------------------------
%
% \bnFont selects the sans-serif family. Set it in your artifact file
% before % Shared definitions for the badness brand artifacts.
%
% \input by the per-artifact .tex files (logo.tex, wordmark.tex, og.tex).
% Anything \providecommand'd here is overridable by the caller — define
% your override BEFORE \input{_common.tex}.

\usepackage{tikz}
\usepackage{xcolor}

% --- Font dispatch ----------------------------------------------------------
%
% \bnFont selects the sans-serif family. Set it in your artifact file
% before \input{_common.tex}. Available values:
%
%   plex       IBM Plex Sans         — modern technical (default)
%   tgheros    TeX Gyre Heros        — Helvetica-clone, neutral
%   source     Source Sans Pro       — humanist, friendly
%   roboto     Roboto                — Google flat-ui geometric
%   fira       Fira Sans             — Mozilla, slightly quirky
%   inter      Inter                 — UI workhorse, very tight
%   montserrat Montserrat            — wide geometric, posterish
%   lato       Lato                  — warm humanist
%   cmbright   Computer Modern Bright — light, scholarly
%   lm         Latin Modern Sans     — LaTeX default, neutral
%
\providecommand{\bnFont}{lm}

\makeatletter
\def\bn@font@plex      {\RequirePackage{plex-sans}}
\def\bn@font@tgheros   {\RequirePackage{tgheros}}
\def\bn@font@source    {\RequirePackage[default]{sourcesanspro}}
\def\bn@font@roboto    {\RequirePackage[sfdefault]{roboto}}
\def\bn@font@fira      {\RequirePackage[sfdefault]{FiraSans}}
\def\bn@font@inter     {\RequirePackage{inter}}
\def\bn@font@montserrat{\RequirePackage[defaultfam,tabular,lining]{montserrat}}
\def\bn@font@lato      {\RequirePackage[default]{lato}}
\def\bn@font@cmbright  {\RequirePackage{cmbright}}
\def\bn@font@lm        {\RequirePackage{lmodern}}
\@nameuse{bn@font@\bnFont}
\makeatother
\renewcommand{\familydefault}{\sfdefault}

% --- Palette ----------------------------------------------------------------
%
% Five role-named slots. \bnDarkMode=1 swaps to a dark variant.
% Override any individual color by \definecolor'ing it AFTER \input{_common.tex}.

\providecommand{\bnDarkMode}{0}

\ifnum\bnDarkMode=1
  % Dark variant
  \definecolor{bnInk}    {HTML}{F2EFE9} % "ink" reads light on dark
  \definecolor{bnBad}    {HTML}{F08A57} % brighter orange so it pops on dark
  \definecolor{bnPaper}  {HTML}{1B1B1B} % near-black canvas
  \definecolor{bnMargin} {HTML}{3D3833} % subtle margin guide
  \definecolor{bnMuted}  {HTML}{A39A8C} % secondary text
\else
  % Light variant (default)
  \definecolor{bnInk}    {HTML}{1B1B1B}
  \definecolor{bnBad}    {HTML}{D9633F}
  \definecolor{bnPaper}  {HTML}{FAFAF7}
  \definecolor{bnMargin} {HTML}{B5B0A5}
  \definecolor{bnMuted}  {HTML}{6F665A}
\fi

% --- Logo geometry (all in cm; tweak freely) --------------------------------

\providecommand{\bnBarCount}{4}     % number of horizontal bars
\providecommand{\bnBarWidth}{4.0}   % length of an in-margin bar
\providecommand{\bnBarHeight}{0.55} % bar thickness
\providecommand{\bnBarGap}{0.32}    % vertical gap between bars
\providecommand{\bnJutWhich}{3}     % 1..bnBarCount — which bar overflows
\providecommand{\bnJutExtra}{1.15}  % how far it sticks out
\providecommand{\bnRound}{0.07}     % corner rounding
\providecommand{\bnShowMargin}{0}   % 1 = draw the dashed margin guide, 0 = hide

% --- The mark itself --------------------------------------------------------
%
% Draws into the current tikzpicture, anchored so the bottom-left of the
% non-jutting bars is at (0,0). The jutting bar extends to the right.

\newcommand{\bnDrawMark}{%
  \pgfmathtruncatemacro{\bnLast}{\bnBarCount-1}%
  \pgfmathsetmacro{\bnTotalH}{\bnBarCount*\bnBarHeight + \bnLast*\bnBarGap}%
  \ifnum\bnShowMargin=1
    % The right-margin guide: a dashed vertical line at the in-margin width,
    % running the full height of the bar stack with a touch of overshoot.
    \draw[bnMargin, dashed, line width=0.4pt, dash pattern=on 2pt off 2pt]
    (\bnBarWidth, -0.12) -- (\bnBarWidth, \bnTotalH+0.12);
  \fi
  \foreach \i in {0,...,\bnLast}{%
      \pgfmathsetmacro{\bny}{(\bnLast-\i)*(\bnBarHeight+\bnBarGap)}%
      \pgfmathtruncatemacro{\bnNum}{\i+1}%
      \ifnum\bnNum=\bnJutWhich
        \fill[bnBad, rounded corners=\bnRound cm]
        (0,\bny) rectangle (\bnBarWidth+\bnJutExtra, \bny+\bnBarHeight);
      \else
        \fill[bnInk, rounded corners=\bnRound cm]
        (0,\bny) rectangle (\bnBarWidth, \bny+\bnBarHeight);
      \fi
    }%
}
. Available values:
%
%   plex       IBM Plex Sans         — modern technical (default)
%   tgheros    TeX Gyre Heros        — Helvetica-clone, neutral
%   source     Source Sans Pro       — humanist, friendly
%   roboto     Roboto                — Google flat-ui geometric
%   fira       Fira Sans             — Mozilla, slightly quirky
%   inter      Inter                 — UI workhorse, very tight
%   montserrat Montserrat            — wide geometric, posterish
%   lato       Lato                  — warm humanist
%   cmbright   Computer Modern Bright — light, scholarly
%   lm         Latin Modern Sans     — LaTeX default, neutral
%
\providecommand{\bnFont}{lm}

\makeatletter
\def\bn@font@plex      {\RequirePackage{plex-sans}}
\def\bn@font@tgheros   {\RequirePackage{tgheros}}
\def\bn@font@source    {\RequirePackage[default]{sourcesanspro}}
\def\bn@font@roboto    {\RequirePackage[sfdefault]{roboto}}
\def\bn@font@fira      {\RequirePackage[sfdefault]{FiraSans}}
\def\bn@font@inter     {\RequirePackage{inter}}
\def\bn@font@montserrat{\RequirePackage[defaultfam,tabular,lining]{montserrat}}
\def\bn@font@lato      {\RequirePackage[default]{lato}}
\def\bn@font@cmbright  {\RequirePackage{cmbright}}
\def\bn@font@lm        {\RequirePackage{lmodern}}
\@nameuse{bn@font@\bnFont}
\makeatother
\renewcommand{\familydefault}{\sfdefault}

% --- Palette ----------------------------------------------------------------
%
% Five role-named slots. \bnDarkMode=1 swaps to a dark variant.
% Override any individual color by \definecolor'ing it AFTER % Shared definitions for the badness brand artifacts.
%
% \input by the per-artifact .tex files (logo.tex, wordmark.tex, og.tex).
% Anything \providecommand'd here is overridable by the caller — define
% your override BEFORE \input{_common.tex}.

\usepackage{tikz}
\usepackage{xcolor}

% --- Font dispatch ----------------------------------------------------------
%
% \bnFont selects the sans-serif family. Set it in your artifact file
% before \input{_common.tex}. Available values:
%
%   plex       IBM Plex Sans         — modern technical (default)
%   tgheros    TeX Gyre Heros        — Helvetica-clone, neutral
%   source     Source Sans Pro       — humanist, friendly
%   roboto     Roboto                — Google flat-ui geometric
%   fira       Fira Sans             — Mozilla, slightly quirky
%   inter      Inter                 — UI workhorse, very tight
%   montserrat Montserrat            — wide geometric, posterish
%   lato       Lato                  — warm humanist
%   cmbright   Computer Modern Bright — light, scholarly
%   lm         Latin Modern Sans     — LaTeX default, neutral
%
\providecommand{\bnFont}{lm}

\makeatletter
\def\bn@font@plex      {\RequirePackage{plex-sans}}
\def\bn@font@tgheros   {\RequirePackage{tgheros}}
\def\bn@font@source    {\RequirePackage[default]{sourcesanspro}}
\def\bn@font@roboto    {\RequirePackage[sfdefault]{roboto}}
\def\bn@font@fira      {\RequirePackage[sfdefault]{FiraSans}}
\def\bn@font@inter     {\RequirePackage{inter}}
\def\bn@font@montserrat{\RequirePackage[defaultfam,tabular,lining]{montserrat}}
\def\bn@font@lato      {\RequirePackage[default]{lato}}
\def\bn@font@cmbright  {\RequirePackage{cmbright}}
\def\bn@font@lm        {\RequirePackage{lmodern}}
\@nameuse{bn@font@\bnFont}
\makeatother
\renewcommand{\familydefault}{\sfdefault}

% --- Palette ----------------------------------------------------------------
%
% Five role-named slots. \bnDarkMode=1 swaps to a dark variant.
% Override any individual color by \definecolor'ing it AFTER \input{_common.tex}.

\providecommand{\bnDarkMode}{0}

\ifnum\bnDarkMode=1
  % Dark variant
  \definecolor{bnInk}    {HTML}{F2EFE9} % "ink" reads light on dark
  \definecolor{bnBad}    {HTML}{F08A57} % brighter orange so it pops on dark
  \definecolor{bnPaper}  {HTML}{1B1B1B} % near-black canvas
  \definecolor{bnMargin} {HTML}{3D3833} % subtle margin guide
  \definecolor{bnMuted}  {HTML}{A39A8C} % secondary text
\else
  % Light variant (default)
  \definecolor{bnInk}    {HTML}{1B1B1B}
  \definecolor{bnBad}    {HTML}{D9633F}
  \definecolor{bnPaper}  {HTML}{FAFAF7}
  \definecolor{bnMargin} {HTML}{B5B0A5}
  \definecolor{bnMuted}  {HTML}{6F665A}
\fi

% --- Logo geometry (all in cm; tweak freely) --------------------------------

\providecommand{\bnBarCount}{4}     % number of horizontal bars
\providecommand{\bnBarWidth}{4.0}   % length of an in-margin bar
\providecommand{\bnBarHeight}{0.55} % bar thickness
\providecommand{\bnBarGap}{0.32}    % vertical gap between bars
\providecommand{\bnJutWhich}{3}     % 1..bnBarCount — which bar overflows
\providecommand{\bnJutExtra}{1.15}  % how far it sticks out
\providecommand{\bnRound}{0.07}     % corner rounding
\providecommand{\bnShowMargin}{0}   % 1 = draw the dashed margin guide, 0 = hide

% --- The mark itself --------------------------------------------------------
%
% Draws into the current tikzpicture, anchored so the bottom-left of the
% non-jutting bars is at (0,0). The jutting bar extends to the right.

\newcommand{\bnDrawMark}{%
  \pgfmathtruncatemacro{\bnLast}{\bnBarCount-1}%
  \pgfmathsetmacro{\bnTotalH}{\bnBarCount*\bnBarHeight + \bnLast*\bnBarGap}%
  \ifnum\bnShowMargin=1
    % The right-margin guide: a dashed vertical line at the in-margin width,
    % running the full height of the bar stack with a touch of overshoot.
    \draw[bnMargin, dashed, line width=0.4pt, dash pattern=on 2pt off 2pt]
    (\bnBarWidth, -0.12) -- (\bnBarWidth, \bnTotalH+0.12);
  \fi
  \foreach \i in {0,...,\bnLast}{%
      \pgfmathsetmacro{\bny}{(\bnLast-\i)*(\bnBarHeight+\bnBarGap)}%
      \pgfmathtruncatemacro{\bnNum}{\i+1}%
      \ifnum\bnNum=\bnJutWhich
        \fill[bnBad, rounded corners=\bnRound cm]
        (0,\bny) rectangle (\bnBarWidth+\bnJutExtra, \bny+\bnBarHeight);
      \else
        \fill[bnInk, rounded corners=\bnRound cm]
        (0,\bny) rectangle (\bnBarWidth, \bny+\bnBarHeight);
      \fi
    }%
}
.

\providecommand{\bnDarkMode}{0}

\ifnum\bnDarkMode=1
  % Dark variant
  \definecolor{bnInk}    {HTML}{F2EFE9} % "ink" reads light on dark
  \definecolor{bnBad}    {HTML}{F08A57} % brighter orange so it pops on dark
  \definecolor{bnPaper}  {HTML}{1B1B1B} % near-black canvas
  \definecolor{bnMargin} {HTML}{3D3833} % subtle margin guide
  \definecolor{bnMuted}  {HTML}{A39A8C} % secondary text
\else
  % Light variant (default)
  \definecolor{bnInk}    {HTML}{1B1B1B}
  \definecolor{bnBad}    {HTML}{D9633F}
  \definecolor{bnPaper}  {HTML}{FAFAF7}
  \definecolor{bnMargin} {HTML}{B5B0A5}
  \definecolor{bnMuted}  {HTML}{6F665A}
\fi

% --- Logo geometry (all in cm; tweak freely) --------------------------------

\providecommand{\bnBarCount}{4}     % number of horizontal bars
\providecommand{\bnBarWidth}{4.0}   % length of an in-margin bar
\providecommand{\bnBarHeight}{0.55} % bar thickness
\providecommand{\bnBarGap}{0.32}    % vertical gap between bars
\providecommand{\bnJutWhich}{3}     % 1..bnBarCount — which bar overflows
\providecommand{\bnJutExtra}{1.15}  % how far it sticks out
\providecommand{\bnRound}{0.07}     % corner rounding
\providecommand{\bnShowMargin}{0}   % 1 = draw the dashed margin guide, 0 = hide

% --- The mark itself --------------------------------------------------------
%
% Draws into the current tikzpicture, anchored so the bottom-left of the
% non-jutting bars is at (0,0). The jutting bar extends to the right.

\newcommand{\bnDrawMark}{%
  \pgfmathtruncatemacro{\bnLast}{\bnBarCount-1}%
  \pgfmathsetmacro{\bnTotalH}{\bnBarCount*\bnBarHeight + \bnLast*\bnBarGap}%
  \ifnum\bnShowMargin=1
    % The right-margin guide: a dashed vertical line at the in-margin width,
    % running the full height of the bar stack with a touch of overshoot.
    \draw[bnMargin, dashed, line width=0.4pt, dash pattern=on 2pt off 2pt]
    (\bnBarWidth, -0.12) -- (\bnBarWidth, \bnTotalH+0.12);
  \fi
  \foreach \i in {0,...,\bnLast}{%
      \pgfmathsetmacro{\bny}{(\bnLast-\i)*(\bnBarHeight+\bnBarGap)}%
      \pgfmathtruncatemacro{\bnNum}{\i+1}%
      \ifnum\bnNum=\bnJutWhich
        \fill[bnBad, rounded corners=\bnRound cm]
        (0,\bny) rectangle (\bnBarWidth+\bnJutExtra, \bny+\bnBarHeight);
      \else
        \fill[bnInk, rounded corners=\bnRound cm]
        (0,\bny) rectangle (\bnBarWidth, \bny+\bnBarHeight);
      \fi
    }%
}
. Available values:
%
%   plex       IBM Plex Sans         — modern technical (default)
%   tgheros    TeX Gyre Heros        — Helvetica-clone, neutral
%   source     Source Sans Pro       — humanist, friendly
%   roboto     Roboto                — Google flat-ui geometric
%   fira       Fira Sans             — Mozilla, slightly quirky
%   inter      Inter                 — UI workhorse, very tight
%   montserrat Montserrat            — wide geometric, posterish
%   lato       Lato                  — warm humanist
%   cmbright   Computer Modern Bright — light, scholarly
%   lm         Latin Modern Sans     — LaTeX default, neutral
%
\providecommand{\bnFont}{lm}

\makeatletter
\def\bn@font@plex      {\RequirePackage{plex-sans}}
\def\bn@font@tgheros   {\RequirePackage{tgheros}}
\def\bn@font@source    {\RequirePackage[default]{sourcesanspro}}
\def\bn@font@roboto    {\RequirePackage[sfdefault]{roboto}}
\def\bn@font@fira      {\RequirePackage[sfdefault]{FiraSans}}
\def\bn@font@inter     {\RequirePackage{inter}}
\def\bn@font@montserrat{\RequirePackage[defaultfam,tabular,lining]{montserrat}}
\def\bn@font@lato      {\RequirePackage[default]{lato}}
\def\bn@font@cmbright  {\RequirePackage{cmbright}}
\def\bn@font@lm        {\RequirePackage{lmodern}}
\@nameuse{bn@font@\bnFont}
\makeatother
\renewcommand{\familydefault}{\sfdefault}

% --- Palette ----------------------------------------------------------------
%
% Five role-named slots. \bnDarkMode=1 swaps to a dark variant.
% Override any individual color by \definecolor'ing it AFTER % Shared definitions for the badness brand artifacts.
%
% \input by the per-artifact .tex files (logo.tex, wordmark.tex, og.tex).
% Anything \providecommand'd here is overridable by the caller — define
% your override BEFORE % Shared definitions for the badness brand artifacts.
%
% \input by the per-artifact .tex files (logo.tex, wordmark.tex, og.tex).
% Anything \providecommand'd here is overridable by the caller — define
% your override BEFORE \input{_common.tex}.

\usepackage{tikz}
\usepackage{xcolor}

% --- Font dispatch ----------------------------------------------------------
%
% \bnFont selects the sans-serif family. Set it in your artifact file
% before \input{_common.tex}. Available values:
%
%   plex       IBM Plex Sans         — modern technical (default)
%   tgheros    TeX Gyre Heros        — Helvetica-clone, neutral
%   source     Source Sans Pro       — humanist, friendly
%   roboto     Roboto                — Google flat-ui geometric
%   fira       Fira Sans             — Mozilla, slightly quirky
%   inter      Inter                 — UI workhorse, very tight
%   montserrat Montserrat            — wide geometric, posterish
%   lato       Lato                  — warm humanist
%   cmbright   Computer Modern Bright — light, scholarly
%   lm         Latin Modern Sans     — LaTeX default, neutral
%
\providecommand{\bnFont}{lm}

\makeatletter
\def\bn@font@plex      {\RequirePackage{plex-sans}}
\def\bn@font@tgheros   {\RequirePackage{tgheros}}
\def\bn@font@source    {\RequirePackage[default]{sourcesanspro}}
\def\bn@font@roboto    {\RequirePackage[sfdefault]{roboto}}
\def\bn@font@fira      {\RequirePackage[sfdefault]{FiraSans}}
\def\bn@font@inter     {\RequirePackage{inter}}
\def\bn@font@montserrat{\RequirePackage[defaultfam,tabular,lining]{montserrat}}
\def\bn@font@lato      {\RequirePackage[default]{lato}}
\def\bn@font@cmbright  {\RequirePackage{cmbright}}
\def\bn@font@lm        {\RequirePackage{lmodern}}
\@nameuse{bn@font@\bnFont}
\makeatother
\renewcommand{\familydefault}{\sfdefault}

% --- Palette ----------------------------------------------------------------
%
% Five role-named slots. \bnDarkMode=1 swaps to a dark variant.
% Override any individual color by \definecolor'ing it AFTER \input{_common.tex}.

\providecommand{\bnDarkMode}{0}

\ifnum\bnDarkMode=1
  % Dark variant
  \definecolor{bnInk}    {HTML}{F2EFE9} % "ink" reads light on dark
  \definecolor{bnBad}    {HTML}{F08A57} % brighter orange so it pops on dark
  \definecolor{bnPaper}  {HTML}{1B1B1B} % near-black canvas
  \definecolor{bnMargin} {HTML}{3D3833} % subtle margin guide
  \definecolor{bnMuted}  {HTML}{A39A8C} % secondary text
\else
  % Light variant (default)
  \definecolor{bnInk}    {HTML}{1B1B1B}
  \definecolor{bnBad}    {HTML}{D9633F}
  \definecolor{bnPaper}  {HTML}{FAFAF7}
  \definecolor{bnMargin} {HTML}{B5B0A5}
  \definecolor{bnMuted}  {HTML}{6F665A}
\fi

% --- Logo geometry (all in cm; tweak freely) --------------------------------

\providecommand{\bnBarCount}{4}     % number of horizontal bars
\providecommand{\bnBarWidth}{4.0}   % length of an in-margin bar
\providecommand{\bnBarHeight}{0.55} % bar thickness
\providecommand{\bnBarGap}{0.32}    % vertical gap between bars
\providecommand{\bnJutWhich}{3}     % 1..bnBarCount — which bar overflows
\providecommand{\bnJutExtra}{1.15}  % how far it sticks out
\providecommand{\bnRound}{0.07}     % corner rounding
\providecommand{\bnShowMargin}{0}   % 1 = draw the dashed margin guide, 0 = hide

% --- The mark itself --------------------------------------------------------
%
% Draws into the current tikzpicture, anchored so the bottom-left of the
% non-jutting bars is at (0,0). The jutting bar extends to the right.

\newcommand{\bnDrawMark}{%
  \pgfmathtruncatemacro{\bnLast}{\bnBarCount-1}%
  \pgfmathsetmacro{\bnTotalH}{\bnBarCount*\bnBarHeight + \bnLast*\bnBarGap}%
  \ifnum\bnShowMargin=1
    % The right-margin guide: a dashed vertical line at the in-margin width,
    % running the full height of the bar stack with a touch of overshoot.
    \draw[bnMargin, dashed, line width=0.4pt, dash pattern=on 2pt off 2pt]
    (\bnBarWidth, -0.12) -- (\bnBarWidth, \bnTotalH+0.12);
  \fi
  \foreach \i in {0,...,\bnLast}{%
      \pgfmathsetmacro{\bny}{(\bnLast-\i)*(\bnBarHeight+\bnBarGap)}%
      \pgfmathtruncatemacro{\bnNum}{\i+1}%
      \ifnum\bnNum=\bnJutWhich
        \fill[bnBad, rounded corners=\bnRound cm]
        (0,\bny) rectangle (\bnBarWidth+\bnJutExtra, \bny+\bnBarHeight);
      \else
        \fill[bnInk, rounded corners=\bnRound cm]
        (0,\bny) rectangle (\bnBarWidth, \bny+\bnBarHeight);
      \fi
    }%
}
.

\usepackage{tikz}
\usepackage{xcolor}

% --- Font dispatch ----------------------------------------------------------
%
% \bnFont selects the sans-serif family. Set it in your artifact file
% before % Shared definitions for the badness brand artifacts.
%
% \input by the per-artifact .tex files (logo.tex, wordmark.tex, og.tex).
% Anything \providecommand'd here is overridable by the caller — define
% your override BEFORE \input{_common.tex}.

\usepackage{tikz}
\usepackage{xcolor}

% --- Font dispatch ----------------------------------------------------------
%
% \bnFont selects the sans-serif family. Set it in your artifact file
% before \input{_common.tex}. Available values:
%
%   plex       IBM Plex Sans         — modern technical (default)
%   tgheros    TeX Gyre Heros        — Helvetica-clone, neutral
%   source     Source Sans Pro       — humanist, friendly
%   roboto     Roboto                — Google flat-ui geometric
%   fira       Fira Sans             — Mozilla, slightly quirky
%   inter      Inter                 — UI workhorse, very tight
%   montserrat Montserrat            — wide geometric, posterish
%   lato       Lato                  — warm humanist
%   cmbright   Computer Modern Bright — light, scholarly
%   lm         Latin Modern Sans     — LaTeX default, neutral
%
\providecommand{\bnFont}{lm}

\makeatletter
\def\bn@font@plex      {\RequirePackage{plex-sans}}
\def\bn@font@tgheros   {\RequirePackage{tgheros}}
\def\bn@font@source    {\RequirePackage[default]{sourcesanspro}}
\def\bn@font@roboto    {\RequirePackage[sfdefault]{roboto}}
\def\bn@font@fira      {\RequirePackage[sfdefault]{FiraSans}}
\def\bn@font@inter     {\RequirePackage{inter}}
\def\bn@font@montserrat{\RequirePackage[defaultfam,tabular,lining]{montserrat}}
\def\bn@font@lato      {\RequirePackage[default]{lato}}
\def\bn@font@cmbright  {\RequirePackage{cmbright}}
\def\bn@font@lm        {\RequirePackage{lmodern}}
\@nameuse{bn@font@\bnFont}
\makeatother
\renewcommand{\familydefault}{\sfdefault}

% --- Palette ----------------------------------------------------------------
%
% Five role-named slots. \bnDarkMode=1 swaps to a dark variant.
% Override any individual color by \definecolor'ing it AFTER \input{_common.tex}.

\providecommand{\bnDarkMode}{0}

\ifnum\bnDarkMode=1
  % Dark variant
  \definecolor{bnInk}    {HTML}{F2EFE9} % "ink" reads light on dark
  \definecolor{bnBad}    {HTML}{F08A57} % brighter orange so it pops on dark
  \definecolor{bnPaper}  {HTML}{1B1B1B} % near-black canvas
  \definecolor{bnMargin} {HTML}{3D3833} % subtle margin guide
  \definecolor{bnMuted}  {HTML}{A39A8C} % secondary text
\else
  % Light variant (default)
  \definecolor{bnInk}    {HTML}{1B1B1B}
  \definecolor{bnBad}    {HTML}{D9633F}
  \definecolor{bnPaper}  {HTML}{FAFAF7}
  \definecolor{bnMargin} {HTML}{B5B0A5}
  \definecolor{bnMuted}  {HTML}{6F665A}
\fi

% --- Logo geometry (all in cm; tweak freely) --------------------------------

\providecommand{\bnBarCount}{4}     % number of horizontal bars
\providecommand{\bnBarWidth}{4.0}   % length of an in-margin bar
\providecommand{\bnBarHeight}{0.55} % bar thickness
\providecommand{\bnBarGap}{0.32}    % vertical gap between bars
\providecommand{\bnJutWhich}{3}     % 1..bnBarCount — which bar overflows
\providecommand{\bnJutExtra}{1.15}  % how far it sticks out
\providecommand{\bnRound}{0.07}     % corner rounding
\providecommand{\bnShowMargin}{0}   % 1 = draw the dashed margin guide, 0 = hide

% --- The mark itself --------------------------------------------------------
%
% Draws into the current tikzpicture, anchored so the bottom-left of the
% non-jutting bars is at (0,0). The jutting bar extends to the right.

\newcommand{\bnDrawMark}{%
  \pgfmathtruncatemacro{\bnLast}{\bnBarCount-1}%
  \pgfmathsetmacro{\bnTotalH}{\bnBarCount*\bnBarHeight + \bnLast*\bnBarGap}%
  \ifnum\bnShowMargin=1
    % The right-margin guide: a dashed vertical line at the in-margin width,
    % running the full height of the bar stack with a touch of overshoot.
    \draw[bnMargin, dashed, line width=0.4pt, dash pattern=on 2pt off 2pt]
    (\bnBarWidth, -0.12) -- (\bnBarWidth, \bnTotalH+0.12);
  \fi
  \foreach \i in {0,...,\bnLast}{%
      \pgfmathsetmacro{\bny}{(\bnLast-\i)*(\bnBarHeight+\bnBarGap)}%
      \pgfmathtruncatemacro{\bnNum}{\i+1}%
      \ifnum\bnNum=\bnJutWhich
        \fill[bnBad, rounded corners=\bnRound cm]
        (0,\bny) rectangle (\bnBarWidth+\bnJutExtra, \bny+\bnBarHeight);
      \else
        \fill[bnInk, rounded corners=\bnRound cm]
        (0,\bny) rectangle (\bnBarWidth, \bny+\bnBarHeight);
      \fi
    }%
}
. Available values:
%
%   plex       IBM Plex Sans         — modern technical (default)
%   tgheros    TeX Gyre Heros        — Helvetica-clone, neutral
%   source     Source Sans Pro       — humanist, friendly
%   roboto     Roboto                — Google flat-ui geometric
%   fira       Fira Sans             — Mozilla, slightly quirky
%   inter      Inter                 — UI workhorse, very tight
%   montserrat Montserrat            — wide geometric, posterish
%   lato       Lato                  — warm humanist
%   cmbright   Computer Modern Bright — light, scholarly
%   lm         Latin Modern Sans     — LaTeX default, neutral
%
\providecommand{\bnFont}{lm}

\makeatletter
\def\bn@font@plex      {\RequirePackage{plex-sans}}
\def\bn@font@tgheros   {\RequirePackage{tgheros}}
\def\bn@font@source    {\RequirePackage[default]{sourcesanspro}}
\def\bn@font@roboto    {\RequirePackage[sfdefault]{roboto}}
\def\bn@font@fira      {\RequirePackage[sfdefault]{FiraSans}}
\def\bn@font@inter     {\RequirePackage{inter}}
\def\bn@font@montserrat{\RequirePackage[defaultfam,tabular,lining]{montserrat}}
\def\bn@font@lato      {\RequirePackage[default]{lato}}
\def\bn@font@cmbright  {\RequirePackage{cmbright}}
\def\bn@font@lm        {\RequirePackage{lmodern}}
\@nameuse{bn@font@\bnFont}
\makeatother
\renewcommand{\familydefault}{\sfdefault}

% --- Palette ----------------------------------------------------------------
%
% Five role-named slots. \bnDarkMode=1 swaps to a dark variant.
% Override any individual color by \definecolor'ing it AFTER % Shared definitions for the badness brand artifacts.
%
% \input by the per-artifact .tex files (logo.tex, wordmark.tex, og.tex).
% Anything \providecommand'd here is overridable by the caller — define
% your override BEFORE \input{_common.tex}.

\usepackage{tikz}
\usepackage{xcolor}

% --- Font dispatch ----------------------------------------------------------
%
% \bnFont selects the sans-serif family. Set it in your artifact file
% before \input{_common.tex}. Available values:
%
%   plex       IBM Plex Sans         — modern technical (default)
%   tgheros    TeX Gyre Heros        — Helvetica-clone, neutral
%   source     Source Sans Pro       — humanist, friendly
%   roboto     Roboto                — Google flat-ui geometric
%   fira       Fira Sans             — Mozilla, slightly quirky
%   inter      Inter                 — UI workhorse, very tight
%   montserrat Montserrat            — wide geometric, posterish
%   lato       Lato                  — warm humanist
%   cmbright   Computer Modern Bright — light, scholarly
%   lm         Latin Modern Sans     — LaTeX default, neutral
%
\providecommand{\bnFont}{lm}

\makeatletter
\def\bn@font@plex      {\RequirePackage{plex-sans}}
\def\bn@font@tgheros   {\RequirePackage{tgheros}}
\def\bn@font@source    {\RequirePackage[default]{sourcesanspro}}
\def\bn@font@roboto    {\RequirePackage[sfdefault]{roboto}}
\def\bn@font@fira      {\RequirePackage[sfdefault]{FiraSans}}
\def\bn@font@inter     {\RequirePackage{inter}}
\def\bn@font@montserrat{\RequirePackage[defaultfam,tabular,lining]{montserrat}}
\def\bn@font@lato      {\RequirePackage[default]{lato}}
\def\bn@font@cmbright  {\RequirePackage{cmbright}}
\def\bn@font@lm        {\RequirePackage{lmodern}}
\@nameuse{bn@font@\bnFont}
\makeatother
\renewcommand{\familydefault}{\sfdefault}

% --- Palette ----------------------------------------------------------------
%
% Five role-named slots. \bnDarkMode=1 swaps to a dark variant.
% Override any individual color by \definecolor'ing it AFTER \input{_common.tex}.

\providecommand{\bnDarkMode}{0}

\ifnum\bnDarkMode=1
  % Dark variant
  \definecolor{bnInk}    {HTML}{F2EFE9} % "ink" reads light on dark
  \definecolor{bnBad}    {HTML}{F08A57} % brighter orange so it pops on dark
  \definecolor{bnPaper}  {HTML}{1B1B1B} % near-black canvas
  \definecolor{bnMargin} {HTML}{3D3833} % subtle margin guide
  \definecolor{bnMuted}  {HTML}{A39A8C} % secondary text
\else
  % Light variant (default)
  \definecolor{bnInk}    {HTML}{1B1B1B}
  \definecolor{bnBad}    {HTML}{D9633F}
  \definecolor{bnPaper}  {HTML}{FAFAF7}
  \definecolor{bnMargin} {HTML}{B5B0A5}
  \definecolor{bnMuted}  {HTML}{6F665A}
\fi

% --- Logo geometry (all in cm; tweak freely) --------------------------------

\providecommand{\bnBarCount}{4}     % number of horizontal bars
\providecommand{\bnBarWidth}{4.0}   % length of an in-margin bar
\providecommand{\bnBarHeight}{0.55} % bar thickness
\providecommand{\bnBarGap}{0.32}    % vertical gap between bars
\providecommand{\bnJutWhich}{3}     % 1..bnBarCount — which bar overflows
\providecommand{\bnJutExtra}{1.15}  % how far it sticks out
\providecommand{\bnRound}{0.07}     % corner rounding
\providecommand{\bnShowMargin}{0}   % 1 = draw the dashed margin guide, 0 = hide

% --- The mark itself --------------------------------------------------------
%
% Draws into the current tikzpicture, anchored so the bottom-left of the
% non-jutting bars is at (0,0). The jutting bar extends to the right.

\newcommand{\bnDrawMark}{%
  \pgfmathtruncatemacro{\bnLast}{\bnBarCount-1}%
  \pgfmathsetmacro{\bnTotalH}{\bnBarCount*\bnBarHeight + \bnLast*\bnBarGap}%
  \ifnum\bnShowMargin=1
    % The right-margin guide: a dashed vertical line at the in-margin width,
    % running the full height of the bar stack with a touch of overshoot.
    \draw[bnMargin, dashed, line width=0.4pt, dash pattern=on 2pt off 2pt]
    (\bnBarWidth, -0.12) -- (\bnBarWidth, \bnTotalH+0.12);
  \fi
  \foreach \i in {0,...,\bnLast}{%
      \pgfmathsetmacro{\bny}{(\bnLast-\i)*(\bnBarHeight+\bnBarGap)}%
      \pgfmathtruncatemacro{\bnNum}{\i+1}%
      \ifnum\bnNum=\bnJutWhich
        \fill[bnBad, rounded corners=\bnRound cm]
        (0,\bny) rectangle (\bnBarWidth+\bnJutExtra, \bny+\bnBarHeight);
      \else
        \fill[bnInk, rounded corners=\bnRound cm]
        (0,\bny) rectangle (\bnBarWidth, \bny+\bnBarHeight);
      \fi
    }%
}
.

\providecommand{\bnDarkMode}{0}

\ifnum\bnDarkMode=1
  % Dark variant
  \definecolor{bnInk}    {HTML}{F2EFE9} % "ink" reads light on dark
  \definecolor{bnBad}    {HTML}{F08A57} % brighter orange so it pops on dark
  \definecolor{bnPaper}  {HTML}{1B1B1B} % near-black canvas
  \definecolor{bnMargin} {HTML}{3D3833} % subtle margin guide
  \definecolor{bnMuted}  {HTML}{A39A8C} % secondary text
\else
  % Light variant (default)
  \definecolor{bnInk}    {HTML}{1B1B1B}
  \definecolor{bnBad}    {HTML}{D9633F}
  \definecolor{bnPaper}  {HTML}{FAFAF7}
  \definecolor{bnMargin} {HTML}{B5B0A5}
  \definecolor{bnMuted}  {HTML}{6F665A}
\fi

% --- Logo geometry (all in cm; tweak freely) --------------------------------

\providecommand{\bnBarCount}{4}     % number of horizontal bars
\providecommand{\bnBarWidth}{4.0}   % length of an in-margin bar
\providecommand{\bnBarHeight}{0.55} % bar thickness
\providecommand{\bnBarGap}{0.32}    % vertical gap between bars
\providecommand{\bnJutWhich}{3}     % 1..bnBarCount — which bar overflows
\providecommand{\bnJutExtra}{1.15}  % how far it sticks out
\providecommand{\bnRound}{0.07}     % corner rounding
\providecommand{\bnShowMargin}{0}   % 1 = draw the dashed margin guide, 0 = hide

% --- The mark itself --------------------------------------------------------
%
% Draws into the current tikzpicture, anchored so the bottom-left of the
% non-jutting bars is at (0,0). The jutting bar extends to the right.

\newcommand{\bnDrawMark}{%
  \pgfmathtruncatemacro{\bnLast}{\bnBarCount-1}%
  \pgfmathsetmacro{\bnTotalH}{\bnBarCount*\bnBarHeight + \bnLast*\bnBarGap}%
  \ifnum\bnShowMargin=1
    % The right-margin guide: a dashed vertical line at the in-margin width,
    % running the full height of the bar stack with a touch of overshoot.
    \draw[bnMargin, dashed, line width=0.4pt, dash pattern=on 2pt off 2pt]
    (\bnBarWidth, -0.12) -- (\bnBarWidth, \bnTotalH+0.12);
  \fi
  \foreach \i in {0,...,\bnLast}{%
      \pgfmathsetmacro{\bny}{(\bnLast-\i)*(\bnBarHeight+\bnBarGap)}%
      \pgfmathtruncatemacro{\bnNum}{\i+1}%
      \ifnum\bnNum=\bnJutWhich
        \fill[bnBad, rounded corners=\bnRound cm]
        (0,\bny) rectangle (\bnBarWidth+\bnJutExtra, \bny+\bnBarHeight);
      \else
        \fill[bnInk, rounded corners=\bnRound cm]
        (0,\bny) rectangle (\bnBarWidth, \bny+\bnBarHeight);
      \fi
    }%
}
.

\providecommand{\bnDarkMode}{0}

\ifnum\bnDarkMode=1
  % Dark variant
  \definecolor{bnInk}    {HTML}{F2EFE9} % "ink" reads light on dark
  \definecolor{bnBad}    {HTML}{F08A57} % brighter orange so it pops on dark
  \definecolor{bnPaper}  {HTML}{1B1B1B} % near-black canvas
  \definecolor{bnMargin} {HTML}{3D3833} % subtle margin guide
  \definecolor{bnMuted}  {HTML}{A39A8C} % secondary text
\else
  % Light variant (default)
  \definecolor{bnInk}    {HTML}{1B1B1B}
  \definecolor{bnBad}    {HTML}{D9633F}
  \definecolor{bnPaper}  {HTML}{FAFAF7}
  \definecolor{bnMargin} {HTML}{B5B0A5}
  \definecolor{bnMuted}  {HTML}{6F665A}
\fi

% --- Logo geometry (all in cm; tweak freely) --------------------------------

\providecommand{\bnBarCount}{4}     % number of horizontal bars
\providecommand{\bnBarWidth}{4.0}   % length of an in-margin bar
\providecommand{\bnBarHeight}{0.55} % bar thickness
\providecommand{\bnBarGap}{0.32}    % vertical gap between bars
\providecommand{\bnJutWhich}{3}     % 1..bnBarCount — which bar overflows
\providecommand{\bnJutExtra}{1.15}  % how far it sticks out
\providecommand{\bnRound}{0.07}     % corner rounding
\providecommand{\bnShowMargin}{0}   % 1 = draw the dashed margin guide, 0 = hide

% --- The mark itself --------------------------------------------------------
%
% Draws into the current tikzpicture, anchored so the bottom-left of the
% non-jutting bars is at (0,0). The jutting bar extends to the right.

\newcommand{\bnDrawMark}{%
  \pgfmathtruncatemacro{\bnLast}{\bnBarCount-1}%
  \pgfmathsetmacro{\bnTotalH}{\bnBarCount*\bnBarHeight + \bnLast*\bnBarGap}%
  \ifnum\bnShowMargin=1
    % The right-margin guide: a dashed vertical line at the in-margin width,
    % running the full height of the bar stack with a touch of overshoot.
    \draw[bnMargin, dashed, line width=0.4pt, dash pattern=on 2pt off 2pt]
    (\bnBarWidth, -0.12) -- (\bnBarWidth, \bnTotalH+0.12);
  \fi
  \foreach \i in {0,...,\bnLast}{%
      \pgfmathsetmacro{\bny}{(\bnLast-\i)*(\bnBarHeight+\bnBarGap)}%
      \pgfmathtruncatemacro{\bnNum}{\i+1}%
      \ifnum\bnNum=\bnJutWhich
        \fill[bnBad, rounded corners=\bnRound cm]
        (0,\bny) rectangle (\bnBarWidth+\bnJutExtra, \bny+\bnBarHeight);
      \else
        \fill[bnInk, rounded corners=\bnRound cm]
        (0,\bny) rectangle (\bnBarWidth, \bny+\bnBarHeight);
      \fi
    }%
}
. Available values:
%
%   plex       IBM Plex Sans         — modern technical (default)
%   tgheros    TeX Gyre Heros        — Helvetica-clone, neutral
%   source     Source Sans Pro       — humanist, friendly
%   roboto     Roboto                — Google flat-ui geometric
%   fira       Fira Sans             — Mozilla, slightly quirky
%   inter      Inter                 — UI workhorse, very tight
%   montserrat Montserrat            — wide geometric, posterish
%   lato       Lato                  — warm humanist
%   cmbright   Computer Modern Bright — light, scholarly
%   lm         Latin Modern Sans     — LaTeX default, neutral
%
\providecommand{\bnFont}{lm}

\makeatletter
\def\bn@font@plex      {\RequirePackage{plex-sans}}
\def\bn@font@tgheros   {\RequirePackage{tgheros}}
\def\bn@font@source    {\RequirePackage[default]{sourcesanspro}}
\def\bn@font@roboto    {\RequirePackage[sfdefault]{roboto}}
\def\bn@font@fira      {\RequirePackage[sfdefault]{FiraSans}}
\def\bn@font@inter     {\RequirePackage{inter}}
\def\bn@font@montserrat{\RequirePackage[defaultfam,tabular,lining]{montserrat}}
\def\bn@font@lato      {\RequirePackage[default]{lato}}
\def\bn@font@cmbright  {\RequirePackage{cmbright}}
\def\bn@font@lm        {\RequirePackage{lmodern}}
\@nameuse{bn@font@\bnFont}
\makeatother
\renewcommand{\familydefault}{\sfdefault}

% --- Palette ----------------------------------------------------------------
%
% Five role-named slots. \bnDarkMode=1 swaps to a dark variant.
% Override any individual color by \definecolor'ing it AFTER % Shared definitions for the badness brand artifacts.
%
% \input by the per-artifact .tex files (logo.tex, wordmark.tex, og.tex).
% Anything \providecommand'd here is overridable by the caller — define
% your override BEFORE % Shared definitions for the badness brand artifacts.
%
% \input by the per-artifact .tex files (logo.tex, wordmark.tex, og.tex).
% Anything \providecommand'd here is overridable by the caller — define
% your override BEFORE % Shared definitions for the badness brand artifacts.
%
% \input by the per-artifact .tex files (logo.tex, wordmark.tex, og.tex).
% Anything \providecommand'd here is overridable by the caller — define
% your override BEFORE \input{_common.tex}.

\usepackage{tikz}
\usepackage{xcolor}

% --- Font dispatch ----------------------------------------------------------
%
% \bnFont selects the sans-serif family. Set it in your artifact file
% before \input{_common.tex}. Available values:
%
%   plex       IBM Plex Sans         — modern technical (default)
%   tgheros    TeX Gyre Heros        — Helvetica-clone, neutral
%   source     Source Sans Pro       — humanist, friendly
%   roboto     Roboto                — Google flat-ui geometric
%   fira       Fira Sans             — Mozilla, slightly quirky
%   inter      Inter                 — UI workhorse, very tight
%   montserrat Montserrat            — wide geometric, posterish
%   lato       Lato                  — warm humanist
%   cmbright   Computer Modern Bright — light, scholarly
%   lm         Latin Modern Sans     — LaTeX default, neutral
%
\providecommand{\bnFont}{lm}

\makeatletter
\def\bn@font@plex      {\RequirePackage{plex-sans}}
\def\bn@font@tgheros   {\RequirePackage{tgheros}}
\def\bn@font@source    {\RequirePackage[default]{sourcesanspro}}
\def\bn@font@roboto    {\RequirePackage[sfdefault]{roboto}}
\def\bn@font@fira      {\RequirePackage[sfdefault]{FiraSans}}
\def\bn@font@inter     {\RequirePackage{inter}}
\def\bn@font@montserrat{\RequirePackage[defaultfam,tabular,lining]{montserrat}}
\def\bn@font@lato      {\RequirePackage[default]{lato}}
\def\bn@font@cmbright  {\RequirePackage{cmbright}}
\def\bn@font@lm        {\RequirePackage{lmodern}}
\@nameuse{bn@font@\bnFont}
\makeatother
\renewcommand{\familydefault}{\sfdefault}

% --- Palette ----------------------------------------------------------------
%
% Five role-named slots. \bnDarkMode=1 swaps to a dark variant.
% Override any individual color by \definecolor'ing it AFTER \input{_common.tex}.

\providecommand{\bnDarkMode}{0}

\ifnum\bnDarkMode=1
  % Dark variant
  \definecolor{bnInk}    {HTML}{F2EFE9} % "ink" reads light on dark
  \definecolor{bnBad}    {HTML}{F08A57} % brighter orange so it pops on dark
  \definecolor{bnPaper}  {HTML}{1B1B1B} % near-black canvas
  \definecolor{bnMargin} {HTML}{3D3833} % subtle margin guide
  \definecolor{bnMuted}  {HTML}{A39A8C} % secondary text
\else
  % Light variant (default)
  \definecolor{bnInk}    {HTML}{1B1B1B}
  \definecolor{bnBad}    {HTML}{D9633F}
  \definecolor{bnPaper}  {HTML}{FAFAF7}
  \definecolor{bnMargin} {HTML}{B5B0A5}
  \definecolor{bnMuted}  {HTML}{6F665A}
\fi

% --- Logo geometry (all in cm; tweak freely) --------------------------------

\providecommand{\bnBarCount}{4}     % number of horizontal bars
\providecommand{\bnBarWidth}{4.0}   % length of an in-margin bar
\providecommand{\bnBarHeight}{0.55} % bar thickness
\providecommand{\bnBarGap}{0.32}    % vertical gap between bars
\providecommand{\bnJutWhich}{3}     % 1..bnBarCount — which bar overflows
\providecommand{\bnJutExtra}{1.15}  % how far it sticks out
\providecommand{\bnRound}{0.07}     % corner rounding
\providecommand{\bnShowMargin}{0}   % 1 = draw the dashed margin guide, 0 = hide

% --- The mark itself --------------------------------------------------------
%
% Draws into the current tikzpicture, anchored so the bottom-left of the
% non-jutting bars is at (0,0). The jutting bar extends to the right.

\newcommand{\bnDrawMark}{%
  \pgfmathtruncatemacro{\bnLast}{\bnBarCount-1}%
  \pgfmathsetmacro{\bnTotalH}{\bnBarCount*\bnBarHeight + \bnLast*\bnBarGap}%
  \ifnum\bnShowMargin=1
    % The right-margin guide: a dashed vertical line at the in-margin width,
    % running the full height of the bar stack with a touch of overshoot.
    \draw[bnMargin, dashed, line width=0.4pt, dash pattern=on 2pt off 2pt]
    (\bnBarWidth, -0.12) -- (\bnBarWidth, \bnTotalH+0.12);
  \fi
  \foreach \i in {0,...,\bnLast}{%
      \pgfmathsetmacro{\bny}{(\bnLast-\i)*(\bnBarHeight+\bnBarGap)}%
      \pgfmathtruncatemacro{\bnNum}{\i+1}%
      \ifnum\bnNum=\bnJutWhich
        \fill[bnBad, rounded corners=\bnRound cm]
        (0,\bny) rectangle (\bnBarWidth+\bnJutExtra, \bny+\bnBarHeight);
      \else
        \fill[bnInk, rounded corners=\bnRound cm]
        (0,\bny) rectangle (\bnBarWidth, \bny+\bnBarHeight);
      \fi
    }%
}
.

\usepackage{tikz}
\usepackage{xcolor}

% --- Font dispatch ----------------------------------------------------------
%
% \bnFont selects the sans-serif family. Set it in your artifact file
% before % Shared definitions for the badness brand artifacts.
%
% \input by the per-artifact .tex files (logo.tex, wordmark.tex, og.tex).
% Anything \providecommand'd here is overridable by the caller — define
% your override BEFORE \input{_common.tex}.

\usepackage{tikz}
\usepackage{xcolor}

% --- Font dispatch ----------------------------------------------------------
%
% \bnFont selects the sans-serif family. Set it in your artifact file
% before \input{_common.tex}. Available values:
%
%   plex       IBM Plex Sans         — modern technical (default)
%   tgheros    TeX Gyre Heros        — Helvetica-clone, neutral
%   source     Source Sans Pro       — humanist, friendly
%   roboto     Roboto                — Google flat-ui geometric
%   fira       Fira Sans             — Mozilla, slightly quirky
%   inter      Inter                 — UI workhorse, very tight
%   montserrat Montserrat            — wide geometric, posterish
%   lato       Lato                  — warm humanist
%   cmbright   Computer Modern Bright — light, scholarly
%   lm         Latin Modern Sans     — LaTeX default, neutral
%
\providecommand{\bnFont}{lm}

\makeatletter
\def\bn@font@plex      {\RequirePackage{plex-sans}}
\def\bn@font@tgheros   {\RequirePackage{tgheros}}
\def\bn@font@source    {\RequirePackage[default]{sourcesanspro}}
\def\bn@font@roboto    {\RequirePackage[sfdefault]{roboto}}
\def\bn@font@fira      {\RequirePackage[sfdefault]{FiraSans}}
\def\bn@font@inter     {\RequirePackage{inter}}
\def\bn@font@montserrat{\RequirePackage[defaultfam,tabular,lining]{montserrat}}
\def\bn@font@lato      {\RequirePackage[default]{lato}}
\def\bn@font@cmbright  {\RequirePackage{cmbright}}
\def\bn@font@lm        {\RequirePackage{lmodern}}
\@nameuse{bn@font@\bnFont}
\makeatother
\renewcommand{\familydefault}{\sfdefault}

% --- Palette ----------------------------------------------------------------
%
% Five role-named slots. \bnDarkMode=1 swaps to a dark variant.
% Override any individual color by \definecolor'ing it AFTER \input{_common.tex}.

\providecommand{\bnDarkMode}{0}

\ifnum\bnDarkMode=1
  % Dark variant
  \definecolor{bnInk}    {HTML}{F2EFE9} % "ink" reads light on dark
  \definecolor{bnBad}    {HTML}{F08A57} % brighter orange so it pops on dark
  \definecolor{bnPaper}  {HTML}{1B1B1B} % near-black canvas
  \definecolor{bnMargin} {HTML}{3D3833} % subtle margin guide
  \definecolor{bnMuted}  {HTML}{A39A8C} % secondary text
\else
  % Light variant (default)
  \definecolor{bnInk}    {HTML}{1B1B1B}
  \definecolor{bnBad}    {HTML}{D9633F}
  \definecolor{bnPaper}  {HTML}{FAFAF7}
  \definecolor{bnMargin} {HTML}{B5B0A5}
  \definecolor{bnMuted}  {HTML}{6F665A}
\fi

% --- Logo geometry (all in cm; tweak freely) --------------------------------

\providecommand{\bnBarCount}{4}     % number of horizontal bars
\providecommand{\bnBarWidth}{4.0}   % length of an in-margin bar
\providecommand{\bnBarHeight}{0.55} % bar thickness
\providecommand{\bnBarGap}{0.32}    % vertical gap between bars
\providecommand{\bnJutWhich}{3}     % 1..bnBarCount — which bar overflows
\providecommand{\bnJutExtra}{1.15}  % how far it sticks out
\providecommand{\bnRound}{0.07}     % corner rounding
\providecommand{\bnShowMargin}{0}   % 1 = draw the dashed margin guide, 0 = hide

% --- The mark itself --------------------------------------------------------
%
% Draws into the current tikzpicture, anchored so the bottom-left of the
% non-jutting bars is at (0,0). The jutting bar extends to the right.

\newcommand{\bnDrawMark}{%
  \pgfmathtruncatemacro{\bnLast}{\bnBarCount-1}%
  \pgfmathsetmacro{\bnTotalH}{\bnBarCount*\bnBarHeight + \bnLast*\bnBarGap}%
  \ifnum\bnShowMargin=1
    % The right-margin guide: a dashed vertical line at the in-margin width,
    % running the full height of the bar stack with a touch of overshoot.
    \draw[bnMargin, dashed, line width=0.4pt, dash pattern=on 2pt off 2pt]
    (\bnBarWidth, -0.12) -- (\bnBarWidth, \bnTotalH+0.12);
  \fi
  \foreach \i in {0,...,\bnLast}{%
      \pgfmathsetmacro{\bny}{(\bnLast-\i)*(\bnBarHeight+\bnBarGap)}%
      \pgfmathtruncatemacro{\bnNum}{\i+1}%
      \ifnum\bnNum=\bnJutWhich
        \fill[bnBad, rounded corners=\bnRound cm]
        (0,\bny) rectangle (\bnBarWidth+\bnJutExtra, \bny+\bnBarHeight);
      \else
        \fill[bnInk, rounded corners=\bnRound cm]
        (0,\bny) rectangle (\bnBarWidth, \bny+\bnBarHeight);
      \fi
    }%
}
. Available values:
%
%   plex       IBM Plex Sans         — modern technical (default)
%   tgheros    TeX Gyre Heros        — Helvetica-clone, neutral
%   source     Source Sans Pro       — humanist, friendly
%   roboto     Roboto                — Google flat-ui geometric
%   fira       Fira Sans             — Mozilla, slightly quirky
%   inter      Inter                 — UI workhorse, very tight
%   montserrat Montserrat            — wide geometric, posterish
%   lato       Lato                  — warm humanist
%   cmbright   Computer Modern Bright — light, scholarly
%   lm         Latin Modern Sans     — LaTeX default, neutral
%
\providecommand{\bnFont}{lm}

\makeatletter
\def\bn@font@plex      {\RequirePackage{plex-sans}}
\def\bn@font@tgheros   {\RequirePackage{tgheros}}
\def\bn@font@source    {\RequirePackage[default]{sourcesanspro}}
\def\bn@font@roboto    {\RequirePackage[sfdefault]{roboto}}
\def\bn@font@fira      {\RequirePackage[sfdefault]{FiraSans}}
\def\bn@font@inter     {\RequirePackage{inter}}
\def\bn@font@montserrat{\RequirePackage[defaultfam,tabular,lining]{montserrat}}
\def\bn@font@lato      {\RequirePackage[default]{lato}}
\def\bn@font@cmbright  {\RequirePackage{cmbright}}
\def\bn@font@lm        {\RequirePackage{lmodern}}
\@nameuse{bn@font@\bnFont}
\makeatother
\renewcommand{\familydefault}{\sfdefault}

% --- Palette ----------------------------------------------------------------
%
% Five role-named slots. \bnDarkMode=1 swaps to a dark variant.
% Override any individual color by \definecolor'ing it AFTER % Shared definitions for the badness brand artifacts.
%
% \input by the per-artifact .tex files (logo.tex, wordmark.tex, og.tex).
% Anything \providecommand'd here is overridable by the caller — define
% your override BEFORE \input{_common.tex}.

\usepackage{tikz}
\usepackage{xcolor}

% --- Font dispatch ----------------------------------------------------------
%
% \bnFont selects the sans-serif family. Set it in your artifact file
% before \input{_common.tex}. Available values:
%
%   plex       IBM Plex Sans         — modern technical (default)
%   tgheros    TeX Gyre Heros        — Helvetica-clone, neutral
%   source     Source Sans Pro       — humanist, friendly
%   roboto     Roboto                — Google flat-ui geometric
%   fira       Fira Sans             — Mozilla, slightly quirky
%   inter      Inter                 — UI workhorse, very tight
%   montserrat Montserrat            — wide geometric, posterish
%   lato       Lato                  — warm humanist
%   cmbright   Computer Modern Bright — light, scholarly
%   lm         Latin Modern Sans     — LaTeX default, neutral
%
\providecommand{\bnFont}{lm}

\makeatletter
\def\bn@font@plex      {\RequirePackage{plex-sans}}
\def\bn@font@tgheros   {\RequirePackage{tgheros}}
\def\bn@font@source    {\RequirePackage[default]{sourcesanspro}}
\def\bn@font@roboto    {\RequirePackage[sfdefault]{roboto}}
\def\bn@font@fira      {\RequirePackage[sfdefault]{FiraSans}}
\def\bn@font@inter     {\RequirePackage{inter}}
\def\bn@font@montserrat{\RequirePackage[defaultfam,tabular,lining]{montserrat}}
\def\bn@font@lato      {\RequirePackage[default]{lato}}
\def\bn@font@cmbright  {\RequirePackage{cmbright}}
\def\bn@font@lm        {\RequirePackage{lmodern}}
\@nameuse{bn@font@\bnFont}
\makeatother
\renewcommand{\familydefault}{\sfdefault}

% --- Palette ----------------------------------------------------------------
%
% Five role-named slots. \bnDarkMode=1 swaps to a dark variant.
% Override any individual color by \definecolor'ing it AFTER \input{_common.tex}.

\providecommand{\bnDarkMode}{0}

\ifnum\bnDarkMode=1
  % Dark variant
  \definecolor{bnInk}    {HTML}{F2EFE9} % "ink" reads light on dark
  \definecolor{bnBad}    {HTML}{F08A57} % brighter orange so it pops on dark
  \definecolor{bnPaper}  {HTML}{1B1B1B} % near-black canvas
  \definecolor{bnMargin} {HTML}{3D3833} % subtle margin guide
  \definecolor{bnMuted}  {HTML}{A39A8C} % secondary text
\else
  % Light variant (default)
  \definecolor{bnInk}    {HTML}{1B1B1B}
  \definecolor{bnBad}    {HTML}{D9633F}
  \definecolor{bnPaper}  {HTML}{FAFAF7}
  \definecolor{bnMargin} {HTML}{B5B0A5}
  \definecolor{bnMuted}  {HTML}{6F665A}
\fi

% --- Logo geometry (all in cm; tweak freely) --------------------------------

\providecommand{\bnBarCount}{4}     % number of horizontal bars
\providecommand{\bnBarWidth}{4.0}   % length of an in-margin bar
\providecommand{\bnBarHeight}{0.55} % bar thickness
\providecommand{\bnBarGap}{0.32}    % vertical gap between bars
\providecommand{\bnJutWhich}{3}     % 1..bnBarCount — which bar overflows
\providecommand{\bnJutExtra}{1.15}  % how far it sticks out
\providecommand{\bnRound}{0.07}     % corner rounding
\providecommand{\bnShowMargin}{0}   % 1 = draw the dashed margin guide, 0 = hide

% --- The mark itself --------------------------------------------------------
%
% Draws into the current tikzpicture, anchored so the bottom-left of the
% non-jutting bars is at (0,0). The jutting bar extends to the right.

\newcommand{\bnDrawMark}{%
  \pgfmathtruncatemacro{\bnLast}{\bnBarCount-1}%
  \pgfmathsetmacro{\bnTotalH}{\bnBarCount*\bnBarHeight + \bnLast*\bnBarGap}%
  \ifnum\bnShowMargin=1
    % The right-margin guide: a dashed vertical line at the in-margin width,
    % running the full height of the bar stack with a touch of overshoot.
    \draw[bnMargin, dashed, line width=0.4pt, dash pattern=on 2pt off 2pt]
    (\bnBarWidth, -0.12) -- (\bnBarWidth, \bnTotalH+0.12);
  \fi
  \foreach \i in {0,...,\bnLast}{%
      \pgfmathsetmacro{\bny}{(\bnLast-\i)*(\bnBarHeight+\bnBarGap)}%
      \pgfmathtruncatemacro{\bnNum}{\i+1}%
      \ifnum\bnNum=\bnJutWhich
        \fill[bnBad, rounded corners=\bnRound cm]
        (0,\bny) rectangle (\bnBarWidth+\bnJutExtra, \bny+\bnBarHeight);
      \else
        \fill[bnInk, rounded corners=\bnRound cm]
        (0,\bny) rectangle (\bnBarWidth, \bny+\bnBarHeight);
      \fi
    }%
}
.

\providecommand{\bnDarkMode}{0}

\ifnum\bnDarkMode=1
  % Dark variant
  \definecolor{bnInk}    {HTML}{F2EFE9} % "ink" reads light on dark
  \definecolor{bnBad}    {HTML}{F08A57} % brighter orange so it pops on dark
  \definecolor{bnPaper}  {HTML}{1B1B1B} % near-black canvas
  \definecolor{bnMargin} {HTML}{3D3833} % subtle margin guide
  \definecolor{bnMuted}  {HTML}{A39A8C} % secondary text
\else
  % Light variant (default)
  \definecolor{bnInk}    {HTML}{1B1B1B}
  \definecolor{bnBad}    {HTML}{D9633F}
  \definecolor{bnPaper}  {HTML}{FAFAF7}
  \definecolor{bnMargin} {HTML}{B5B0A5}
  \definecolor{bnMuted}  {HTML}{6F665A}
\fi

% --- Logo geometry (all in cm; tweak freely) --------------------------------

\providecommand{\bnBarCount}{4}     % number of horizontal bars
\providecommand{\bnBarWidth}{4.0}   % length of an in-margin bar
\providecommand{\bnBarHeight}{0.55} % bar thickness
\providecommand{\bnBarGap}{0.32}    % vertical gap between bars
\providecommand{\bnJutWhich}{3}     % 1..bnBarCount — which bar overflows
\providecommand{\bnJutExtra}{1.15}  % how far it sticks out
\providecommand{\bnRound}{0.07}     % corner rounding
\providecommand{\bnShowMargin}{0}   % 1 = draw the dashed margin guide, 0 = hide

% --- The mark itself --------------------------------------------------------
%
% Draws into the current tikzpicture, anchored so the bottom-left of the
% non-jutting bars is at (0,0). The jutting bar extends to the right.

\newcommand{\bnDrawMark}{%
  \pgfmathtruncatemacro{\bnLast}{\bnBarCount-1}%
  \pgfmathsetmacro{\bnTotalH}{\bnBarCount*\bnBarHeight + \bnLast*\bnBarGap}%
  \ifnum\bnShowMargin=1
    % The right-margin guide: a dashed vertical line at the in-margin width,
    % running the full height of the bar stack with a touch of overshoot.
    \draw[bnMargin, dashed, line width=0.4pt, dash pattern=on 2pt off 2pt]
    (\bnBarWidth, -0.12) -- (\bnBarWidth, \bnTotalH+0.12);
  \fi
  \foreach \i in {0,...,\bnLast}{%
      \pgfmathsetmacro{\bny}{(\bnLast-\i)*(\bnBarHeight+\bnBarGap)}%
      \pgfmathtruncatemacro{\bnNum}{\i+1}%
      \ifnum\bnNum=\bnJutWhich
        \fill[bnBad, rounded corners=\bnRound cm]
        (0,\bny) rectangle (\bnBarWidth+\bnJutExtra, \bny+\bnBarHeight);
      \else
        \fill[bnInk, rounded corners=\bnRound cm]
        (0,\bny) rectangle (\bnBarWidth, \bny+\bnBarHeight);
      \fi
    }%
}
.

\usepackage{tikz}
\usepackage{xcolor}

% --- Font dispatch ----------------------------------------------------------
%
% \bnFont selects the sans-serif family. Set it in your artifact file
% before % Shared definitions for the badness brand artifacts.
%
% \input by the per-artifact .tex files (logo.tex, wordmark.tex, og.tex).
% Anything \providecommand'd here is overridable by the caller — define
% your override BEFORE % Shared definitions for the badness brand artifacts.
%
% \input by the per-artifact .tex files (logo.tex, wordmark.tex, og.tex).
% Anything \providecommand'd here is overridable by the caller — define
% your override BEFORE \input{_common.tex}.

\usepackage{tikz}
\usepackage{xcolor}

% --- Font dispatch ----------------------------------------------------------
%
% \bnFont selects the sans-serif family. Set it in your artifact file
% before \input{_common.tex}. Available values:
%
%   plex       IBM Plex Sans         — modern technical (default)
%   tgheros    TeX Gyre Heros        — Helvetica-clone, neutral
%   source     Source Sans Pro       — humanist, friendly
%   roboto     Roboto                — Google flat-ui geometric
%   fira       Fira Sans             — Mozilla, slightly quirky
%   inter      Inter                 — UI workhorse, very tight
%   montserrat Montserrat            — wide geometric, posterish
%   lato       Lato                  — warm humanist
%   cmbright   Computer Modern Bright — light, scholarly
%   lm         Latin Modern Sans     — LaTeX default, neutral
%
\providecommand{\bnFont}{lm}

\makeatletter
\def\bn@font@plex      {\RequirePackage{plex-sans}}
\def\bn@font@tgheros   {\RequirePackage{tgheros}}
\def\bn@font@source    {\RequirePackage[default]{sourcesanspro}}
\def\bn@font@roboto    {\RequirePackage[sfdefault]{roboto}}
\def\bn@font@fira      {\RequirePackage[sfdefault]{FiraSans}}
\def\bn@font@inter     {\RequirePackage{inter}}
\def\bn@font@montserrat{\RequirePackage[defaultfam,tabular,lining]{montserrat}}
\def\bn@font@lato      {\RequirePackage[default]{lato}}
\def\bn@font@cmbright  {\RequirePackage{cmbright}}
\def\bn@font@lm        {\RequirePackage{lmodern}}
\@nameuse{bn@font@\bnFont}
\makeatother
\renewcommand{\familydefault}{\sfdefault}

% --- Palette ----------------------------------------------------------------
%
% Five role-named slots. \bnDarkMode=1 swaps to a dark variant.
% Override any individual color by \definecolor'ing it AFTER \input{_common.tex}.

\providecommand{\bnDarkMode}{0}

\ifnum\bnDarkMode=1
  % Dark variant
  \definecolor{bnInk}    {HTML}{F2EFE9} % "ink" reads light on dark
  \definecolor{bnBad}    {HTML}{F08A57} % brighter orange so it pops on dark
  \definecolor{bnPaper}  {HTML}{1B1B1B} % near-black canvas
  \definecolor{bnMargin} {HTML}{3D3833} % subtle margin guide
  \definecolor{bnMuted}  {HTML}{A39A8C} % secondary text
\else
  % Light variant (default)
  \definecolor{bnInk}    {HTML}{1B1B1B}
  \definecolor{bnBad}    {HTML}{D9633F}
  \definecolor{bnPaper}  {HTML}{FAFAF7}
  \definecolor{bnMargin} {HTML}{B5B0A5}
  \definecolor{bnMuted}  {HTML}{6F665A}
\fi

% --- Logo geometry (all in cm; tweak freely) --------------------------------

\providecommand{\bnBarCount}{4}     % number of horizontal bars
\providecommand{\bnBarWidth}{4.0}   % length of an in-margin bar
\providecommand{\bnBarHeight}{0.55} % bar thickness
\providecommand{\bnBarGap}{0.32}    % vertical gap between bars
\providecommand{\bnJutWhich}{3}     % 1..bnBarCount — which bar overflows
\providecommand{\bnJutExtra}{1.15}  % how far it sticks out
\providecommand{\bnRound}{0.07}     % corner rounding
\providecommand{\bnShowMargin}{0}   % 1 = draw the dashed margin guide, 0 = hide

% --- The mark itself --------------------------------------------------------
%
% Draws into the current tikzpicture, anchored so the bottom-left of the
% non-jutting bars is at (0,0). The jutting bar extends to the right.

\newcommand{\bnDrawMark}{%
  \pgfmathtruncatemacro{\bnLast}{\bnBarCount-1}%
  \pgfmathsetmacro{\bnTotalH}{\bnBarCount*\bnBarHeight + \bnLast*\bnBarGap}%
  \ifnum\bnShowMargin=1
    % The right-margin guide: a dashed vertical line at the in-margin width,
    % running the full height of the bar stack with a touch of overshoot.
    \draw[bnMargin, dashed, line width=0.4pt, dash pattern=on 2pt off 2pt]
    (\bnBarWidth, -0.12) -- (\bnBarWidth, \bnTotalH+0.12);
  \fi
  \foreach \i in {0,...,\bnLast}{%
      \pgfmathsetmacro{\bny}{(\bnLast-\i)*(\bnBarHeight+\bnBarGap)}%
      \pgfmathtruncatemacro{\bnNum}{\i+1}%
      \ifnum\bnNum=\bnJutWhich
        \fill[bnBad, rounded corners=\bnRound cm]
        (0,\bny) rectangle (\bnBarWidth+\bnJutExtra, \bny+\bnBarHeight);
      \else
        \fill[bnInk, rounded corners=\bnRound cm]
        (0,\bny) rectangle (\bnBarWidth, \bny+\bnBarHeight);
      \fi
    }%
}
.

\usepackage{tikz}
\usepackage{xcolor}

% --- Font dispatch ----------------------------------------------------------
%
% \bnFont selects the sans-serif family. Set it in your artifact file
% before % Shared definitions for the badness brand artifacts.
%
% \input by the per-artifact .tex files (logo.tex, wordmark.tex, og.tex).
% Anything \providecommand'd here is overridable by the caller — define
% your override BEFORE \input{_common.tex}.

\usepackage{tikz}
\usepackage{xcolor}

% --- Font dispatch ----------------------------------------------------------
%
% \bnFont selects the sans-serif family. Set it in your artifact file
% before \input{_common.tex}. Available values:
%
%   plex       IBM Plex Sans         — modern technical (default)
%   tgheros    TeX Gyre Heros        — Helvetica-clone, neutral
%   source     Source Sans Pro       — humanist, friendly
%   roboto     Roboto                — Google flat-ui geometric
%   fira       Fira Sans             — Mozilla, slightly quirky
%   inter      Inter                 — UI workhorse, very tight
%   montserrat Montserrat            — wide geometric, posterish
%   lato       Lato                  — warm humanist
%   cmbright   Computer Modern Bright — light, scholarly
%   lm         Latin Modern Sans     — LaTeX default, neutral
%
\providecommand{\bnFont}{lm}

\makeatletter
\def\bn@font@plex      {\RequirePackage{plex-sans}}
\def\bn@font@tgheros   {\RequirePackage{tgheros}}
\def\bn@font@source    {\RequirePackage[default]{sourcesanspro}}
\def\bn@font@roboto    {\RequirePackage[sfdefault]{roboto}}
\def\bn@font@fira      {\RequirePackage[sfdefault]{FiraSans}}
\def\bn@font@inter     {\RequirePackage{inter}}
\def\bn@font@montserrat{\RequirePackage[defaultfam,tabular,lining]{montserrat}}
\def\bn@font@lato      {\RequirePackage[default]{lato}}
\def\bn@font@cmbright  {\RequirePackage{cmbright}}
\def\bn@font@lm        {\RequirePackage{lmodern}}
\@nameuse{bn@font@\bnFont}
\makeatother
\renewcommand{\familydefault}{\sfdefault}

% --- Palette ----------------------------------------------------------------
%
% Five role-named slots. \bnDarkMode=1 swaps to a dark variant.
% Override any individual color by \definecolor'ing it AFTER \input{_common.tex}.

\providecommand{\bnDarkMode}{0}

\ifnum\bnDarkMode=1
  % Dark variant
  \definecolor{bnInk}    {HTML}{F2EFE9} % "ink" reads light on dark
  \definecolor{bnBad}    {HTML}{F08A57} % brighter orange so it pops on dark
  \definecolor{bnPaper}  {HTML}{1B1B1B} % near-black canvas
  \definecolor{bnMargin} {HTML}{3D3833} % subtle margin guide
  \definecolor{bnMuted}  {HTML}{A39A8C} % secondary text
\else
  % Light variant (default)
  \definecolor{bnInk}    {HTML}{1B1B1B}
  \definecolor{bnBad}    {HTML}{D9633F}
  \definecolor{bnPaper}  {HTML}{FAFAF7}
  \definecolor{bnMargin} {HTML}{B5B0A5}
  \definecolor{bnMuted}  {HTML}{6F665A}
\fi

% --- Logo geometry (all in cm; tweak freely) --------------------------------

\providecommand{\bnBarCount}{4}     % number of horizontal bars
\providecommand{\bnBarWidth}{4.0}   % length of an in-margin bar
\providecommand{\bnBarHeight}{0.55} % bar thickness
\providecommand{\bnBarGap}{0.32}    % vertical gap between bars
\providecommand{\bnJutWhich}{3}     % 1..bnBarCount — which bar overflows
\providecommand{\bnJutExtra}{1.15}  % how far it sticks out
\providecommand{\bnRound}{0.07}     % corner rounding
\providecommand{\bnShowMargin}{0}   % 1 = draw the dashed margin guide, 0 = hide

% --- The mark itself --------------------------------------------------------
%
% Draws into the current tikzpicture, anchored so the bottom-left of the
% non-jutting bars is at (0,0). The jutting bar extends to the right.

\newcommand{\bnDrawMark}{%
  \pgfmathtruncatemacro{\bnLast}{\bnBarCount-1}%
  \pgfmathsetmacro{\bnTotalH}{\bnBarCount*\bnBarHeight + \bnLast*\bnBarGap}%
  \ifnum\bnShowMargin=1
    % The right-margin guide: a dashed vertical line at the in-margin width,
    % running the full height of the bar stack with a touch of overshoot.
    \draw[bnMargin, dashed, line width=0.4pt, dash pattern=on 2pt off 2pt]
    (\bnBarWidth, -0.12) -- (\bnBarWidth, \bnTotalH+0.12);
  \fi
  \foreach \i in {0,...,\bnLast}{%
      \pgfmathsetmacro{\bny}{(\bnLast-\i)*(\bnBarHeight+\bnBarGap)}%
      \pgfmathtruncatemacro{\bnNum}{\i+1}%
      \ifnum\bnNum=\bnJutWhich
        \fill[bnBad, rounded corners=\bnRound cm]
        (0,\bny) rectangle (\bnBarWidth+\bnJutExtra, \bny+\bnBarHeight);
      \else
        \fill[bnInk, rounded corners=\bnRound cm]
        (0,\bny) rectangle (\bnBarWidth, \bny+\bnBarHeight);
      \fi
    }%
}
. Available values:
%
%   plex       IBM Plex Sans         — modern technical (default)
%   tgheros    TeX Gyre Heros        — Helvetica-clone, neutral
%   source     Source Sans Pro       — humanist, friendly
%   roboto     Roboto                — Google flat-ui geometric
%   fira       Fira Sans             — Mozilla, slightly quirky
%   inter      Inter                 — UI workhorse, very tight
%   montserrat Montserrat            — wide geometric, posterish
%   lato       Lato                  — warm humanist
%   cmbright   Computer Modern Bright — light, scholarly
%   lm         Latin Modern Sans     — LaTeX default, neutral
%
\providecommand{\bnFont}{lm}

\makeatletter
\def\bn@font@plex      {\RequirePackage{plex-sans}}
\def\bn@font@tgheros   {\RequirePackage{tgheros}}
\def\bn@font@source    {\RequirePackage[default]{sourcesanspro}}
\def\bn@font@roboto    {\RequirePackage[sfdefault]{roboto}}
\def\bn@font@fira      {\RequirePackage[sfdefault]{FiraSans}}
\def\bn@font@inter     {\RequirePackage{inter}}
\def\bn@font@montserrat{\RequirePackage[defaultfam,tabular,lining]{montserrat}}
\def\bn@font@lato      {\RequirePackage[default]{lato}}
\def\bn@font@cmbright  {\RequirePackage{cmbright}}
\def\bn@font@lm        {\RequirePackage{lmodern}}
\@nameuse{bn@font@\bnFont}
\makeatother
\renewcommand{\familydefault}{\sfdefault}

% --- Palette ----------------------------------------------------------------
%
% Five role-named slots. \bnDarkMode=1 swaps to a dark variant.
% Override any individual color by \definecolor'ing it AFTER % Shared definitions for the badness brand artifacts.
%
% \input by the per-artifact .tex files (logo.tex, wordmark.tex, og.tex).
% Anything \providecommand'd here is overridable by the caller — define
% your override BEFORE \input{_common.tex}.

\usepackage{tikz}
\usepackage{xcolor}

% --- Font dispatch ----------------------------------------------------------
%
% \bnFont selects the sans-serif family. Set it in your artifact file
% before \input{_common.tex}. Available values:
%
%   plex       IBM Plex Sans         — modern technical (default)
%   tgheros    TeX Gyre Heros        — Helvetica-clone, neutral
%   source     Source Sans Pro       — humanist, friendly
%   roboto     Roboto                — Google flat-ui geometric
%   fira       Fira Sans             — Mozilla, slightly quirky
%   inter      Inter                 — UI workhorse, very tight
%   montserrat Montserrat            — wide geometric, posterish
%   lato       Lato                  — warm humanist
%   cmbright   Computer Modern Bright — light, scholarly
%   lm         Latin Modern Sans     — LaTeX default, neutral
%
\providecommand{\bnFont}{lm}

\makeatletter
\def\bn@font@plex      {\RequirePackage{plex-sans}}
\def\bn@font@tgheros   {\RequirePackage{tgheros}}
\def\bn@font@source    {\RequirePackage[default]{sourcesanspro}}
\def\bn@font@roboto    {\RequirePackage[sfdefault]{roboto}}
\def\bn@font@fira      {\RequirePackage[sfdefault]{FiraSans}}
\def\bn@font@inter     {\RequirePackage{inter}}
\def\bn@font@montserrat{\RequirePackage[defaultfam,tabular,lining]{montserrat}}
\def\bn@font@lato      {\RequirePackage[default]{lato}}
\def\bn@font@cmbright  {\RequirePackage{cmbright}}
\def\bn@font@lm        {\RequirePackage{lmodern}}
\@nameuse{bn@font@\bnFont}
\makeatother
\renewcommand{\familydefault}{\sfdefault}

% --- Palette ----------------------------------------------------------------
%
% Five role-named slots. \bnDarkMode=1 swaps to a dark variant.
% Override any individual color by \definecolor'ing it AFTER \input{_common.tex}.

\providecommand{\bnDarkMode}{0}

\ifnum\bnDarkMode=1
  % Dark variant
  \definecolor{bnInk}    {HTML}{F2EFE9} % "ink" reads light on dark
  \definecolor{bnBad}    {HTML}{F08A57} % brighter orange so it pops on dark
  \definecolor{bnPaper}  {HTML}{1B1B1B} % near-black canvas
  \definecolor{bnMargin} {HTML}{3D3833} % subtle margin guide
  \definecolor{bnMuted}  {HTML}{A39A8C} % secondary text
\else
  % Light variant (default)
  \definecolor{bnInk}    {HTML}{1B1B1B}
  \definecolor{bnBad}    {HTML}{D9633F}
  \definecolor{bnPaper}  {HTML}{FAFAF7}
  \definecolor{bnMargin} {HTML}{B5B0A5}
  \definecolor{bnMuted}  {HTML}{6F665A}
\fi

% --- Logo geometry (all in cm; tweak freely) --------------------------------

\providecommand{\bnBarCount}{4}     % number of horizontal bars
\providecommand{\bnBarWidth}{4.0}   % length of an in-margin bar
\providecommand{\bnBarHeight}{0.55} % bar thickness
\providecommand{\bnBarGap}{0.32}    % vertical gap between bars
\providecommand{\bnJutWhich}{3}     % 1..bnBarCount — which bar overflows
\providecommand{\bnJutExtra}{1.15}  % how far it sticks out
\providecommand{\bnRound}{0.07}     % corner rounding
\providecommand{\bnShowMargin}{0}   % 1 = draw the dashed margin guide, 0 = hide

% --- The mark itself --------------------------------------------------------
%
% Draws into the current tikzpicture, anchored so the bottom-left of the
% non-jutting bars is at (0,0). The jutting bar extends to the right.

\newcommand{\bnDrawMark}{%
  \pgfmathtruncatemacro{\bnLast}{\bnBarCount-1}%
  \pgfmathsetmacro{\bnTotalH}{\bnBarCount*\bnBarHeight + \bnLast*\bnBarGap}%
  \ifnum\bnShowMargin=1
    % The right-margin guide: a dashed vertical line at the in-margin width,
    % running the full height of the bar stack with a touch of overshoot.
    \draw[bnMargin, dashed, line width=0.4pt, dash pattern=on 2pt off 2pt]
    (\bnBarWidth, -0.12) -- (\bnBarWidth, \bnTotalH+0.12);
  \fi
  \foreach \i in {0,...,\bnLast}{%
      \pgfmathsetmacro{\bny}{(\bnLast-\i)*(\bnBarHeight+\bnBarGap)}%
      \pgfmathtruncatemacro{\bnNum}{\i+1}%
      \ifnum\bnNum=\bnJutWhich
        \fill[bnBad, rounded corners=\bnRound cm]
        (0,\bny) rectangle (\bnBarWidth+\bnJutExtra, \bny+\bnBarHeight);
      \else
        \fill[bnInk, rounded corners=\bnRound cm]
        (0,\bny) rectangle (\bnBarWidth, \bny+\bnBarHeight);
      \fi
    }%
}
.

\providecommand{\bnDarkMode}{0}

\ifnum\bnDarkMode=1
  % Dark variant
  \definecolor{bnInk}    {HTML}{F2EFE9} % "ink" reads light on dark
  \definecolor{bnBad}    {HTML}{F08A57} % brighter orange so it pops on dark
  \definecolor{bnPaper}  {HTML}{1B1B1B} % near-black canvas
  \definecolor{bnMargin} {HTML}{3D3833} % subtle margin guide
  \definecolor{bnMuted}  {HTML}{A39A8C} % secondary text
\else
  % Light variant (default)
  \definecolor{bnInk}    {HTML}{1B1B1B}
  \definecolor{bnBad}    {HTML}{D9633F}
  \definecolor{bnPaper}  {HTML}{FAFAF7}
  \definecolor{bnMargin} {HTML}{B5B0A5}
  \definecolor{bnMuted}  {HTML}{6F665A}
\fi

% --- Logo geometry (all in cm; tweak freely) --------------------------------

\providecommand{\bnBarCount}{4}     % number of horizontal bars
\providecommand{\bnBarWidth}{4.0}   % length of an in-margin bar
\providecommand{\bnBarHeight}{0.55} % bar thickness
\providecommand{\bnBarGap}{0.32}    % vertical gap between bars
\providecommand{\bnJutWhich}{3}     % 1..bnBarCount — which bar overflows
\providecommand{\bnJutExtra}{1.15}  % how far it sticks out
\providecommand{\bnRound}{0.07}     % corner rounding
\providecommand{\bnShowMargin}{0}   % 1 = draw the dashed margin guide, 0 = hide

% --- The mark itself --------------------------------------------------------
%
% Draws into the current tikzpicture, anchored so the bottom-left of the
% non-jutting bars is at (0,0). The jutting bar extends to the right.

\newcommand{\bnDrawMark}{%
  \pgfmathtruncatemacro{\bnLast}{\bnBarCount-1}%
  \pgfmathsetmacro{\bnTotalH}{\bnBarCount*\bnBarHeight + \bnLast*\bnBarGap}%
  \ifnum\bnShowMargin=1
    % The right-margin guide: a dashed vertical line at the in-margin width,
    % running the full height of the bar stack with a touch of overshoot.
    \draw[bnMargin, dashed, line width=0.4pt, dash pattern=on 2pt off 2pt]
    (\bnBarWidth, -0.12) -- (\bnBarWidth, \bnTotalH+0.12);
  \fi
  \foreach \i in {0,...,\bnLast}{%
      \pgfmathsetmacro{\bny}{(\bnLast-\i)*(\bnBarHeight+\bnBarGap)}%
      \pgfmathtruncatemacro{\bnNum}{\i+1}%
      \ifnum\bnNum=\bnJutWhich
        \fill[bnBad, rounded corners=\bnRound cm]
        (0,\bny) rectangle (\bnBarWidth+\bnJutExtra, \bny+\bnBarHeight);
      \else
        \fill[bnInk, rounded corners=\bnRound cm]
        (0,\bny) rectangle (\bnBarWidth, \bny+\bnBarHeight);
      \fi
    }%
}
. Available values:
%
%   plex       IBM Plex Sans         — modern technical (default)
%   tgheros    TeX Gyre Heros        — Helvetica-clone, neutral
%   source     Source Sans Pro       — humanist, friendly
%   roboto     Roboto                — Google flat-ui geometric
%   fira       Fira Sans             — Mozilla, slightly quirky
%   inter      Inter                 — UI workhorse, very tight
%   montserrat Montserrat            — wide geometric, posterish
%   lato       Lato                  — warm humanist
%   cmbright   Computer Modern Bright — light, scholarly
%   lm         Latin Modern Sans     — LaTeX default, neutral
%
\providecommand{\bnFont}{lm}

\makeatletter
\def\bn@font@plex      {\RequirePackage{plex-sans}}
\def\bn@font@tgheros   {\RequirePackage{tgheros}}
\def\bn@font@source    {\RequirePackage[default]{sourcesanspro}}
\def\bn@font@roboto    {\RequirePackage[sfdefault]{roboto}}
\def\bn@font@fira      {\RequirePackage[sfdefault]{FiraSans}}
\def\bn@font@inter     {\RequirePackage{inter}}
\def\bn@font@montserrat{\RequirePackage[defaultfam,tabular,lining]{montserrat}}
\def\bn@font@lato      {\RequirePackage[default]{lato}}
\def\bn@font@cmbright  {\RequirePackage{cmbright}}
\def\bn@font@lm        {\RequirePackage{lmodern}}
\@nameuse{bn@font@\bnFont}
\makeatother
\renewcommand{\familydefault}{\sfdefault}

% --- Palette ----------------------------------------------------------------
%
% Five role-named slots. \bnDarkMode=1 swaps to a dark variant.
% Override any individual color by \definecolor'ing it AFTER % Shared definitions for the badness brand artifacts.
%
% \input by the per-artifact .tex files (logo.tex, wordmark.tex, og.tex).
% Anything \providecommand'd here is overridable by the caller — define
% your override BEFORE % Shared definitions for the badness brand artifacts.
%
% \input by the per-artifact .tex files (logo.tex, wordmark.tex, og.tex).
% Anything \providecommand'd here is overridable by the caller — define
% your override BEFORE \input{_common.tex}.

\usepackage{tikz}
\usepackage{xcolor}

% --- Font dispatch ----------------------------------------------------------
%
% \bnFont selects the sans-serif family. Set it in your artifact file
% before \input{_common.tex}. Available values:
%
%   plex       IBM Plex Sans         — modern technical (default)
%   tgheros    TeX Gyre Heros        — Helvetica-clone, neutral
%   source     Source Sans Pro       — humanist, friendly
%   roboto     Roboto                — Google flat-ui geometric
%   fira       Fira Sans             — Mozilla, slightly quirky
%   inter      Inter                 — UI workhorse, very tight
%   montserrat Montserrat            — wide geometric, posterish
%   lato       Lato                  — warm humanist
%   cmbright   Computer Modern Bright — light, scholarly
%   lm         Latin Modern Sans     — LaTeX default, neutral
%
\providecommand{\bnFont}{lm}

\makeatletter
\def\bn@font@plex      {\RequirePackage{plex-sans}}
\def\bn@font@tgheros   {\RequirePackage{tgheros}}
\def\bn@font@source    {\RequirePackage[default]{sourcesanspro}}
\def\bn@font@roboto    {\RequirePackage[sfdefault]{roboto}}
\def\bn@font@fira      {\RequirePackage[sfdefault]{FiraSans}}
\def\bn@font@inter     {\RequirePackage{inter}}
\def\bn@font@montserrat{\RequirePackage[defaultfam,tabular,lining]{montserrat}}
\def\bn@font@lato      {\RequirePackage[default]{lato}}
\def\bn@font@cmbright  {\RequirePackage{cmbright}}
\def\bn@font@lm        {\RequirePackage{lmodern}}
\@nameuse{bn@font@\bnFont}
\makeatother
\renewcommand{\familydefault}{\sfdefault}

% --- Palette ----------------------------------------------------------------
%
% Five role-named slots. \bnDarkMode=1 swaps to a dark variant.
% Override any individual color by \definecolor'ing it AFTER \input{_common.tex}.

\providecommand{\bnDarkMode}{0}

\ifnum\bnDarkMode=1
  % Dark variant
  \definecolor{bnInk}    {HTML}{F2EFE9} % "ink" reads light on dark
  \definecolor{bnBad}    {HTML}{F08A57} % brighter orange so it pops on dark
  \definecolor{bnPaper}  {HTML}{1B1B1B} % near-black canvas
  \definecolor{bnMargin} {HTML}{3D3833} % subtle margin guide
  \definecolor{bnMuted}  {HTML}{A39A8C} % secondary text
\else
  % Light variant (default)
  \definecolor{bnInk}    {HTML}{1B1B1B}
  \definecolor{bnBad}    {HTML}{D9633F}
  \definecolor{bnPaper}  {HTML}{FAFAF7}
  \definecolor{bnMargin} {HTML}{B5B0A5}
  \definecolor{bnMuted}  {HTML}{6F665A}
\fi

% --- Logo geometry (all in cm; tweak freely) --------------------------------

\providecommand{\bnBarCount}{4}     % number of horizontal bars
\providecommand{\bnBarWidth}{4.0}   % length of an in-margin bar
\providecommand{\bnBarHeight}{0.55} % bar thickness
\providecommand{\bnBarGap}{0.32}    % vertical gap between bars
\providecommand{\bnJutWhich}{3}     % 1..bnBarCount — which bar overflows
\providecommand{\bnJutExtra}{1.15}  % how far it sticks out
\providecommand{\bnRound}{0.07}     % corner rounding
\providecommand{\bnShowMargin}{0}   % 1 = draw the dashed margin guide, 0 = hide

% --- The mark itself --------------------------------------------------------
%
% Draws into the current tikzpicture, anchored so the bottom-left of the
% non-jutting bars is at (0,0). The jutting bar extends to the right.

\newcommand{\bnDrawMark}{%
  \pgfmathtruncatemacro{\bnLast}{\bnBarCount-1}%
  \pgfmathsetmacro{\bnTotalH}{\bnBarCount*\bnBarHeight + \bnLast*\bnBarGap}%
  \ifnum\bnShowMargin=1
    % The right-margin guide: a dashed vertical line at the in-margin width,
    % running the full height of the bar stack with a touch of overshoot.
    \draw[bnMargin, dashed, line width=0.4pt, dash pattern=on 2pt off 2pt]
    (\bnBarWidth, -0.12) -- (\bnBarWidth, \bnTotalH+0.12);
  \fi
  \foreach \i in {0,...,\bnLast}{%
      \pgfmathsetmacro{\bny}{(\bnLast-\i)*(\bnBarHeight+\bnBarGap)}%
      \pgfmathtruncatemacro{\bnNum}{\i+1}%
      \ifnum\bnNum=\bnJutWhich
        \fill[bnBad, rounded corners=\bnRound cm]
        (0,\bny) rectangle (\bnBarWidth+\bnJutExtra, \bny+\bnBarHeight);
      \else
        \fill[bnInk, rounded corners=\bnRound cm]
        (0,\bny) rectangle (\bnBarWidth, \bny+\bnBarHeight);
      \fi
    }%
}
.

\usepackage{tikz}
\usepackage{xcolor}

% --- Font dispatch ----------------------------------------------------------
%
% \bnFont selects the sans-serif family. Set it in your artifact file
% before % Shared definitions for the badness brand artifacts.
%
% \input by the per-artifact .tex files (logo.tex, wordmark.tex, og.tex).
% Anything \providecommand'd here is overridable by the caller — define
% your override BEFORE \input{_common.tex}.

\usepackage{tikz}
\usepackage{xcolor}

% --- Font dispatch ----------------------------------------------------------
%
% \bnFont selects the sans-serif family. Set it in your artifact file
% before \input{_common.tex}. Available values:
%
%   plex       IBM Plex Sans         — modern technical (default)
%   tgheros    TeX Gyre Heros        — Helvetica-clone, neutral
%   source     Source Sans Pro       — humanist, friendly
%   roboto     Roboto                — Google flat-ui geometric
%   fira       Fira Sans             — Mozilla, slightly quirky
%   inter      Inter                 — UI workhorse, very tight
%   montserrat Montserrat            — wide geometric, posterish
%   lato       Lato                  — warm humanist
%   cmbright   Computer Modern Bright — light, scholarly
%   lm         Latin Modern Sans     — LaTeX default, neutral
%
\providecommand{\bnFont}{lm}

\makeatletter
\def\bn@font@plex      {\RequirePackage{plex-sans}}
\def\bn@font@tgheros   {\RequirePackage{tgheros}}
\def\bn@font@source    {\RequirePackage[default]{sourcesanspro}}
\def\bn@font@roboto    {\RequirePackage[sfdefault]{roboto}}
\def\bn@font@fira      {\RequirePackage[sfdefault]{FiraSans}}
\def\bn@font@inter     {\RequirePackage{inter}}
\def\bn@font@montserrat{\RequirePackage[defaultfam,tabular,lining]{montserrat}}
\def\bn@font@lato      {\RequirePackage[default]{lato}}
\def\bn@font@cmbright  {\RequirePackage{cmbright}}
\def\bn@font@lm        {\RequirePackage{lmodern}}
\@nameuse{bn@font@\bnFont}
\makeatother
\renewcommand{\familydefault}{\sfdefault}

% --- Palette ----------------------------------------------------------------
%
% Five role-named slots. \bnDarkMode=1 swaps to a dark variant.
% Override any individual color by \definecolor'ing it AFTER \input{_common.tex}.

\providecommand{\bnDarkMode}{0}

\ifnum\bnDarkMode=1
  % Dark variant
  \definecolor{bnInk}    {HTML}{F2EFE9} % "ink" reads light on dark
  \definecolor{bnBad}    {HTML}{F08A57} % brighter orange so it pops on dark
  \definecolor{bnPaper}  {HTML}{1B1B1B} % near-black canvas
  \definecolor{bnMargin} {HTML}{3D3833} % subtle margin guide
  \definecolor{bnMuted}  {HTML}{A39A8C} % secondary text
\else
  % Light variant (default)
  \definecolor{bnInk}    {HTML}{1B1B1B}
  \definecolor{bnBad}    {HTML}{D9633F}
  \definecolor{bnPaper}  {HTML}{FAFAF7}
  \definecolor{bnMargin} {HTML}{B5B0A5}
  \definecolor{bnMuted}  {HTML}{6F665A}
\fi

% --- Logo geometry (all in cm; tweak freely) --------------------------------

\providecommand{\bnBarCount}{4}     % number of horizontal bars
\providecommand{\bnBarWidth}{4.0}   % length of an in-margin bar
\providecommand{\bnBarHeight}{0.55} % bar thickness
\providecommand{\bnBarGap}{0.32}    % vertical gap between bars
\providecommand{\bnJutWhich}{3}     % 1..bnBarCount — which bar overflows
\providecommand{\bnJutExtra}{1.15}  % how far it sticks out
\providecommand{\bnRound}{0.07}     % corner rounding
\providecommand{\bnShowMargin}{0}   % 1 = draw the dashed margin guide, 0 = hide

% --- The mark itself --------------------------------------------------------
%
% Draws into the current tikzpicture, anchored so the bottom-left of the
% non-jutting bars is at (0,0). The jutting bar extends to the right.

\newcommand{\bnDrawMark}{%
  \pgfmathtruncatemacro{\bnLast}{\bnBarCount-1}%
  \pgfmathsetmacro{\bnTotalH}{\bnBarCount*\bnBarHeight + \bnLast*\bnBarGap}%
  \ifnum\bnShowMargin=1
    % The right-margin guide: a dashed vertical line at the in-margin width,
    % running the full height of the bar stack with a touch of overshoot.
    \draw[bnMargin, dashed, line width=0.4pt, dash pattern=on 2pt off 2pt]
    (\bnBarWidth, -0.12) -- (\bnBarWidth, \bnTotalH+0.12);
  \fi
  \foreach \i in {0,...,\bnLast}{%
      \pgfmathsetmacro{\bny}{(\bnLast-\i)*(\bnBarHeight+\bnBarGap)}%
      \pgfmathtruncatemacro{\bnNum}{\i+1}%
      \ifnum\bnNum=\bnJutWhich
        \fill[bnBad, rounded corners=\bnRound cm]
        (0,\bny) rectangle (\bnBarWidth+\bnJutExtra, \bny+\bnBarHeight);
      \else
        \fill[bnInk, rounded corners=\bnRound cm]
        (0,\bny) rectangle (\bnBarWidth, \bny+\bnBarHeight);
      \fi
    }%
}
. Available values:
%
%   plex       IBM Plex Sans         — modern technical (default)
%   tgheros    TeX Gyre Heros        — Helvetica-clone, neutral
%   source     Source Sans Pro       — humanist, friendly
%   roboto     Roboto                — Google flat-ui geometric
%   fira       Fira Sans             — Mozilla, slightly quirky
%   inter      Inter                 — UI workhorse, very tight
%   montserrat Montserrat            — wide geometric, posterish
%   lato       Lato                  — warm humanist
%   cmbright   Computer Modern Bright — light, scholarly
%   lm         Latin Modern Sans     — LaTeX default, neutral
%
\providecommand{\bnFont}{lm}

\makeatletter
\def\bn@font@plex      {\RequirePackage{plex-sans}}
\def\bn@font@tgheros   {\RequirePackage{tgheros}}
\def\bn@font@source    {\RequirePackage[default]{sourcesanspro}}
\def\bn@font@roboto    {\RequirePackage[sfdefault]{roboto}}
\def\bn@font@fira      {\RequirePackage[sfdefault]{FiraSans}}
\def\bn@font@inter     {\RequirePackage{inter}}
\def\bn@font@montserrat{\RequirePackage[defaultfam,tabular,lining]{montserrat}}
\def\bn@font@lato      {\RequirePackage[default]{lato}}
\def\bn@font@cmbright  {\RequirePackage{cmbright}}
\def\bn@font@lm        {\RequirePackage{lmodern}}
\@nameuse{bn@font@\bnFont}
\makeatother
\renewcommand{\familydefault}{\sfdefault}

% --- Palette ----------------------------------------------------------------
%
% Five role-named slots. \bnDarkMode=1 swaps to a dark variant.
% Override any individual color by \definecolor'ing it AFTER % Shared definitions for the badness brand artifacts.
%
% \input by the per-artifact .tex files (logo.tex, wordmark.tex, og.tex).
% Anything \providecommand'd here is overridable by the caller — define
% your override BEFORE \input{_common.tex}.

\usepackage{tikz}
\usepackage{xcolor}

% --- Font dispatch ----------------------------------------------------------
%
% \bnFont selects the sans-serif family. Set it in your artifact file
% before \input{_common.tex}. Available values:
%
%   plex       IBM Plex Sans         — modern technical (default)
%   tgheros    TeX Gyre Heros        — Helvetica-clone, neutral
%   source     Source Sans Pro       — humanist, friendly
%   roboto     Roboto                — Google flat-ui geometric
%   fira       Fira Sans             — Mozilla, slightly quirky
%   inter      Inter                 — UI workhorse, very tight
%   montserrat Montserrat            — wide geometric, posterish
%   lato       Lato                  — warm humanist
%   cmbright   Computer Modern Bright — light, scholarly
%   lm         Latin Modern Sans     — LaTeX default, neutral
%
\providecommand{\bnFont}{lm}

\makeatletter
\def\bn@font@plex      {\RequirePackage{plex-sans}}
\def\bn@font@tgheros   {\RequirePackage{tgheros}}
\def\bn@font@source    {\RequirePackage[default]{sourcesanspro}}
\def\bn@font@roboto    {\RequirePackage[sfdefault]{roboto}}
\def\bn@font@fira      {\RequirePackage[sfdefault]{FiraSans}}
\def\bn@font@inter     {\RequirePackage{inter}}
\def\bn@font@montserrat{\RequirePackage[defaultfam,tabular,lining]{montserrat}}
\def\bn@font@lato      {\RequirePackage[default]{lato}}
\def\bn@font@cmbright  {\RequirePackage{cmbright}}
\def\bn@font@lm        {\RequirePackage{lmodern}}
\@nameuse{bn@font@\bnFont}
\makeatother
\renewcommand{\familydefault}{\sfdefault}

% --- Palette ----------------------------------------------------------------
%
% Five role-named slots. \bnDarkMode=1 swaps to a dark variant.
% Override any individual color by \definecolor'ing it AFTER \input{_common.tex}.

\providecommand{\bnDarkMode}{0}

\ifnum\bnDarkMode=1
  % Dark variant
  \definecolor{bnInk}    {HTML}{F2EFE9} % "ink" reads light on dark
  \definecolor{bnBad}    {HTML}{F08A57} % brighter orange so it pops on dark
  \definecolor{bnPaper}  {HTML}{1B1B1B} % near-black canvas
  \definecolor{bnMargin} {HTML}{3D3833} % subtle margin guide
  \definecolor{bnMuted}  {HTML}{A39A8C} % secondary text
\else
  % Light variant (default)
  \definecolor{bnInk}    {HTML}{1B1B1B}
  \definecolor{bnBad}    {HTML}{D9633F}
  \definecolor{bnPaper}  {HTML}{FAFAF7}
  \definecolor{bnMargin} {HTML}{B5B0A5}
  \definecolor{bnMuted}  {HTML}{6F665A}
\fi

% --- Logo geometry (all in cm; tweak freely) --------------------------------

\providecommand{\bnBarCount}{4}     % number of horizontal bars
\providecommand{\bnBarWidth}{4.0}   % length of an in-margin bar
\providecommand{\bnBarHeight}{0.55} % bar thickness
\providecommand{\bnBarGap}{0.32}    % vertical gap between bars
\providecommand{\bnJutWhich}{3}     % 1..bnBarCount — which bar overflows
\providecommand{\bnJutExtra}{1.15}  % how far it sticks out
\providecommand{\bnRound}{0.07}     % corner rounding
\providecommand{\bnShowMargin}{0}   % 1 = draw the dashed margin guide, 0 = hide

% --- The mark itself --------------------------------------------------------
%
% Draws into the current tikzpicture, anchored so the bottom-left of the
% non-jutting bars is at (0,0). The jutting bar extends to the right.

\newcommand{\bnDrawMark}{%
  \pgfmathtruncatemacro{\bnLast}{\bnBarCount-1}%
  \pgfmathsetmacro{\bnTotalH}{\bnBarCount*\bnBarHeight + \bnLast*\bnBarGap}%
  \ifnum\bnShowMargin=1
    % The right-margin guide: a dashed vertical line at the in-margin width,
    % running the full height of the bar stack with a touch of overshoot.
    \draw[bnMargin, dashed, line width=0.4pt, dash pattern=on 2pt off 2pt]
    (\bnBarWidth, -0.12) -- (\bnBarWidth, \bnTotalH+0.12);
  \fi
  \foreach \i in {0,...,\bnLast}{%
      \pgfmathsetmacro{\bny}{(\bnLast-\i)*(\bnBarHeight+\bnBarGap)}%
      \pgfmathtruncatemacro{\bnNum}{\i+1}%
      \ifnum\bnNum=\bnJutWhich
        \fill[bnBad, rounded corners=\bnRound cm]
        (0,\bny) rectangle (\bnBarWidth+\bnJutExtra, \bny+\bnBarHeight);
      \else
        \fill[bnInk, rounded corners=\bnRound cm]
        (0,\bny) rectangle (\bnBarWidth, \bny+\bnBarHeight);
      \fi
    }%
}
.

\providecommand{\bnDarkMode}{0}

\ifnum\bnDarkMode=1
  % Dark variant
  \definecolor{bnInk}    {HTML}{F2EFE9} % "ink" reads light on dark
  \definecolor{bnBad}    {HTML}{F08A57} % brighter orange so it pops on dark
  \definecolor{bnPaper}  {HTML}{1B1B1B} % near-black canvas
  \definecolor{bnMargin} {HTML}{3D3833} % subtle margin guide
  \definecolor{bnMuted}  {HTML}{A39A8C} % secondary text
\else
  % Light variant (default)
  \definecolor{bnInk}    {HTML}{1B1B1B}
  \definecolor{bnBad}    {HTML}{D9633F}
  \definecolor{bnPaper}  {HTML}{FAFAF7}
  \definecolor{bnMargin} {HTML}{B5B0A5}
  \definecolor{bnMuted}  {HTML}{6F665A}
\fi

% --- Logo geometry (all in cm; tweak freely) --------------------------------

\providecommand{\bnBarCount}{4}     % number of horizontal bars
\providecommand{\bnBarWidth}{4.0}   % length of an in-margin bar
\providecommand{\bnBarHeight}{0.55} % bar thickness
\providecommand{\bnBarGap}{0.32}    % vertical gap between bars
\providecommand{\bnJutWhich}{3}     % 1..bnBarCount — which bar overflows
\providecommand{\bnJutExtra}{1.15}  % how far it sticks out
\providecommand{\bnRound}{0.07}     % corner rounding
\providecommand{\bnShowMargin}{0}   % 1 = draw the dashed margin guide, 0 = hide

% --- The mark itself --------------------------------------------------------
%
% Draws into the current tikzpicture, anchored so the bottom-left of the
% non-jutting bars is at (0,0). The jutting bar extends to the right.

\newcommand{\bnDrawMark}{%
  \pgfmathtruncatemacro{\bnLast}{\bnBarCount-1}%
  \pgfmathsetmacro{\bnTotalH}{\bnBarCount*\bnBarHeight + \bnLast*\bnBarGap}%
  \ifnum\bnShowMargin=1
    % The right-margin guide: a dashed vertical line at the in-margin width,
    % running the full height of the bar stack with a touch of overshoot.
    \draw[bnMargin, dashed, line width=0.4pt, dash pattern=on 2pt off 2pt]
    (\bnBarWidth, -0.12) -- (\bnBarWidth, \bnTotalH+0.12);
  \fi
  \foreach \i in {0,...,\bnLast}{%
      \pgfmathsetmacro{\bny}{(\bnLast-\i)*(\bnBarHeight+\bnBarGap)}%
      \pgfmathtruncatemacro{\bnNum}{\i+1}%
      \ifnum\bnNum=\bnJutWhich
        \fill[bnBad, rounded corners=\bnRound cm]
        (0,\bny) rectangle (\bnBarWidth+\bnJutExtra, \bny+\bnBarHeight);
      \else
        \fill[bnInk, rounded corners=\bnRound cm]
        (0,\bny) rectangle (\bnBarWidth, \bny+\bnBarHeight);
      \fi
    }%
}
.

\providecommand{\bnDarkMode}{0}

\ifnum\bnDarkMode=1
  % Dark variant
  \definecolor{bnInk}    {HTML}{F2EFE9} % "ink" reads light on dark
  \definecolor{bnBad}    {HTML}{F08A57} % brighter orange so it pops on dark
  \definecolor{bnPaper}  {HTML}{1B1B1B} % near-black canvas
  \definecolor{bnMargin} {HTML}{3D3833} % subtle margin guide
  \definecolor{bnMuted}  {HTML}{A39A8C} % secondary text
\else
  % Light variant (default)
  \definecolor{bnInk}    {HTML}{1B1B1B}
  \definecolor{bnBad}    {HTML}{D9633F}
  \definecolor{bnPaper}  {HTML}{FAFAF7}
  \definecolor{bnMargin} {HTML}{B5B0A5}
  \definecolor{bnMuted}  {HTML}{6F665A}
\fi

% --- Logo geometry (all in cm; tweak freely) --------------------------------

\providecommand{\bnBarCount}{4}     % number of horizontal bars
\providecommand{\bnBarWidth}{4.0}   % length of an in-margin bar
\providecommand{\bnBarHeight}{0.55} % bar thickness
\providecommand{\bnBarGap}{0.32}    % vertical gap between bars
\providecommand{\bnJutWhich}{3}     % 1..bnBarCount — which bar overflows
\providecommand{\bnJutExtra}{1.15}  % how far it sticks out
\providecommand{\bnRound}{0.07}     % corner rounding
\providecommand{\bnShowMargin}{0}   % 1 = draw the dashed margin guide, 0 = hide

% --- The mark itself --------------------------------------------------------
%
% Draws into the current tikzpicture, anchored so the bottom-left of the
% non-jutting bars is at (0,0). The jutting bar extends to the right.

\newcommand{\bnDrawMark}{%
  \pgfmathtruncatemacro{\bnLast}{\bnBarCount-1}%
  \pgfmathsetmacro{\bnTotalH}{\bnBarCount*\bnBarHeight + \bnLast*\bnBarGap}%
  \ifnum\bnShowMargin=1
    % The right-margin guide: a dashed vertical line at the in-margin width,
    % running the full height of the bar stack with a touch of overshoot.
    \draw[bnMargin, dashed, line width=0.4pt, dash pattern=on 2pt off 2pt]
    (\bnBarWidth, -0.12) -- (\bnBarWidth, \bnTotalH+0.12);
  \fi
  \foreach \i in {0,...,\bnLast}{%
      \pgfmathsetmacro{\bny}{(\bnLast-\i)*(\bnBarHeight+\bnBarGap)}%
      \pgfmathtruncatemacro{\bnNum}{\i+1}%
      \ifnum\bnNum=\bnJutWhich
        \fill[bnBad, rounded corners=\bnRound cm]
        (0,\bny) rectangle (\bnBarWidth+\bnJutExtra, \bny+\bnBarHeight);
      \else
        \fill[bnInk, rounded corners=\bnRound cm]
        (0,\bny) rectangle (\bnBarWidth, \bny+\bnBarHeight);
      \fi
    }%
}
.

\providecommand{\bnDarkMode}{0}

\ifnum\bnDarkMode=1
  % Dark variant
  \definecolor{bnInk}    {HTML}{F2EFE9} % "ink" reads light on dark
  \definecolor{bnBad}    {HTML}{F08A57} % brighter orange so it pops on dark
  \definecolor{bnPaper}  {HTML}{1B1B1B} % near-black canvas
  \definecolor{bnMargin} {HTML}{3D3833} % subtle margin guide
  \definecolor{bnMuted}  {HTML}{A39A8C} % secondary text
\else
  % Light variant (default)
  \definecolor{bnInk}    {HTML}{1B1B1B}
  \definecolor{bnBad}    {HTML}{D9633F}
  \definecolor{bnPaper}  {HTML}{FAFAF7}
  \definecolor{bnMargin} {HTML}{B5B0A5}
  \definecolor{bnMuted}  {HTML}{6F665A}
\fi

% --- Logo geometry (all in cm; tweak freely) --------------------------------

\providecommand{\bnBarCount}{4}     % number of horizontal bars
\providecommand{\bnBarWidth}{4.0}   % length of an in-margin bar
\providecommand{\bnBarHeight}{0.55} % bar thickness
\providecommand{\bnBarGap}{0.32}    % vertical gap between bars
\providecommand{\bnJutWhich}{3}     % 1..bnBarCount — which bar overflows
\providecommand{\bnJutExtra}{1.15}  % how far it sticks out
\providecommand{\bnRound}{0.07}     % corner rounding
\providecommand{\bnShowMargin}{0}   % 1 = draw the dashed margin guide, 0 = hide

% --- The mark itself --------------------------------------------------------
%
% Draws into the current tikzpicture, anchored so the bottom-left of the
% non-jutting bars is at (0,0). The jutting bar extends to the right.

\newcommand{\bnDrawMark}{%
  \pgfmathtruncatemacro{\bnLast}{\bnBarCount-1}%
  \pgfmathsetmacro{\bnTotalH}{\bnBarCount*\bnBarHeight + \bnLast*\bnBarGap}%
  \ifnum\bnShowMargin=1
    % The right-margin guide: a dashed vertical line at the in-margin width,
    % running the full height of the bar stack with a touch of overshoot.
    \draw[bnMargin, dashed, line width=0.4pt, dash pattern=on 2pt off 2pt]
    (\bnBarWidth, -0.12) -- (\bnBarWidth, \bnTotalH+0.12);
  \fi
  \foreach \i in {0,...,\bnLast}{%
      \pgfmathsetmacro{\bny}{(\bnLast-\i)*(\bnBarHeight+\bnBarGap)}%
      \pgfmathtruncatemacro{\bnNum}{\i+1}%
      \ifnum\bnNum=\bnJutWhich
        \fill[bnBad, rounded corners=\bnRound cm]
        (0,\bny) rectangle (\bnBarWidth+\bnJutExtra, \bny+\bnBarHeight);
      \else
        \fill[bnInk, rounded corners=\bnRound cm]
        (0,\bny) rectangle (\bnBarWidth, \bny+\bnBarHeight);
      \fi
    }%
}
